
\documentclass[11pt,oneside]{article}

\usepackage[T1]{fontenc}
\usepackage[utf8]{inputenc}
\usepackage{lmodern}
\usepackage{microtype}
\usepackage{geometry}
\geometry{margin=1in}
\usepackage{amsmath,amssymb,amsthm,mathtools}
\usepackage{array,booktabs,longtable,tabularx}
\usepackage{enumitem}
\usepackage{listings}
\usepackage{xcolor}
\usepackage{hyperref}
\usepackage{cleveref}
\usepackage{fancyhdr}
\usepackage{titlesec}

\hypersetup{
  colorlinks=true,
  linkcolor=blue!45!black,
  urlcolor=blue!45!black,
  citecolor=blue!45!black,
  pdftitle={Inquiry Calculus v2.0 --- Typed Relational Inquiry, Interrogative Succession, and Regenerative Control},
  pdfauthor={},
  pdfsubject={Typed relational programming language for inquiry}
}

\setlength{\headheight}{14pt}
\pagestyle{fancy}
\fancyhf{}
\fancyhead[L]{Inquiry Calculus}
\fancyhead[R]{Unified Canonical Specification v2.0}
\fancyfoot[C]{\thepage}

\setlist[itemize]{leftmargin=1.5em}
\setlist[enumerate]{leftmargin=1.7em}
\setlength{\parindent}{0pt}
\setlength{\parskip}{0.55em}
\setlength{\emergencystretch}{2em}
\setcounter{secnumdepth}{3}
\setcounter{tocdepth}{2}
\makeatletter
\renewcommand{\@pnumwidth}{3em}
\renewcommand{\@tocrmarg}{4em}
\renewcommand*\l@section{\@dottedtocline{1}{1.5em}{3em}}
\renewcommand*\l@subsection{\@dottedtocline{2}{4.5em}{3.8em}}
\makeatother

\newtheorem{definition}{Definition}[section]
\newtheorem{law}[definition]{Law}
\newtheorem{theorem}[definition]{Theorem}
\newtheorem{proposition}[definition]{Proposition}
\newtheorem{lemma}[definition]{Lemma}
\newtheorem{corollary}[definition]{Corollary}
\newtheorem{obligation}[definition]{Obligation}
\newtheorem{remark}[definition]{Remark}
\newtheorem{example}[definition]{Example}

\newcommand{\B}{\mathbb B}
\newcommand{\Rel}{\mathsf{Rel}}
\newcommand{\Prop}{\mathsf{Prop}}
\newcommand{\Prog}{\mathsf{Prog}}
\newcommand{\Raw}{\mathsf{Raw}}
\newcommand{\Result}{\mathsf{Result}}
\newcommand{\OpenIR}{\mathsf{OpenIR}}
\newcommand{\RelIR}{\mathsf{RelIR}}
\newcommand{\ProbeIR}{\mathsf{ProbeIR}}
\newcommand{\BoundaryIR}{\mathsf{BoundaryIR}}
\newcommand{\TraversalIR}{\mathsf{TraversalIR}}
\newcommand{\PromptIR}{\mathsf{PromptIR}}
\newcommand{\PatchIR}{\mathsf{PatchIR}}
\newcommand{\StateRef}{\mathsf{StateRef}}
\newcommand{\ReturnRef}{\mathsf{ReturnRef}}
\newcommand{\OpRef}{\mathsf{OpRef}}
\newcommand{\BoundaryRef}{\mathsf{BoundaryRef}}
\newcommand{\EventRef}{\mathsf{EventRef}}
\newcommand{\Fib}{\operatorname{Fib}}
\newcommand{\Comp}{\operatorname{Comp}}
\newcommand{\Dom}{\operatorname{Dom}}
\newcommand{\Img}{\operatorname{Im}}
\newcommand{\Ker}{\operatorname{Ker}}
\newcommand{\Obs}{\operatorname{Obs}}
\newcommand{\Con}{\operatorname{Con}}
\newcommand{\Sep}{\operatorname{Sep}}
\newcommand{\Prof}{\operatorname{Prof}}
\newcommand{\Depart}{\operatorname{Depart}}
\newcommand{\NegField}{\operatorname{NegField}}
\newcommand{\RetField}{\operatorname{RetField}}
\newcommand{\PosNeg}{\operatorname{PosNeg}}
\newcommand{\Sig}{\operatorname{Sig}}
\newcommand{\Occ}{\operatorname{Occ}}
\newcommand{\Trav}{\operatorname{Trav}}
\newcommand{\Stand}{\operatorname{Stand}}
\newcommand{\Supp}{\operatorname{Supp}}
\newcommand{\Closed}{\operatorname{Closed}}
\newcommand{\Unlock}{\operatorname{Unlock}}
\newcommand{\Applicable}{\operatorname{Applicable}}
\newcommand{\Preserve}{\operatorname{Preserve}}
\newcommand{\Gain}{\operatorname{Gain}}
\newcommand{\Recover}{\operatorname{Recover}}
\newcommand{\Reopen}{\operatorname{Reopen}}
\newcommand{\Compile}{\operatorname{Compile}}
\newcommand{\Render}{\operatorname{Render}}
\newcommand{\Elab}{\operatorname{Elab}}
\newcommand{\Normalize}{\operatorname{Normalize}}
\newcommand{\Decode}{\operatorname{Decode}}
\newcommand{\Interpret}{\operatorname{Interpret}}
\newcommand{\Chk}{\operatorname{Check}}
\newcommand{\Actualize}{\operatorname{Actualize}}
\newcommand{\Formable}{\operatorname{Formable}}
\newcommand{\Executable}{\operatorname{Executable}}
\newcommand{\Answerable}{\operatorname{Answerable}}
\newcommand{\Productive}{\operatorname{Productive}}
\newcommand{\ResolvedQ}{\operatorname{ResolvedQ}}
\newcommand{\Ready}{\operatorname{Ready}}
\newcommand{\RequiredDischarge}{\operatorname{RequiredDischarge}}
\newcommand{\QSucc}{\operatorname{QSucc}}
\newcommand{\QStep}{\operatorname{QStep}}
\newcommand{\HeadQ}{\operatorname{HeadQ}}
\newcommand{\RootFrontier}{\operatorname{RootFrontier}}
\newcommand{\IFP}{\operatorname{IFP}}
\newcommand{\CueRaw}{\operatorname{CueRaw}}
\newcommand{\CueSupp}{\operatorname{CueSupp}}
\newcommand{\AskRef}{\mathfrak a}
\newcommand{\SupportedRoute}{\operatorname{SupportedByDeclaredRoute}}
\newcommand{\ND}{\operatorname{ND}}
\newcommand{\ResPath}{\varrho}
\newcommand{\den}[1]{[\![#1]\!]}
\newcommand{\eqH}{\equiv_{\mathcal H}}
\newcommand{\bind}{\mathbin{\mathsf{bind}}}
\newcommand{\numeq}[1]{\tag{#1}}

\crefname{definition}{definition}{definitions}
\crefname{law}{law}{laws}
\crefname{theorem}{theorem}{theorems}
\crefname{proposition}{proposition}{propositions}
\crefname{corollary}{corollary}{corollaries}
\crefname{obligation}{obligation}{obligations}

\title{\textbf{Inquiry Calculus v2.0}\\
\large Typed Relational Inquiry, Interrogative Succession, and Regenerative Control\\
\normalsize A Unified, Well-Defined, Well-Typed Relational Programming Language}
\author{}
\date{Version 2.0 --- Unified Canonical Specification\\26 August 2026}

\begin{document}
\maketitle

\begin{abstract}
This document specifies Inquiry Calculus v2.0 as the single forward semantic,
programming-language, compiler, runtime, and regenerative-control authority for the project.
Version 2.0 conservatively embeds the standing v1.1 typed relational calculus and consolidates its
positive-negation successor, paired-actuality, and interrogative-succession results into one
coherent language architecture.  The major version marks that consolidation; it does not license
accidental semantic strengthening.  The primary represented unit is a typed \emph{form}: any
consequentially represented value, relation, distinction, question, operator, program, or reified
construction that may lawfully enter further typed relations.  Object-role and operator-role are
therefore relational roles under typed interfaces rather than primitive ontological species.

Every pair of represented forms shares at least the coarsest relation of joint admission to the current represented field.  Inquiry seeks consequentially finer relations.  A realized distinction is not exhausted by a binary label or by a four-cell display.  It admits a recursively compositional candidate-boundary form
\[
D=(X,Y,B_D,\pi_X,\pi_Y,\Gamma_D)
\]
whose projections supply candidate side incidence, but neither projection is thereby an exterior or a negation.  Reciprocal inquiry begins from an explicit supported determination presentation, positively witnesses departure through a typed oriented negation use, fills that negation, and returns through the reverse section of the same relation.  After an explicit seed/reorientation, the process repeats from the exterior.  The resulting dependent sixfold occurrence
\[
\Xi_D(\omega)=(S_X,O_X,R_X;S_Y,O_Y,R_Y)
\]
retains the two negation-use identities, their semantic and execution coverage, the use-specific return fibers, protected recovery, seed support, residuals, and provenance.  The six roles are not independent slots, and \(\Gamma_D\) checks joint compatibility only after they have been generated.  The oriented double cover \(\widetilde B_D=B_D\times\{+,-\}\) records orientation without synthesizing a reverse negation.  The familiar four-cell square remains only a projection of the dependent occurrence.

A question is a partially bound typed relation.  Partial binding is a regenerative abstraction:
bind the known ports, preserve the relation, and the remaining open section/fiber is the question.
Removing a filling while preserving all relations through which it mattered leaves a structured
hole.  The intersection of the solution fibers of a relational web is the field of lawful
replacements.  A representation understands or determines a form, relative to a protected
horizon, when this residual web regenerates a unique protected equivalence class and the
discriminators needed to reopen the abstraction remain available.

This yields a direct account of generative inquiry and scientific representation invention.  If a current representation identifies two cases whose protected consequences differ, the representation is insufficient.  The calculus then opens higher-order questions for a separating context, a finer representation, a more appropriate grain, a new probe or instrument, or a binding extension.  Tool construction is therefore probe-basis extension: implementing a discriminator that makes a consequentially hidden difference observable.  The dual operation is abstraction: when an active difference never changes protected consequence, it becomes a candidate for folding.

The source inquiry language is an answer-dependent question-program calculus
\[
K::=\mathsf{Return}_I(a)\mid\mathsf{Ask}(q,\kappa),
\]
which may execute directly in an ordinary LLM chat when its open ports are generative, and lowers through typed discharge authority to the runtime effect language
\[
P::=\mathsf{Return}\mid\mathsf{Branch}\mid\mathsf{Probe}.
\]
A semantic question compiles to deterministic computation, LLM generation, an actual probe, an independent checker, a warrant operation, or a mixed program according to the modes of its open ports.  Raw return, decoded completion, interpretation, checking, warrant, and standing remain distinct.  Realized history remains authoritative evidence; learned distinctions, representations, question paths, and operator folds are compact regenerative structures over that history rather than replacements for it.

The derived interrogative algebra makes answer-conditioned succession explicit without adding a second semantics.  A dynamic successor is indexed by the checked \(\mathsf{Ask}\) occurrence and its first-order continuation, because the same semantic question and supported answer under different continuations may lead to different successor questions.  A live frontier retains both discretionary productive questions and questions required to discharge an explicit \(\mathsf{Probe}\), \(\mathsf{Check}\), \(\mathsf{Warrant}\), support, or reconstruction obligation.  Surface interrogatives and route labels remain erasable renderings of ordinary relations, questions, and \(\mathsf{IProg}\).

Question-side and return-side actuality are likewise derived projections of one ordinary authoritative event spine.  The canonical reconstruction chain is a checked \(\mathsf{Ask}\) occurrence, its ordinary event and raw return, the exact resolution path, the whole supported answer set, the occurrence-specific continuation, and the next \(\mathsf{Ask}\) occurrence.  Resume is not replay, endpoint equality is not path identity, and no paired or question-route projection creates a second history.

The result is a recursively closed, well-typed calculus for forming distinctions, positively determining typed exteriority, returning through use-specific fibers, opening holes, generating and executing questions, inventing representations and probes, traversing reciprocal boundaries, preserving actuality, learning regenerative forms and operator paths, folding and reopening structure, and applying the same operations to its own presentation without self-warrant.
\end{abstract}

\tableofcontents
\newpage

\part{Standing, Authority, and Constitutional Laws}

\section{Standing of this document}

\subsection{Forward authority}

This specification is the one cohesive document from which the active Inquiry Calculus is
reconstructed.  Earlier semantic documents remain regression ancestry, derivation evidence, and
sources of counterexamples.  They do not silently override this document merely because they
contain older notation or a larger architecture.

The complete immediately preceding v1.1 authority surface is retained at the pre-consolidation Git
coordinate \texttt{a936ebf}; the decision ledger records its full identity and the independently
received interrogative-source digest.  It remains ancestry and regression evidence rather than a
coequal forward authority.

A later successor may replace this document only by a predecessor-judged revision that:
\begin{enumerate}
\item preserves every still-valid protected capability;
\item gives an explicit disposition for any predecessor obligation it removes;
\item establishes at least one typed strict gain or a necessary correction;
\item preserves authoritative ancestry and reopening routes;
\item receives standing independently of the candidate's own assertion that it is better.
\end{enumerate}

\subsection{Conservative embedding of v1.1}

\begin{theorem}[v1.1 conservative embedding]\label{thm:v11-embedding}
There is a typed injective translation
\[
\boxed{
\operatorname{Embed}_{1.1\to2.0}:
\mathcal L_{1.1}\hookrightarrow\mathcal L_{2.0}
}
\]
that maps every admitted v1.1 typed form, relation, partial-binding question, supported answer,
\(\mathsf{Return}_I/\mathsf{Ask}\) source program, runtime program, discharge mode, event/raw
return, support/standing construction, positive-negation reciprocal occurrence, fold, bridge, and
predecessor-judged revision to the same typed semantic construction in v2.0.  For every v1.1 term
\(t\) and binding carried by the translation,
\[
\boxed{
\den{\operatorname{Embed}_{1.1\to2.0}(t)}_{2.0}
=
\den{t}_{1.1}.
}
\]
For a predecessor supported answer with semantic member set \(S\), the proof-carrying v2.0
translation satisfies
\[
\boxed{
\left|\operatorname{EmbedSupp}_{1.1\to2.0}(S,\vartheta,E,w)\right|=S,
}
\]
retains the predecessor route/evidence witness, and selects the same source continuation branch.
The predecessor \(\mathsf{Code}(A)\) carrier maps to the explicitly typed reified-code family of
\cref{def:program-code-families}; it is neither dropped nor identified with execution.
No v1.1 discharge port is weakened, no generated form becomes actual or standing, and no existing
event, use identity, return fiber, support route, or reopening condition is erased by embedding.
\end{theorem}

\begin{proof}
The v2.0 substrate retains the v1.1 type/form universe, typed relation denotation, named-port
question constructor, supported-answer semantics, source and runtime grammars, five discharge
modes, ordinary event spine, least-fixed-point standing, positive-negation contracts, and
fold/bridge/revision laws unchanged.  The new interrogative and paired-actuality structures are
typed derived annotations or projections over those constructions and erase by
\cref{thm:interrogative-lowering,law:one-occurrence-spine}.  Structural inclusion therefore
preserves typing and denotation, and no new equality is imposed on old forms.
\end{proof}

\begin{corollary}[Evidence transport without conformance inflation]\label{cor:v11-evidence-transport}
An executable v1.1 fixture remains evidence for the embedded v2.0 behavior at exactly its original
scope, applicability, grain, horizon, versions, and coverage.  The embedding alone establishes no
new interrogative, controller, or consolidation fixture as passing.
\end{corollary}

\begin{law}[Major version non-strengthening]\label{law:v2-nonstrengthening}
The v2.0 name records a coherent language and control consolidation.  It does not authorize
strengthening converse into inverse, possibility into actuality, parsing into support, checking
into warrant, ledger order into causality, or derived projections into new authority.
\end{law}

\subsection{Five statuses}

Every construct in this document is explicitly one of:
\begin{enumerate}
\item \textbf{constitutional}: a law of the calculus;
\item \textbf{canonical-restated}: a constitutional or previously admitted construct repeated locally so that a derivation is readable, without becoming a second definition;
\item \textbf{derived}: regenerable from constitutional and canonical-restated structure;
\item \textbf{binding-supplied}: additional structure supplied by a domain or runtime binding;
\item \textbf{implementation-only}: an encoding needed by a compiler/runtime but not an ontological primitive.
\end{enumerate}

A convenient implementation sort never becomes constitutional merely because code requires a record for it.

Unless a local status line says otherwise, definitions, typing judgments, laws, and theorems in the
typed semantic-language parts are \textbf{constitutional}; a repeated presentation explicitly
identified as such is \textbf{canonical-restated}; constructions in the derived-structures and
interrogative-algebra sections are \textbf{derived}; domain and mathematical-home choices are
\textbf{binding-supplied}; and concrete runtime record layouts, indexes, storage, and backend
encodings are \textbf{implementation-only}.  A cross-reference to a constitutional contract from an
implementation-only record does not transfer the record's layout into semantics.  These defaults
give every named construction exactly one status while permitting a local explicit override.

Material consolidated from the paired-actuality and interrogative-succession analyses below is
explicitly \textbf{derived}.  Its inclusion in the single canonical source removes competing
authority paths; it does not promote a trace projection, root annotation, route classification,
or controller view into a constitutional primitive.

\section{Typed language architecture}\label{sec:language-architecture}

The organizing spine of v2.0 is the typed compilation and evidence stack
\[
\boxed{
\begin{array}{c}
\text{surface interrogative / binding-native syntax}\\
\downarrow\ \Elab\\
\text{derived interrogative roots, macros, and route annotations}\\
\downarrow\ \operatorname{Lower}_Q\\
\text{typed open relation }(?_I R[\beta])\\
\downarrow\\
\text{source }\mathsf{IProg}:\mathsf{Return}_I\mid\mathsf{Ask}\\
\downarrow\ \Compile\\
\text{runtime/effect IR}:\mathsf{Return}\mid\mathsf{Branch}\mid\mathsf{Probe}\\
\downarrow\\
\text{actual execution / immutable }\Raw\text{ return / ordinary event}\\
\downarrow\\
\text{resolution / support / check / warrant / standing}\\
\downarrow\\
\text{residual and occurrence-indexed successor-question frontier}.
\end{array}
}
\]

Each layer has one responsibility.  Surface syntax renders; derived roots construct ordinary
questions; \(\operatorname{Lower}_Q\) erases annotations; the source program carries typed
answer-dependent control; the runtime performs only its admitted effects; ordinary history records
actuality; resolution and authority relations determine what the return supports; the residual
constructs the next live occurrence.  No downstream convenience may redefine an upstream semantic
relation.

\begin{law}[Typed non-collapse across the language stack]\label{law:language-stack-noncollapse}
Every transition between layers preserves the represented type, binding, occurrence identity,
whole supported answer, discharge authority, provenance, and protected continuation required by
its contract.  In particular, raw provider bytes do not directly select a semantic continuation,
route annotations do not execute, and a controller policy does not create standing.
\end{law}

\section{Constitutional laws}

\begin{law}[Relational substrate]
\[
\boxed{\textbf{RELATIONS ARE THE SEMANTIC SUBSTRATE.}}
\]
``Object'', ``state'', ``question'', ``answer'', ``boundary'', ``memory'', ``operator'', ``method'', ``perception'', ``environment'', ``cause'', ``problem'', ``program'', and ``successor'' name roles occupied by typed relational arrangements.  They do not, merely by being nouns, introduce independent semantic species.
\end{law}


\begin{law}[Recursive form closure]
Let
\[
\mathsf{Form}_{\B}
=
\sum_{A:\mathsf{Ty}_{\B}}\den{A}_{\B}
\]
be the binding-local universe of typed represented forms.  Any consequentially represented relation, distinction, question, operator, program, method, type witness, or binding witness may be reified as a form and may enter further well-typed relations.  A form may play an operator role only through an admitted typed operational interpretation.  Hence:
\[
\boxed{
\begin{gathered}
\textbf{OBJECT-ROLE AND OPERATOR-ROLE ARE RELATIONAL ROLES,}\\
\textbf{NOT DISJOINT PRIMITIVE SPECIES.}
\end{gathered}}
\]
Typing is never erased by reification.
\end{law}

\begin{law}[Coarsest common relation]
For every represented pair \(x,y\in\mathsf{Form}_{\B}\),
\[
\boxed{\top_{\B}(x,y)}
\qquad\text{where}\qquad
\top_{\B}=\mathsf{Form}_{\B}\times\mathsf{Form}_{\B}.
\]
Thus ``unrelated'' never means that no relation exists.  It means that no consequentially useful proper refinement of the coarsest admitted relation has yet been established for the inquiry at hand.
\end{law}

\begin{law}[Refine where consequences split; fold where they do not]
If a current representation \(\eta\) identifies two forms whose protected consequences differ,
\[
\eta(x)=\eta(y)
\quad\land\quad
x\not\equiv_{\mathcal H}y,
\]
then the representation is insufficient and inquiry must preserve the obligation to seek a separating context, representation, grain, probe, or binding.  Conversely, an active distinction that no protected continuation can inspect is a candidate for licensed folding.  Hence:
\[
\boxed{\textbf{REFINE WHERE CONSEQUENCES SPLIT; FOLD WHERE THEY DO NOT.}}
\]
\end{law}

\begin{law}[Regenerative distinction]
A realized distinction occurrence is not merely a label placed between two pre-existing objects.  It is a typed regenerative form from which its candidate sides, the supported determination presentation, the relation uses that positively witness departure, and every protected return/recovery obligation are recoverable at the admitted grain.  Distinctions may themselves be reified as sides of further distinctions, opened at finer grain, or interpreted as operators under typed interfaces.
\end{law}

\begin{law}[Relational exclusion]
An instruction or renderer may suppress a form without establishing any semantic incompatibility.  A question establishes relational exclusion only when its answer cells are disjoint under the same relation, binding, scope, applicability, and grain.  Therefore:
\[
\boxed{\textbf{SURFACE EXCLUSION IS NOT RELATIONAL EXCLUSION.}}
\]
An apparent violation of a relational exclusion generates inquiry into the side claims, decoder, scope, grain, question frame, or missing distinction rather than merely stronger suppression.
\end{law}

\begin{law}[Protected discrimination]
For an admitted protected future \(\mathcal H\),
\[
\boxed{
\begin{gathered}
\textbf{DO NOT IDENTIFY WHAT A PROTECTED}\\
\textbf{CONTINUATION CAN DISTINGUISH.}
\end{gathered}
}
\]
Conversely, a distinction may be hidden from active use when no protected continuation can inspect it, provided required ancestry, residual information, recovery, and reopening remain available.
\end{law}

\begin{law}[Arrangement and succession]
\[
\boxed{\textbf{ARRANGEMENT}\;\perp\;\textbf{SUCCESSION}.}
\]
Relations that coexist in a represented slice and relations that say how slices succeed are independent axes.  Storage order, realized/domain succession, and inquiry traversal order are also distinct unless a binding proves them coincident.
\end{law}

\begin{law}[Function is relation]
\[
\boxed{\textbf{FUNCTION}\subset\textbf{RELATION}.}
\]
Pure deterministic computation is executed before open inquiry whenever its result is exactly available.
\end{law}

\begin{law}[Question law]
A semantic question is an open typed relation.  Its inquiry role is:
\[
\boxed{
\textbf{QUESTION}
=
\textbf{DIFFERENTIATION}
+
\textbf{CONSEQUENTIAL OCCLUSION}.
}
\]
It directs resolution toward an unresolved distinction and permits certified consequence-null context to be omitted from the active view.
\end{law}

\begin{law}[Boundary law]
A consequential local distinction is represented by a factorization
\[
\boxed{
D^+ \xleftarrow{\pi_+} B_D \xrightarrow{\pi_-} D^- .
}
\]
The boundary carrier is neither side by role.  A boundary position holds both candidate projections.  Projection alone establishes neither exteriority nor a typed negation use.
\end{law}

\begin{law}[Positive-negation recurrence]\label{law:positive-negation-recurrence}
The constitutive reciprocal act is:
\[
\boxed{
\begin{gathered}
\textbf{POSITIVELY DETERMINE A LAWFUL TYPED NEGATION}\\
\textbf{OF THE LIVE DETERMINATION, THEN RETURN THROUGH}\\
\textbf{THE RELATION THAT LICENSED IT.}
\end{gathered}
}
\]
For source determination \(x\) and an admitted oriented relation use \(u\), the dependency is
\[
x\longrightarrow N_u[x]
\longrightarrow ?y[N_u(x,y)]
\longrightarrow y
\longrightarrow N_u^{-1}[y].
\]
The exterior filling may then be seeded as a new source and the recurrence applied in the reciprocal orientation.  Boolean complement, boundary projection, failed search, non-equivalence, and unknown status do not supply the departure witness required by this law.
\end{law}

\begin{law}[Return is not reconciliation]\label{law:return-not-reconciliation}
Pure reciprocal return computes the reverse section of the same typed relation use.  It does not mutate the standing source web.  A protected change of standing meaning requires separately actualized, checked, and warranted retraction, revision, changed applicability, changed grain, changed binding, or evidence that the predecessor presentation was underdetermined.
\end{law}

\begin{law}[Traversal law]
Inquiry primarily traverses the opened boundary:
\[
\boxed{T_D:B_D\rightsquigarrow B_D.}
\]
Crossing from one side to the other is an optional test, not the definition of traversal.
\end{law}

\begin{law}[Actuality separation]
\[
\boxed{
\begin{gathered}
\text{represented possibility}\neq\text{probe code}\neq\text{attempt}\neq\text{raw return}\\
{}\neq\text{decoded completion}\neq\text{interpretation}\neq\text{checked claim}\\
{}\neq\text{standing relation}.
\end{gathered}
}
\]
No transition may erase these distinctions.
\end{law}

\begin{law}[Unified probe]
There is one actualizable interaction form.  A probe is interpreted as
\[
o:Z\rightsquigarrow R_o\times Z .
\]
Observation-like and action-like behavior are derived projections of a probe occurrence, not disjoint primitives.
\end{law}

\begin{law}[Linked history without semantic fusion]
Historical co-occurrence and succession must be preserved at the chosen grain, but relations that occurred together remain available for independent future recomposition.  Historical linkage is not permanent semantic fusion.
\end{law}

\begin{law}[Operator-first learning]
Retained state need not be destructively compressed.  Recurrent typed probe/traversal paths may be compiled into smaller reusable operators when protected behavior is preserved.  Every promoted operator retains provenance and an unlock/reopening contract.
\end{law}

\begin{law}[No self-warrant]
Questioning, generation, recursion, and support may be cyclic.  Positive standing cannot arise from a rootless support cycle.  Standing is the least fixed point generated from grounded ingress and independently discharged support.
\end{law}

\begin{law}[Successor judged by predecessor]
A candidate successor may not revise the acceptance conditions by which its own promotion is judged.  Protection changes are themselves explicit, typed, independently judged revisions.
\end{law}

\part{Typed Relational Substrate}

\section{Bindings}

\begin{definition}[Binding]
A binding is represented by
\[
\boxed{
\B=
(\mathsf{Ty},\mathsf{Int},\mathsf G,\mathsf{Rel},\circ,
\mathsf{ProbeSem},\mathsf{Check},\mathsf{Warrant},
\mathcal H,\mathcal R).
}
\]
Its components are:
\begin{itemize}
\item \(\mathsf{Ty}\): admitted base and constructed types;
\item \(\mathsf{Int}\): compositional interfaces;
\item \(\mathsf G\): primitive typed operator generators;
\item \(\mathsf{Rel}\): admitted relation schemas;
\item \(\circ\): lawful partial serial composition;
\item \(\mathsf{ProbeSem}\): semantics of actualizable probe code;
\item \(\mathsf{Check}\): independently admitted checking relations;
\item \(\mathsf{Warrant}\): standing/support policy;
\item \(\mathcal H\): protected future continuations, observations, or consequences;
\item \(\mathcal R\): explicit resource/effectivity horizon.
\end{itemize}
\end{definition}

Order, metric, probability, topology, Boolean algebra, linear structure, causality, control, concurrency, game semantics, or a particular notion of cost are not constitutional.  A binding may supply them.

\begin{definition}[Binding role]
A represented arrangement \(B\) plays the binding role for domain fragment \(D\) and horizon \(\mathcal H\) when
\[
\operatorname{Binds}(B,D,\mathcal H)
\]
and the relations required to type, compose, actualize, decode, check, and compare the protected forms are available.
\end{definition}

\begin{definition}[Rebinding]
When a protected residual cannot be expressed in the incumbent binding,
\[
\operatorname{Expressible}(\B,d)=\bot,
\]
rebinding is the ordinary higher-order question
\[
?\B'[
\operatorname{ExtendsOrBridges}(\B,\B')
\land
\operatorname{Expressible}(\B',d)
].
\]
A rebind must state which old terms, consequences, histories, and operator contracts remain transportable.
\end{definition}

\section{Types and judgments}

\subsection{Reference type grammar}

The reference implementation admits:
\[
\begin{aligned}
A,B ::= {}&
\alpha
\mid \mathbf 1
\mid \mathbf{Bool}
\mid \mathbf{Nat}
\mid A\times B
\mid A+B\\
&\mid \Sigma x:A.B(x)
\mid \Pi x:A.B(x)
\mid \mathsf{Fin}(A)
\mid \mathsf{List}(A)\\
&\mid \Raw(A)
\mid \Result(A)
\mid \mathsf{IProg}(A)
\mid \Prog(A)
\mid \mathsf{Code}(A).
\end{aligned}
\]
This is a programming-language representation.  A native binding need not claim that the represented domain intrinsically carries every one of these structures.

\begin{definition}[Program and code families]\label{def:program-code-families}
\(\mathsf{IProg}(A)\) is the first-order source syntax returning \(A\),
\(\Prog(A)\) is the executable runtime syntax returning \(A\), and
\(\mathsf{Code}(A)\) is nonexecuting reified program syntax.  Their constructors and typing rules
are given below.  There are typed quotations
\[
\mathsf{quote}_I:\mathsf{IProg}(A)\to\mathsf{Code}(A),
\qquad
\mathsf{quote}_P:\Prog(A)\to\mathsf{Code}(A),
\]
and a binding- and compiler-version-indexed partial interpretation
\[
\mathsf{Act}^{code}_{\B,v,A}:
\mathsf{Code}(A)
\rightharpoonup
\mathsf{IProg}(A)+\Prog(A).
\]
Quotation alone neither executes nor warrants code.  The partial interpretation is available only
through an admitted binding \(\B\) and compiler version \(v\).  Whenever \(v\) admits the
corresponding quotation form, it satisfies the typed quotation laws
\[
\mathsf{Act}^{code}_{\B,v,A}(\mathsf{quote}_I(K))=\mathsf{inl}(K),
\qquad
\mathsf{Act}^{code}_{\B,v,A}(\mathsf{quote}_P(P))=\mathsf{inr}(P).
\]
Outside the admitted syntax/version domain the interpretation is undefined rather than guessed.
Thus the predecessor \(\mathsf{Code}(A)\)
constructor has a total place in the v2.0 type grammar without becoming a second program semantics.
\end{definition}

\subsection{Implementation code sorts}

The reference compiler uses, as implementation-only sorts,
\[
\begin{aligned}
&\RelIR[A,B],\quad
\OpenIR[\Gamma;\Delta],\quad
\ProbeIR[R],\\
&\BoundaryIR,\quad
\TraversalIR,\quad
\PromptIR,\quad
\PatchIR[k],\\
&\StateRef,\quad
\ReturnRef,\quad
\OpRef,\quad
\BoundaryRef,\quad
\EventRef.
\end{aligned}
\]

Patch roles include
\[
k\in
\{
\mathsf{Semantic},
\mathsf{Traversal},
\mathsf{Compiler},
\mathsf{Binding},
\mathsf{Protection},
\mathsf{Implementation}
\}.
\]


\section{Recursive represented forms}

\begin{definition}[Typed represented form]
The binding-local universe of represented forms is the dependent sum
\[
\boxed{
\mathsf{Form}_{\B}
=
\sum_{A:\mathsf{Ty}_{\B}}\den{A}_{\B}.
}
\]
An element is written \(\langle A,x\rangle\), or simply \(x:A\) when the type is inferable.  The dependent sum allows differently typed values to participate in a common higher relation without identifying their types.
\end{definition}

\begin{definition}[Reification]
Every admitted typed operator term may be reified:
\[
\boxed{
\operatorname{Reify}_{A,B}:
\mathsf{Tm}_{\B}(A,B)
\to
\mathsf{Form}_{\B}.
}
\]
Write \(\ulcorner f\urcorner\) for \(\operatorname{Reify}(f)\).  Relation schemas, questions, distinction points, programs, methods, and presentation fragments may likewise be reified through their admitted code or semantic carrier.
\end{definition}

\begin{definition}[Operational interpretation]
A represented form may play an operator role only where the binding supplies a partial typed interpretation
\[
\boxed{
\operatorname{Act}_{A,B}:
\mathsf{Form}_{\B}
\rightharpoonup
\mathsf{Tm}_{\B}(A,B).
}
\]
Reification and operational interpretation need not be mutual inverses.  A form may admit several operational interpretations under different interfaces, scopes, or bindings.
\end{definition}

\begin{law}[Role closure without type collapse]
If \(f:A\rightsquigarrow B\), then \(\ulcorner f\urcorner\) may be used as an operand of another relation.  If a form \(z\) admits \(\operatorname{Act}_{C,D}(z)\), then the same represented form may play an operator role from \(C\) to \(D\).  This closure does not imply that every form is executable or that all forms share one type.
\end{law}

\begin{law}[Construction ancestry does not restrict recomposition]
If \(z\) was constructed from \(x,y\), authoritative ancestry may preserve that fact, but it does not prevent \(x\), \(y\), or \(z\) from later participating independently in other lawful relations and distinctions.  Construction ancestry is history, not permanent semantic containment.
\end{law}

\begin{definition}[Grain-relative opening]
For grains \(g\preceq g'\), a form may be opaque at \(g\) and open into a finer relational presentation at \(g'\):
\[
\boxed{
\operatorname{Open}_{g\to g'}:
\mathsf{Form}_{\B}
\rightsquigarrow
\mathsf{Presentation}_{\B}.
}
\]
No absolute atomicity is assumed.  A point is simply a form whose internal distinctions are not active at the current grain.
\end{definition}

\begin{law}[Atomicity is grain-relative]
\[
\boxed{\textbf{ATOMICITY IS RELATIVE TO THE CURRENT GRAIN OF INQUIRY.}}
\]
A later discriminator may lawfully reopen a previously opaque form into consequential internal structure.
\end{law}

\subsection{Judgments}

The principal judgments are:
\[
\B;\Gamma\vdash t:A,
\]
\[
\B;\Gamma\vdash\phi:\Prop,
\]
\[
\B;\Gamma\vdash R:A\rightsquigarrow B,
\]
\[
\B;\Gamma\vdash ?\Delta[\phi]\;\mathsf{question},
\]
\[
\B;\Gamma\vdash o:\ProbeIR[R],
\]
\[
\B;\Gamma\vdash K:\mathsf{IProg}(A),
\]
\[
\B;\Gamma\vdash P:\Prog(A),
\]
\[
\B;\Gamma\vdash \ResPath:A\leadsto B,
\]
and
\[
X_t\vdash \Stand(\lambda).
\]

\section{Relations}

\begin{definition}[Typed relation]
A binary relation from \(A\) to \(B\) is
\[
R:A\rightsquigarrow B
\quad\text{with}\quad
\den{R}_{\B}\subseteq \den A_{\B}\times\den B_{\B}.
\]
An \(n\)-ary relation schema is
\[
R\hookrightarrow A_1\times\cdots\times A_n.
\]
\end{definition}

\begin{definition}[Function]
A function is a total single-valued relation:
\[
f:A\to B
\iff
R_f:A\rightsquigarrow B
\land
\forall a\,\exists!b\,R_f(a,b).
\]
\end{definition}

\begin{definition}[Identity]
\[
\mathsf{id}_A:A\rightsquigarrow A,
\qquad
\mathsf{id}_A(a,a').
\iff a=a'.
\]
\end{definition}

\begin{definition}[Serial composition]
For
\[
R:A\rightsquigarrow B,\qquad S:B\rightsquigarrow C,
\]
\[
\boxed{
(S\circ R)(a,c)
\iff
\exists b:B.\ R(a,b)\land S(b,c).
}
\]
Thus relational composition is join plus hiding of the mediator.
\end{definition}

\begin{definition}[Converse]
Where admitted by the representation,
\[
R^\dagger:B\rightsquigarrow A,
\qquad
R^\dagger(b,a)\iff R(a,b).
\]
Converse is not an inverse and makes no uniqueness claim.
\end{definition}

\begin{definition}[Image and inverse image]
For \(X\subseteq A\) and \(Y\subseteq B\),
\[
R[X]=\{b:\exists a\in X.\ R(a,b)\},
\]
\[
R^{-1}[Y]=\{a:\exists b\in Y.\ R(a,b)\}.
\]
\end{definition}

\begin{definition}[Parallel terms]
Two operators are parallel when they have the same input and output interfaces.  Direct protected contextual equivalence is defined only for parallel terms or for terms connected by an explicit comparison bridge.
\end{definition}

\subsection{Free typed operator language}

Given primitive generators \(\mathsf G_{\B}\), define typed terms:
\[
\mathsf{Tm}_{\B}(A,B)
\]
as the least family containing identities and primitive generators and closed under lawful composition.  Syntactically:
\[
p ::= \mathsf{id}_A \mid g \mid p_2\circ p_1,
\]
subject to interface matching.

Parenthesization is quotientable by associativity only after the binding establishes the corresponding compositional law.  Path provenance may remain protected even when denotational composites coincide.


\section{Coarsest relation and relational refinement}

\begin{definition}[Coarsest represented relation]
On \(\mathsf{Form}_{\B}\), define
\[
\boxed{
\top_{\B}
=
\mathsf{Form}_{\B}\times\mathsf{Form}_{\B}.
}
\]
For represented forms \(x,y\), \(\top_{\B}(x,y)\) says only that they are jointly admitted to the current represented field.  It does not assert physical coexistence, causality, similarity, or any stronger domain relation.
\end{definition}

\begin{definition}[Binding-level coexistence]
If the binding supplies an actuality/existence predicate \(Exists_{\B}\), one may define the stronger relation
\[
Coexists_{\B}(x,y)
\iff
Exists_{\B}(x)\land Exists_{\B}(y).
\]
This construction is binding-supplied rather than constitutional.
\end{definition}

\begin{definition}[Relation-refinement preorder]
For relations \(R,S\) on a common carrier, define
\[
\boxed{
R\preceq_{\mathrm{rel}}S
\iff
\den S\subseteq\den R.
}
\]
Thus \(S\) is at least as discriminating/informative as \(R\).  The coarsest relation is least informative:
\[
\top_{\B}\preceq_{\mathrm{rel}}R
\]
for every relation \(R\) on the same represented carrier.
\end{definition}

\begin{definition}[Nontrivial relation discovery]
Given \(R(x,y)\), the higher-order question for a proper refinement is
\[
\boxed{
?R'[
R'(x,y)
\land
R\preceq_{\mathrm{rel}}R'
\land
\den{R'}\subsetneq\den R
].
}
\]
A discovered refinement is useful only if it changes a protected consequence, supports a protected discriminator, or supplies a regenerative factor needed by later inquiry.
\end{definition}

\begin{law}[No vacuity from universal relatedness]
The fact that \(\top_{\B}(x,y)\) always holds contributes almost no discrimination.  Meaningful inquiry concerns the refinement structure of relations and which refinements survive variation, probing, and protected consequence tests.
\end{law}

\section{Logical and relational expression language}

\subsection{Formula grammar}

The reference relation language uses:
\[
\begin{aligned}
\phi,\psi ::= {}&
\top
\mid \bot
\mid R(t_1,\ldots,t_n)
\mid t=u\\
&\mid \phi\land\psi
\mid \phi\lor\psi
\mid \phi\Rightarrow\psi
\mid \neg\phi\\
&\mid \exists x:A.\phi
\mid \forall x:A.\phi.
\end{aligned}
\]

The metalanguage is classical set-theoretic unless a binding explicitly chooses another logic.  Expressibility and decidability of complement or quantification inside a binding remain separate questions.

Logical negation is not identified with contextual typed negation or positive reciprocal exterior:
\[
\boxed{
\neg\phi
\neq N_u
\neq N_u[x]
\neq y
\quad\text{in general}.
}
\]
Here \(N_u\) is an ordinary supported oriented relation playing a negation role, \(N_u[x]\) is its source-bound section, and \(y\) is one positively supported filling.  The logical connective remains part of the formula language; it neither creates a \(\mathsf{NegationUse}\) nor supplies its coverage or departure evidence.

\subsection{Minimal generating basis}

In the classical reference dialect one may take
\[
\top,\quad =,\quad \land,\quad \exists,\quad \neg
\]
as a sufficient logical generating basis and derive:
\[
\phi\lor\psi
:=
\neg(\neg\phi\land\neg\psi),
\]
\[
\forall x.\phi
:=
\neg\exists x.\neg\phi,
\]
\[
\phi\Rightarrow\psi
:=
\neg\phi\lor\psi.
\]

These derivations are logical conveniences, not claims that every binding natively represents Boolean complement.

\subsection{Data-only relational query IR}

A small relation-expression IR is:
\[
\boxed{
E ::=
R
\mid \mathsf{Bind}(E,\beta)
\mid \mathsf{Join}(E,E)
\mid \mathsf{Expose}_I(E)
\mid \mathsf{Hide}_I(E)
\mid \mathsf{Rename}_{\sigma}(E)
\mid \mathsf{Guard}_G(E).
}
\]

Its denotation is inherited from ordinary relational join, substitution, existential hiding, variable renaming, and guarded restriction.  A concrete semantic question is an expression with a nonempty exposed port set.

\part{Relation Schemas, Questions, and Refinement}

\section{Relation schemas and named ports}

\begin{definition}[Relation schema]
A typed relation schema is
\[
R\hookrightarrow X_1\times\cdots\times X_n
\]
with named signature
\[
\operatorname{sig}(R)
=
(x_1:X_1,\ldots,x_n:X_n).
\]
Port names carry semantic and typing roles; they are not presentation decoration.
\end{definition}

\begin{definition}[Partial binding]
Let \(I\subseteq\{1,\ldots,n\}\) be the nonempty set of open ports.  A partial binding
\[
\beta
\]
assigns a well-typed value to every port outside \(I\).
\end{definition}

\begin{definition}[Completion fiber]
\[
\boxed{
\Fib_I(R\mid\beta)
=
\left\{
a_I\in\prod_{i\in I}X_i:
R[\beta\oplus a_I]
\right\}.
}
\]
\end{definition}

\begin{definition}[Question]
The canonical question generated by opening \(I\) is
\[
\boxed{
q=?_I R[\beta].
}
\]
Its complete semantic answer carrier is
\[
A(q)=\prod_{i\in I}X_i.
\]
Its valid completion relation is
\[
\Comp_q(a_I)\iff a_I\in\Fib_I(R\mid\beta).
\]
\end{definition}

A fully bound relation instance is a proposition, not a warranted fact:
\[
\boxed{
\text{relation schema}
\to
\text{partial binding}
\to
\text{question}
\to
\text{complete binding}
\to
\text{proposition}.
}
\]

\subsection{One relation generates many questions}

For
\[
\operatorname{CompRel}(f,g,h)\iff h=g\circ f,
\]
the same schema generates:
\[
?h[\operatorname{CompRel}(f,g,h)]
\]
for a composite,
\[
?f[\operatorname{CompRel}(f,g,h)]
\]
for a predecessor/factor, and
\[
?g[\operatorname{CompRel}(f,g,h)]
\]
for a continuation/factor.

No independent ``factor question'' primitive is required.

\section{Discharge authority}

\begin{definition}[Discharge mode]
Every executable open port is assigned a compiler-enforced discharge mode:
\[
\boxed{
\operatorname{mode}(q,i)
\in
\{
\mathsf{Pure},
\mathsf{Generate},
\mathsf{Probe},
\mathsf{Check},
\mathsf{Warrant}
\}.
}
\]

\begin{itemize}
\item \(\mathsf{Pure}\): determined by already standing data and registered deterministic computation;
\item \(\mathsf{Generate}\): a semantic generator may propose a provisional filling;
\item \(\mathsf{Probe}\): only an actualized interaction return may fill the port;
\item \(\mathsf{Check}\): an independently admitted checker must discharge the port;
\item \(\mathsf{Warrant}\): only the standing/support policy may discharge the port.
\end{itemize}
\end{definition}

The same query grammar therefore supports different lawful routes without allowing syntax or an LLM to confer authority on itself.

\begin{law}[Authority preservation]
If a port has mode \(\mathsf{Probe}\), \(\mathsf{Check}\), or \(\mathsf{Warrant}\), compilation to a merely generative LLM completion is ill-typed as a discharge of that port.
\end{law}

\section{Question composition}

\subsection{Dependent composition}

For
\[
q_1=?x:A[\phi(x)]
\]
and
\[
q_2(x)=?y:B(x)[\psi(x,y)],
\]
define
\[
\boxed{
q_1\bind_{\!Q} q_2
=
?(x:A,y:B(x))[
\phi(x)\land\psi(x,y)
].
}
\]

\subsection{Independent combination}

For independent ports,
\[
\boxed{
q_1\otimes_Q q_2
=
?(x:A,y:B)[
\phi(x)\land\psi(y)
].
}
\]

\subsection{Refinement by guard}

\[
\boxed{
q\upharpoonright\psi
=
?x:A[
\phi(x)\land\psi(x)
].
}
\]

\subsection{Answer substitution}

\[
\boxed{
\operatorname{Plug}(?x:A[\phi],a)
=
\phi[a/x].
}
\]

For multi-port questions, substitution is simultaneous over the returned assignment.

\section{Question-conditioned discrimination}

A question does not need to induce a globally binary ontology.  It induces a local answer profile over a live carrier.

\begin{definition}[Question support profile]
Let \(X_q\) be a live carrier to which \(q\) applies.  Define
\[
C_q\subseteq X_q\times A(q)
\]
where \(C_q(x,a)\) means that complete answer \(a\) is semantically compatible with \(x\) under the declared binding, scope, and grain.

Define
\[
\sigma_q(x)
=
\{a\in A(q):C_q(x,a)\}.
\]
\end{definition}

\begin{definition}[Local question equivalence]
\[
\boxed{
x\sim_q y
\iff
\sigma_q(x)=\sigma_q(y).
}
\]
This equivalence records what the question itself cannot distinguish.  It is local and must not be identified with protected global behavioral equivalence.
\end{definition}

When \(C_q\) is functional, write
\[
\pi_q:X_q\to A(q)
\]
and recover
\[
x\sim_q y\iff\pi_q(x)=\pi_q(y).
\]

\section{Question refinement algebra}

\begin{definition}[Precision preorder]
For questions \(q,r\) on a common live carrier,
\[
\boxed{
q\preceq_{\mathrm{prec}}r
}
\]
when \(r\)'s answer profile determines \(q\)'s answer profile.  In the functional specialization this means:
\[
\boxed{
\exists c:A(r)\to A(q)
\quad
\pi_q=c\circ\pi_r.
}
\]
Equivalently,
\[
\Ker(\pi_r)\subseteq\Ker(\pi_q).
\]
Thus \(r\) is at least as discriminating as \(q\).
\end{definition}

\begin{definition}[Joint refinement]
For functional question coordinates,
\[
\pi_{q\otimes r}(x)
=
(\pi_q(x),\pi_r(x)),
\]
hence
\[
\boxed{
\Ker(\pi_{q\otimes r})
=
\Ker(\pi_q)\cap\Ker(\pi_r).
}
\]
\end{definition}

\begin{definition}[Active representation refinement]
Let
\[
\eta_t:X\to V_t
\]
be the current active question-conditioned representation.  Adding a nonredundant question coordinate gives:
\[
\boxed{
\eta_{t+1}
=
\langle\eta_t,\pi_q\rangle.
}
\]
Therefore:
\[
\boxed{
\Ker(\eta_{t+1})
=
\Ker(\eta_t)\cap\Ker(\pi_q).
}
\]
\end{definition}

This equation is the exact form of the statement that a question expands semantic resolution along a selected distinction.

\begin{proposition}[Question redundancy]
A functional question \(q\) adds no active discriminatory coordinate exactly when there exists
\[
\bar\pi_q:V_t\to A(q)
\]
such that
\[
\boxed{
\pi_q=\bar\pi_q\circ\eta_t.
}
\]
Equivalently,
\[
\Ker(\eta_t)\subseteq\Ker(\pi_q).
\]
\end{proposition}

\begin{law}[Precision is not improvement]
Question refinement and representation improvement are different orders.  More precision is not automatically better.  A coarser representation may be a strict improvement when it preserves all protected behavior at lower cost or with better robustness.
\end{law}


\section{Sections, holes, fibers, and regenerative abstraction}

\subsection{Partial application as a relational section}

\begin{definition}[Section]
Let
\[
\rho:X\times Y\rightsquigarrow C
\]
with semantics \(\rho(x,y,c)\).  For fixed \(y:Y\), define the residual \(X\)-section
\[
\boxed{
\operatorname{Sec}^{X}_{\rho}(y):X\rightsquigarrow C
}
\]
by
\[
\operatorname{Sec}^{X}_{\rho}(y)(x,c)
\iff
\rho(x,y,c).
\]
For fixed \(x:X\), define the distinct \(Y\)-section
\[
\boxed{
\operatorname{Sec}^{Y}_{\rho}(x):Y\rightsquigarrow C
}
\]
by
\[
\operatorname{Sec}^{Y}_{\rho}(x)(y,c)
\iff
\rho(x,y,c).
\]
These are left and right residual views of the same relation and are not generally symmetric.
\end{definition}

\begin{definition}[Solution fiber]
For target output region \(T\subseteq C\), define
\[
\boxed{
\operatorname{Sol}^{X}_{\rho}(y;T)
=
\{x:X:\exists c\in T.\,\rho(x,y,c)\}.
}
\]
Equivalently,
\[
\operatorname{Sol}^{X}_{\rho}(y;T)
=
\left(\operatorname{Sec}^{X}_{\rho}(y)\right)^{-1}[T].
\]
For predicate-valued \(\rho:X\times Y\to\Prop\),
\[
\operatorname{Sol}^{X}_{\rho}(y)
=
\{x:X:\rho(x,y)\}.
\]
\end{definition}

\begin{law}[Question as structured hole]
For a relation whose \(Y\)-ports are bound and \(X\)-ports exposed,
\[
\boxed{
?_x\rho[x,y]
\quad\text{has completion fiber}\quad
\operatorname{Sol}^{X}_{\rho}(y).
}
\]
Thus a question is a structured hole in a preserved relation: remove a filling, retain the relation through which it mattered, and ask what can lawfully refill the open port.
\end{law}

\begin{definition}[Relational abstraction]
Abstraction of \(x\) from a relation web does not mean deletion of every trace of \(x\).  It means replacing the filling by an open port while preserving the typed relations constraining that port.  Write
\[
\boxed{
\operatorname{Hole}_x(W)
}
\]
for the web obtained by exposing the occurrences of \(x\) that the abstraction licence requires to remain regeneratively available.
\end{definition}

\begin{law}[Abstraction by removal]
\begin{center}
\fbox{\parbox{0.86\linewidth}{\centering\bfseries
A REGENERATIVE ABSTRACTION REMOVES A FILLING WHILE PRESERVING THE RELATIONS THAT CONSTRAIN WHAT MAY LAWFULLY REFILL IT.
}}
\end{center}
A destructive deletion that erases consequential constraints is not the same operation.
\end{law}

\subsection{Relational webs and indexed meets}

\begin{definition}[Solution field of a web]
Let \(W=\{\rho_i\}_{i\in I}\) be a family of residual relations constraining an open port of type \(X\).  After all other relevant ports are bound, define
\[
\boxed{
\operatorname{Sol}^{X}_{W}
=
\bigcap_{i\in I}\operatorname{Sol}^{X}_{\rho_i}.
}
\]
\end{definition}

\begin{theorem}[Indexed-meet refinement]
If \(W\subseteq W'\), then
\[
\boxed{
\operatorname{Sol}^{X}_{W'}
\subseteq
\operatorname{Sol}^{X}_{W}.
}
\]
Adding constraints can only refine the lawful filler field; removing constraints can only enlarge it.
\end{theorem}

\begin{proof}
The solution field is an intersection indexed by the relation family.  Every member of the larger intersection belongs to every set in the smaller family.
\end{proof}

\begin{remark}
The same order-theoretic shape appears in protected horizon refinement:
\[
\equiv_{\mathcal H}
=
\bigcap_{K\in\mathcal H}\ker(\Con_K).
\]
Enlarging the discriminator family refines equivalence for the same reason: both constructions are indexed meets.
\end{remark}

\subsection{Forced properties and near-miss fields}

\begin{definition}[Property image of a hole]
For property map \(p:X\to P\),
\[
\boxed{
\operatorname{Val}_{p,W}
=
p[\operatorname{Sol}^{X}_{W}].
}
\]
A property value \(v\) is forced by the residual web when
\[
\boxed{
\operatorname{Val}_{p,W}=\{v\}.
}
\]
\end{definition}

Literal uniqueness of the filler, protected uniqueness of the filler class, and uniqueness of a property are distinct:
\[
|\operatorname{Sol}_W|=1,
\qquad
|\operatorname{Sol}_W/{\equiv_{\mathcal H}}|=1,
\qquad
|p[\operatorname{Sol}_W]|=1.
\]
None is silently substituted for another.

\begin{definition}[Protected determination]
\[
\boxed{
Determines_{\mathcal H}(W,x)
\iff
\operatorname{Sol}^{X}_{W}/{\equiv_{\mathcal H}}
=
\{[x]_{\mathcal H}\}.
}
\]
The residual web need not recover a unique literal representation if every surviving filler has the same protected consequences.
\end{definition}

\begin{theorem}[Exact determination through a signature]\label{thm:determine-through}
Let \(\sigma:X\to S\) be an exact deterministic available signature and let
\(\chi:X\to C\) be a protected target signature over the same declared scope,
applicability, grain, binding, and horizon.  Then
\[
\boxed{
\chi=h\circ\sigma\ \text{ for some }h:\sigma[X]\to C
\quad\Longleftrightarrow\quad
\ker\sigma\subseteq\ker\chi.
}
\]
Thus exact recovery, exact consequence-preserving compression, and other exact
determination-through questions may reuse this factorization test when their
signatures satisfy these hypotheses.  Working, partial, nondeterministic, or
incompletely covered signatures require an explicit weaker contract and may return
\(\mathsf{Unknown}\); they do not inherit exactness from this theorem.
\end{theorem}

\begin{proof}
If \(\chi=h\circ\sigma\), equality under \(\sigma\) implies equality under \(\chi\).
Conversely, kernel inclusion makes \(h(\sigma(x))=\chi(x)\) well defined on the image
of \(\sigma\).
\end{proof}

\begin{definition}[Residual ambiguity]
When more than one protected filler class survives,
\[
\left|\operatorname{Sol}^{X}_{W}/{\equiv_{\mathcal H}}\right|>1,
\]
the ambiguity is itself an inquiry obligation.  A canonical separator question is
\[
\boxed{
?\rho[
\rho\text{ separates at least two surviving protected filler classes}
].
}
\]
Adding a successful separator \(\rho\) gives \(W'=W\cup\{\rho\}\) and strictly refines the solution field.
\end{definition}

\subsection{Internal hom and representability as stronger bindings}

If a binding provides a monoidal closed category and \(\rho:X\otimes Y\to C\), currying provides
\[
\Lambda_X\rho:Y\to[X,C],
\qquad
\Lambda_Y\rho:X\to[Y,C].
\]
For fixed \(y\), \(\Lambda_X\rho(y)\) is an internal-hom realization of the same residual information as the \(X\)-section.  Internal hom is therefore binding-supplied stronger structure; the relational section/fiber construction does not require it.

A still stronger representability claim is available when a hole/context induces a functor \(F_H:\mathcal C^{op}\to\mathbf{Set}\) naturally isomorphic to \(\operatorname{Hom}(-,x)\):
\[
\boxed{
F_H(-)\cong\operatorname{Hom}_{\mathcal C}(-,x).
}
\]
This determines \(x\) up to the relevant categorical isomorphism by representability.  A singleton local solution fiber does not by itself establish this stronger claim.

\section{Representation invention and discriminator generation}

\begin{definition}[Representation defect]
Let \(\eta:X\to S\) be the current representation and \(\kappa:X\rightsquigarrow K\) a protected consequence relation.  A witnessed representation defect is a pair \(x,y\) such that
\[
\boxed{
\eta(x)=\eta(y)
\quad\land\quad
x\not\equiv_{\mathcal H}y.
}
\]
In the functional consequence specialization, it suffices to witness \(\kappa(x)\neq\kappa(y)\).
\end{definition}

\begin{law}[Representation insufficiency]
A consequence difference inside one current representation cell is an obligation to search for a context, representation, grain, probe, or binding that refines that cell.  Failure of the current repertoire to produce such a separator is \(\mathsf{Unknown}\), not proof that the underlying forms are equivalent.
\end{law}

\begin{definition}[Separating context question]
For context family \(\mathcal C\),
\[
\boxed{
?C[
C\in\mathcal C
\land
\Con(C[x])\neq\Con(C[y])
].
}
\]
If no represented context succeeds, inquiry may open the higher-order context-extension question
\[
?C'[
C'\notin\mathcal C
\land
\Con(C'[x])\neq\Con(C'[y])
].
\]
\end{definition}

\begin{definition}[Representation question]
A separating representation question is
\[
\boxed{
?\eta'[
\eta'(x)\neq\eta'(y)
\land
\eta'\text{ is relevant to the protected consequence}
].
}
\]
A representation is sufficient for protected consequence \(\kappa\) over scope \(U\) when \(\kappa|_U\) factors through it:
\[
\boxed{
\exists \bar\kappa.\quad
\kappa|_U=\bar\kappa\circ\eta'|_U.
}
\]
Where exact functional factorization is unavailable, use the protected sufficiency condition
\[
\eta'(x)=\eta'(y)\Longrightarrow x\equiv_{\mathcal H}y.
\]
\end{definition}

\begin{definition}[Grain question]
\[
\boxed{
?g[
\eta_g(x)\neq\eta_g(y)
\land
\eta_g\text{ exposes a protected consequence difference}
].
}
\]
The desired grain is not maximally fine by default.  After a separating grain is found, inquiry may seek the coarsest representation that still preserves the protected distinction.
\end{definition}

\begin{definition}[Probe/tool invention question]
For current probe basis \(\mathcal P\), a new probe candidate asks
\[
\boxed{
?p'[
\forall p\in\mathcal P.\ p(x)=p(y)
\land
p'(x)\neq p'(y)
\land
x\not\equiv_{\mathcal H}y
].
}
\]
A microscope, assay, sensor, transform, debugger, benchmark, experiment, or proof procedure may instantiate this role when its return provides the required typed discrimination.
\end{definition}

\begin{law}[Tool construction as probe-basis extension]
\begin{center}
\fbox{\parbox{0.86\linewidth}{\centering\bfseries
INSTRUMENT INVENTION IS AN IMPLEMENTED EXTENSION OF THE PROBE/REPRESENTATION BASIS THAT EXPOSES A PREVIOUSLY HIDDEN CONSEQUENTIAL DISTINCTION.
}}
\end{center}
A generated tool proposal receives no semantic authority until actual use and independent checking warrant the discrimination it claims to expose.
\end{law}

\begin{definition}[Representation-gap localization]
A failure to expose a required distinction may be localized, without making the labels primitive, as:
\[
\begin{gathered}
\mathsf{CONTEXT\_GAP},\quad
\mathsf{REPRESENTATION\_GAP},\quad
\mathsf{GRAIN\_GAP},\\
\mathsf{PROBE\_GAP},\quad
\mathsf{LANGUAGE\_GAP},\quad
\mathsf{BINDING\_GAP}.
\end{gathered}
\]
Each name records which factor in the chain
\[
\text{difference}\to\text{context}\to\text{representation}\to\text{probe}\to\text{return}\to\text{consequence}
\]
currently lacks a lawful realization.
\end{definition}

\part{Positive Negation, Reciprocal Return, and Occlusion Calculus}

\section{Candidate-boundary schemas}

\begin{definition}[Typed distinction schema]\label{def:distinction-schema}
A local typed distinction schema is
\[
\boxed{D=(X,Y,B_D,\pi_X,\pi_Y,\Gamma_D)}
\]
where \(X,Y\) are candidate side carriers, \(B_D\) is a boundary-point carrier,
\[
\pi_X:B_D\to X,
\qquad
\pi_Y:B_D\to Y,
\]
and \(\Gamma_D\) is a binding-supplied downstream compatibility relation for a dependently generated reciprocal occurrence.  The schema does not contain determination presentations, negation uses, seeds, or return fibers; those belong to an occurrence of inquiry over \(D\).
\end{definition}

For \(z:B_D\), write \(x_z=\pi_X(z)\) and \(y_z=\pi_Y(z)\).  The orientation labels \(+\) and \(-\) may be used for \(X\) and \(Y\), but carry no truth value or privileged status.

\begin{definition}[Candidate boundary incidence]\label{def:boundary-incidence}
\[
\boxed{
\partial_D(x,y)
\iff
\exists z:B_D.\,\pi_X(z)=x\land\pi_Y(z)=y.
}
\]
This incidence makes \(x\) and \(y\) available in one candidate chart.  It is not a departure witness, a negation relation, a crossing, or a completed sixfold occurrence.
\end{definition}

\begin{definition}[Boundary-point profile]
\[
\Prof_D(x)=\{z:B_D:\pi_X(z)=x\},
\qquad
x_1\equiv_D^{\mathrm{prof}}x_2
\iff
\Prof_D(x_1)=\Prof_D(x_2).
\]
Profile equality remains a candidate-incidence observation and does not define contextual exteriority.
\end{definition}

\begin{law}[Boundary points remain regenerative candidate forms]\label{law:boundary-point-regeneration}
A boundary point \(z\) may be reified and at its declared grain regenerates
\[
z\rightsquigarrow(D,\pi_X(z),\pi_Y(z)).
\]
This preserves candidate-chart ancestry only.  When a reciprocal occurrence is protected, its determination presentations, use identities, fibers, recovery, seed, events, and residuals must be recovered separately; none is inferred from \(z\).
\end{law}

\section{Live determination and positive departure}

\begin{definition}[Determination presentation]\label{def:determination-presentation}
For distinction \(D\), orientation \(\alpha\in\{X,Y\}\), and live source \(s\), a determination presentation
\[
\boxed{W_D^\alpha(s)}
\]
is an explicit versioned relational web supported as constitutive of the standing source determination, indexed by declared scope, applicability, grain, protected horizon, binding, and provenance.  It is not automatically every retained fact, every standing relation mentioning \(s\), every protected continuation, or a globally unique essence.  Until a stronger admission/minimization law is warranted, the canonical reversible presentation is the support/dependency web of the specific standing claim occupying the source role.  Regenerative minimization must retain that predecessor presentation as ancestry.
\end{definition}

\begin{definition}[Positive departure witness]\label{def:departure-witness}
Let \(s:A\) be the live source under \(W=W_D^\alpha(s)\), and let \(e:E\) be a candidate.  A positive departure witness consists of typed represented observations
\[
d_A:A\rightsquigarrow C_A,
\qquad
d_E:E\rightsquigarrow C_E,
\]
supported fillings \(d_A(s,a)\) and \(d_E(e,b)\), and a standing incompatibility use \(a\perp b\) with \(\perp\hookrightarrow C_A\times C_E\), such that the source observation is relevant to \(W\).  A binding-native direct incompatibility may discharge the same role without a shared observation codomain.  Write
\[
\boxed{\Depart_{D,W}^{\alpha}(s,e)}
\]
when such a non-circular certificate exists at the declared authority, scope, applicability, and grain.
\end{definition}

\begin{law}[Departure is positive and determination-relative]\label{law:departure-relative}
Failure of equality, search, retrieval, generation, or proof does not establish \(\Depart\).  Incomplete observation remains \(\mathsf{Unknown}\), neither interior nor exterior.  Moreover,
\[
\Depart_{D,W}^{\alpha}(s,e)\not\Rightarrow s\not\eqH e,
\qquad
s\not\eqH e\not\Rightarrow\Depart_{D,W}^{\alpha}(s,e).
\]
A positively exterior candidate may remain a useful protected near-departure, while a protected difference may depend on a relation absent from the current presentation.
\end{law}

\begin{definition}[Derived boundary crossing]\label{def:boundary-crossing}
\[
\boxed{
\mathsf{BoundaryCross}
=
\mathsf{DepartureWitness}
+
\mathsf{Traversal/SuccessionProvenance}.
}
\]
Positive exteriority need not include an observed crossing path, and a candidate boundary projection alone supplies neither term.
\end{definition}

\section{Typed negation uses and tagged frontiers}

\begin{definition}[Relation use and negation use]\label{def:negation-use}
A relation use is an immutable occurrence of a typed represented relation with explicit relation identity, bindings, orientation, scope, applicability, grain, support, and warrant.  An oriented typed negation use \(u\) over \(D\) additionally carries
\[
N_u\hookrightarrow A\times E_u,
\qquad
W_D^\alpha(s),
\qquad
\mathsf{coverage}_{sem}(u),
\]
and a soundness derivation mapping every admitted exact incidence \(N_u(s,e)\) to a non-circular witness of \(\Depart_{D,W}^{\alpha}(s,e)\).  Typed negation is thus a supported role played by an ordinary oriented relation, not Boolean complement or a new ontological carrier.  One oriented use never synthesizes its reverse.
\end{definition}

\begin{definition}[Positive-negation filling]\label{def:positive-negation}
For admitted use \(u\), define
\[
\NegField_u(s)=N_u[s]=\{e:E_u:N_u(s,e)\}.
\]
An occurrence
\[
\boxed{\PosNeg(u;s,e,w)}
\]
records \(e\in\NegField_u(s)\), its departure witness \(w\), and the use identity.  \(\PosNeg\) names this role, not an additional carrier.  The distinctions
\[
\neg s\neq N_u\neq N_u[s]\neq e
\]
are mandatory.
\end{definition}

\begin{definition}[Semantic and execution coverage]\label{def:negation-coverage}
Every negation use declares semantic coverage such as exact exhaustive, certified partial, or open, together with its applicability domain and certificate.  Separately, an occurrence declares what portion of that semantic field was materialized or executed.  Therefore
\[
\boxed{\mathsf{coverage}_{sem}(u)\neq\mathsf{coverage}_{exec}(u,\omega).}
\]
An empty exact exhaustive field differs from an empty unsearched field.  No breaker found under partial semantic or execution coverage establishes only \(\mathsf{Unknown}\), not absolute closure.
\end{definition}

\begin{definition}[Tagged negation frontier]\label{def:negation-frontier}
For the admitted use family \(\mathcal U_D^\alpha(s)\), the live exterior frontier is the tagged dependent sum
\[
\boxed{
\mathcal N_D^\alpha(s)
=
\sum_{u\in\mathcal U_D^\alpha(s)}
\{u\}\times\NegField_u(s).
}
\]
If the same exterior is reached through two uses, both occurrences remain.  Heterogeneous target carriers remain typed by the dependent sum.  Exact collective coverage requires an explicit cover certificate; partial members do not become exhaustive merely by union, and an untagged union is unlawful when return provenance is protected.
\end{definition}

Ordinary intersection, union, and composition may generate candidate relations, but do not confer negation authority.  They introduce no exterior candidate absent from their admitted members unless the composed relation independently earns a sound negation-use warrant.  Open use families are traversed under the ordinary resource-bounded fairness rules.

\section{Reverse-section return and protected recovery}

\begin{definition}[Use-specific return fiber]\label{def:return-fiber}
For an admitted incidence \(N_u(s,e)\), pure return is the reverse section
\[
\boxed{
\RetField_u(e)
=
N_u^{-1}[e]
=
\{s':A:N_u(s',e)\}.
}
\]
The source \(s\) necessarily belongs to this fiber, but need not be uniquely recovered.  A selected return role \(R\in\RetField_u(e)\) is not the fiber.  Exact source closure requires
\[
\RetField_u(e)/{\equiv_{\mathcal H}}=\{[s]_{\mathcal H}\}
\]
under the declared semantic and execution coverage.
\end{definition}

\begin{definition}[Protected recovery]\label{def:protected-recovery}
For a protected represented observation \(r:A\rightsquigarrow C\), define its protected signature at \(a\) by
\[
\Sig_{r,\mathcal H}(a)
=
\{[c]_{\mathcal H}:r(a,c)\}.
\]
A return fiber \(F\subseteq A\) locally recovers \(r\) when
\[
\boxed{
\Recover_{\mathcal H}(r,F)
\iff
\forall a,a'\in F.\,
\Sig_{r,\mathcal H}(a)=\Sig_{r,\mathcal H}(a').
}
\]
Recovery is thus protected determination by the return fiber, not a scalar similarity and not a new primitive.  Raw relation differences ignored by \(\mathcal H\) do not defeat recovery; a protected signature difference does.
\end{definition}

\begin{definition}[Source-web recovery profile]\label{def:web-recovery}
For a source presentation \(W=W_D^\alpha(s)\), admitted use \(u\), and exterior \(e\), define
\[
\boxed{
\mathsf{Recov}_{W,\mathcal H}(s,e;u)
=
\{r\in W:\Recover_{\mathcal H}(r,\RetField_u(e))\}.
}
\]
An exterior may recover some constitutive source relations without regenerating the entire protected source class.  This set-valued profile is semantic; no universal scalar recovery percentage is introduced.
\end{definition}

\begin{definition}[Three-valued recovery/loss profile]\label{def:recovery-loss-profile}
For \(F=\RetField_u(e)\) and \(r\in W\), a positive non-recovery witness is a pair
\(a,a'\in F\) such that
\[
\Sig_{r,\mathcal H}(a)\neq\Sig_{r,\mathcal H}(a').
\]
At an executable occurrence \(\omega\), retain
\[
\boxed{
\mathsf{RecoveryStatus}_{\omega}(r,e;u)
\in
\{
\mathsf{Recovered}(c),
\mathsf{NotRecovered}(w),
\mathsf{Unknown}(\delta)
\}.
}
\]
Here \(c\) is a recovery certificate, \(w\) is a protected signature-difference
witness, and \(\delta\) records missing evidence, coverage, or decision capability.
Consequently
\[
W\setminus\mathsf{Recov}_{W,\mathcal H}(s,e;u)
\]
may be called an irrecoverable residue only under a declared exact decision/coverage
certificate.  Otherwise that complement unlawfully merges witnessed loss with
\(\mathsf{Unknown}\).
\end{definition}

\begin{definition}[Derived constitutive characterization]\label{def:characterization-view}
A coverage-indexed characterization view of a source \(s\) is reconstructed from its
supported determination presentation, positively certified admitted variation that
remains in the protected source class, tagged negation frontier, and the three-valued
recovery/loss profiles of admitted exterior occurrences.  It retains scope,
applicability, grain, binding, horizon, coverage, provenance, and open residuals.

This view may serve an externally declared goal horizon or a recursively developed
constitutive inquiry horizon.  A recursively generated discriminator enters the latter
only through ordinary independent actualization, checking, warrant, and admission; it
cannot protect or warrant itself merely because the characterization program generated
it.  The view is relative and versioned, not an absolute essence, new semantic
primitive, or authoritative history species.  A compressed characterization is lawful
only under inquiry-regenerative sufficiency.
\end{definition}

\begin{definition}[Near-negation recovery order]\label{def:near-negation-order}
For two admitted tagged exterior occurrences \((u_1,e_1)\) and \((u_2,e_2)\), write
\[
(u_1,e_1)\succeq_{W,\mathcal H}(u_2,e_2)
\iff
\mathsf{Recov}_{W,\mathcal H}(s,e_2;u_2)
\subseteq
\mathsf{Recov}_{W,\mathcal H}(s,e_1;u_1).
\]
This is only a protected recovery preorder.  Cost, risk, and time may be combined by an explicit product or Pareto order, never by a primitive scalar negation distance.
\end{definition}

\begin{definition}[Negation-family return signature]\label{def:family-signature}
For each admitted negation use \(u_i\), let
\[
\bar\sigma_i(s)
=
(\mathsf{applicability}_i(s),N_{u_i}[s])
\]
with applicability retained whenever protected.  The tagged family signature is
\[
\boxed{
\sigma_{\mathcal U}(s)=\prod_i\bar\sigma_i(s).
}
\]
For deterministic exact signatures,
\[
\boxed{
\ker\sigma_{\mathcal U}=\bigcap_i\ker\bar\sigma_i.
}
\]
Adding a lawful coordinate may refine but cannot coarsen the exact family observational partition.  Information accumulates through this product, not through an untagged union of negation relations.
\end{definition}

\begin{law}[Family information is not joint actuality]\label{law:family-jointness}
The product \(\sigma_{\mathcal U}\) is a tagged derived information view over its
declared family.  Component signatures supported in different occurrences or
contexts do not thereby form one jointly actualizable return.  Any binding that
interprets \((\bar\sigma_1,\ldots,\bar\sigma_n)\) as a single realized composite
observation must supply supported co-applicability and joint-realizability evidence,
with the component provenance and contexts preserved.  Absent that evidence, family
factorization may combine retained information but may not manufacture a composite
actual event.
\end{law}

\begin{definition}[Schema and family recovery]\label{def:family-recovery}
A protected source observation \(\chi_{r,\mathcal H}\) is recoverable from the admitted negation-use family when
\[
\boxed{
\ker\sigma_{\mathcal U}
\subseteq
\ker\chi_{r,\mathcal H}.
}
\]
Equivalently, where factorization is available, there is an \(h\) such that
\[
\boxed{
\chi_{r,\mathcal H}=h\circ\sigma_{\mathcal U}.
}
\]
Family recovery may therefore succeed when no member signature alone recovers the observation; the canonical three-state breaker witnesses this joint gain.  Occurrence-local recovery remains fixed by the use, fiber, and evidence available at that historical occurrence and is not retroactively strengthened by a later learned use.
This is the specialization of \cref{thm:determine-through} to the exact deterministic
family signature; \cref{law:family-jointness} still prevents informational product
recovery from being read as proof of simultaneous actualizability.
\end{definition}

\begin{theorem}[Compatible monotone return cannot revise a determined source]\label{thm:monotone-return}
If a web \(W\) already determines exactly \([s]_{\mathcal H}\), and a compatible constraint \(C\) admits \(s\), then
\[
\operatorname{Sol}(W\land C)
=
\operatorname{Sol}(W)\cap\operatorname{Sol}(C)
\]
still has only \([s]_{\mathcal H}\) as its protected class.  Hence adding a compatible exterior constraint cannot transform an already determined source into a protectedly different source.  A state-changing redetermination must expose a warranted revision route of the kind required by \cref{law:return-not-reconciliation}.
\end{theorem}

\section{Dependent reciprocal occurrence and sixfold view}

\begin{definition}[Seed and reorientation]\label{def:reciprocal-seed}
After \(\PosNeg(u_X;S_X,O_X,w_X)\), a supported typed seed relation
\[
\sigma_{X\to Y}:E_{u_X}\rightsquigarrow Y
\]
selects or determines \(S_Y\) from \(O_X\).  Identity may be used only when the carriers and intended roles justify it.  \(O_X\) and \(S_Y\) remain distinct roles and the seed's support and provenance remain explicit.
\end{definition}

\begin{definition}[Dependent sixfold occurrence]\label{def:dependent-sixfold}
A reciprocal occurrence \(\omega\) over \(D\) is generated in the following dependency order:
\[
\boxed{
\begin{aligned}
W_D^X(S_X),\ u_X
&\longrightarrow
\PosNeg(u_X;S_X,O_X,w_X)
\longrightarrow
F_X=\RetField_{u_X}(O_X)
\ni R_X,\\
O_X
&\xrightarrow{\sigma_{X\to Y}}
S_Y,\\
W_D^Y(S_Y),\ u_Y
&\longrightarrow
\PosNeg(u_Y;S_Y,O_Y,w_Y)
\longrightarrow
F_Y=\RetField_{u_Y}(O_Y)
\ni R_Y.
\end{aligned}
}
\]
The derived role view is
\[
\boxed{\Xi_D(\omega)=(S_X,O_X,R_X;S_Y,O_Y,R_Y).}
\]
The occurrence also carries \(u_X,u_Y\), both coverage records, \(F_X,F_Y\), local and family recovery results, seed support, boundary-point provenance when one exists, ordinary event references, and unresolved residuals.  It is reconstructed from ordinary questions, fibers, support, and events; it is not a separate authoritative history species.
\end{definition}

\begin{law}[Roles remain distinct and asymmetric]\label{law:sixfold-role-distinction}
The six roles in \cref{def:dependent-sixfold} are dependent rather than independent.  In general
\[
R_X\neq O_Y,
\qquad
R_Y\neq O_X.
\]
Stable return in one orientation may coexist with unstable return in the other, and one-way negation does not imply reciprocal negation.  Exact reciprocal closure is relative to both negation semantic coverage and occurrence execution coverage.
\end{law}

\begin{definition}[Protected reciprocal return stability]\label{def:reciprocal-return-stability}
For a completed occurrence \(\omega\), protected reciprocal return stability holds when
\[
F_X/{\equiv_{\mathcal H}}=\{[S_X]_{\mathcal H}\},
\qquad
F_Y/{\equiv_{\mathcal H}}=\{[S_Y]_{\mathcal H}\},
\]
and every declared protected recovery obligation succeeds under the occurrence's semantic and execution coverage.  Literal uniqueness is not required unless the horizon distinguishes literal representatives.  This is the successor form of protected mutual determination; it is a property of the two same-use fibers, not of boundary incidence or \(\Gamma_D\) alone.
\end{definition}

\begin{law}[Gamma is downstream]\label{law:gamma-downstream}
Only after the dependent roles, use identities, fibers, recovery results, and seed have been supplied may the binding check
\[
\boxed{\Gamma_D(\omega).}
\]
\(\Gamma_D\) cannot manufacture a missing role, negation use, departure witness, return filling, or recovery result.  Compatibility failure becomes an ordinary protected residual rather than retroactive evidence that an ungenerated role existed.
\end{law}

The principal reciprocal residuals are return ambiguity, protected recovery failure, seed mismatch, reciprocal mismatch, and downstream compatibility failure.  They enter the generic separator problem.  If a witnessed protected distinction has no expression in the admitted representation/probe language, the result is a representation or binding gap rather than endless search within the same language.

\section{Orientation, recursion, and re-entry}

\begin{definition}[Oriented double cover]\label{def:double-cover}
\[
\widetilde B_D=B_D\times\{+,-\},
\qquad
p(z,\alpha)=z,
\]
with involution \(\iota(z,+)=(z,-)\) and \(\iota(z,-)=(z,+)\).  The cover records which orientation is under inquiry.  Applying \(\iota\) changes orientation only; it does not construct \(u_Y\), prove departure, or execute the reciprocal branch.
\end{definition}

\begin{definition}[Distinction reification]
A boundary point or completed reciprocal occurrence may be reified without erasing its type or provenance:
\[
\operatorname{QuoteDist}_D:
B_D+\mathsf{RecipOccurrence}_D
\to\mathsf{Form}_{\B}.
\]
It may then occupy a side of another typed distinction.  A schematic derived former
\[
\Delta:\mathsf{Form}_{\B}\times\mathsf{Form}_{\B}
\rightharpoonup\mathsf{Form}_{\B}
\]
is lawful only when its typed distinction schema, scope, presentation, and support obligations are discharged.  Construction ancestry is one projection of the global relational graph; constituents remain independently recomposable.
\end{definition}

\section{Candidate polarity and traversal}

Given candidate incidence \(\partial_D\), one may derive the face-level polarity
\[
A^{\perp_D}=\{y\in Y:\forall x\in A,\ \partial_D(x,y)\},
\qquad
E^{\perp_D}=\{x\in X:\forall y\in E,\ \partial_D(x,y)\},
\]
and \(\operatorname{cl}_D(A)=A^{\perp_D\perp_D}\).  This is a chart operation only.  It supplies neither the positive evidence of \cref{def:departure-witness} nor the use authority of \cref{def:negation-use}, and it does not replace use-specific return fibers or the dependent sixfold occurrence.

\begin{law}[Candidate profile is not double closure]
In general,
\[
\operatorname{cl}_D(\{x\})
\neq
[x]_{\equiv_D^{\mathrm{prof}}}.
\]
Candidate polarity remains a useful face-level construction but cannot erase exact boundary-profile differences or any protected occurrence-level distinction.
\end{law}

\begin{definition}[Boundary traversal]
\[
T_D:B_D\rightsquigarrow B_D,
\qquad
\tau_D:(z,\alpha)\rightsquigarrow(z',\alpha)
\iff T_D(z,z').
\]
Base traversal, orientation change, reciprocal-question execution, and binding-defined realized/domain succession are four distinct relations.  A crossing is derived only as in \cref{def:boundary-crossing}.
\end{definition}

\section{Question-conditioned occlusion}

\begin{definition}[Active view]
For question \(q\), state/configuration carrier \(S\), and grain \(g\), an active view is a relation or function
\[
\boxed{
\eta_{q,g}:S\rightsquigarrow V_{q,g}.
}
\]
It supplies the context made active for compiling and executing \(q\).
\end{definition}

Relative to \(q\) and protected horizon \(\mathcal H\), available retained structure is partitioned operationally into
\[
\boxed{
\Sigma
=
A_{q,\mathcal H}
\;\dot\cup\;
U_{q,\mathcal H}
\;\dot\cup\;
I_{q,\mathcal H}.
}
\]
Here \(A\) is established relevant, \(U\) is relevance-undetermined, and \(I\) is certified consequence-null under an explicit guard.  Only \(I\) may be aggressively occluded; \(U\) remains reachable exploration reserve.

\begin{law}[Unknown relevance]
\[
\boxed{
\text{not currently active}
\neq
\text{warranted irrelevant}.
}
\]
\end{law}

\begin{definition}[Occlusion licence]
A local view \(\eta_{q,g}\) is exact for protected local future \(\mathcal H_q\) when
\[
\eta_{q,g}(s)=\eta_{q,g}(s')
\Longrightarrow
\forall K\in\mathcal H_q,\,
\Con_K(s)=\Con_K(s').
\]
Every aggressively occluded item carries an unlock/reopening predicate.
\end{definition}

\section{Question-frame defect}

\begin{definition}[Question-frame defect]
Let \(\equiv_{\mathcal H}\) be protected equivalence on the applicable live carrier.  Define
\[
\boxed{
FD_{\mathcal H}(q)
=
\{(x,y):x\sim_q y\land x\not\equiv_{\mathcal H}y\}.
}
\]
\end{definition}

A nonempty frame defect is lawful for a local question coordinate when surrounding retained history/state preserves the omitted protected distinction.  If the question view is proposed as a complete replacement representation, require
\[
\boxed{FD_{\mathcal H}(q)=\varnothing.}
\]

\section{Square and cube as projections of dependent occurrences}

\begin{definition}[Four-cell square projection]
For a completed reciprocal occurrence \(\omega\), the traditional square is the projection
\[
\boxed{
Sq_D(\omega)
=
\begin{pmatrix}
S_X & O_X\\
O_Y & S_Y
\end{pmatrix}.
}
\]
It intentionally omits \(R_X,R_Y\), both negation-use identities, coverage, fibers, recovery, seed support, residuals, and provenance.  It is therefore a useful view, not a complete definition of the occurrence or boundary point.
\end{definition}

\begin{definition}[Reciprocal succession prism]
If realized succession carries reciprocal occurrence \(\omega_t\) to \(\omega_{t+1}\), the full derived cube/prism is
\[
\boxed{
\mathcal C_t(D)
=
(\Xi_D(\omega_t),\lambda_t,\Xi_D(\omega_{t+1})).
}
\]
The square-cube display
\[
(Sq_D(\omega_t),\lambda_t,Sq_D(\omega_{t+1}))
\]
is a projection of this richer succession object.
\end{definition}

\begin{law}[Arrangement, reciprocal orientation, and succession remain distinct]
The dependent sixfold view describes one completed reciprocal occurrence; the double cover records orientation; \(T_D\) records inquiry traversal; \(\lambda_{\mathrm{real}}\) records binding-defined realized/domain succession.  None is identified with another without an explicit bridge.
\end{law}

\part{Probe, Perception, and Executable Inquiry}

\section{The unified probe relation}

\begin{definition}[Probe semantics]
Let \(Z\) be a coupled agent/environment configuration carrier at a declared grain and let \(R_o\) be the raw return carrier for probe operator \(o\).  The denotation of \(o\) is:
\[
\boxed{
\den{o}_{\B}
\subseteq
Z\times R_o\times Z.
}
\]
Equivalently:
\[
o:Z\rightsquigarrow R_o\times Z.
\]
\end{definition}

A realized occurrence is written:
\[
\boxed{
z_t\xrightarrow[o_t]{r_t}z_{t+1}.
}
\]

The event records both the return and the successor configuration.  No separate primitive ``observation transition'' or ``action transition'' is required.

\subsection{Observation-like and action-like projections}

Let
\[
\pi_W:Z\to W
\]
project the coupled configuration to an environment/world component relevant to a binding.  An occurrence is observation-like relative to horizon \(\mathcal H\) when:
\[
\pi_W(z_t)\equiv_{\mathcal H}\pi_W(z_{t+1}),
\]
and action-like when the protected world projection differs.  The same probe may be both observation-like in one projection and action-like in another.

Therefore:
\[
\boxed{
\textbf{OBSERVATION AND ACTION ARE DERIVED ROLES OF PROBING.}
}
\]

\section{Perception}

\begin{definition}[Live compatible region]
Let \(L_t\subseteq Z\) be the currently live region of configurations compatible with retained standing information at the chosen grain.
\end{definition}

After occurrence of \(o\) with raw return \(r\), define:
\[
\boxed{
L_{t+1}^{o,r}
=
\{
z':
\exists z\in L_t.\,
\den{o}_{\B}(z,r,z')
\}.
}
\]

\begin{definition}[Perceptual discrimination]
The raw return \(r\) produces a perceptual discrimination exactly to the extent that:
\[
L_{t+1}^{o,r}
\subsetneq
\bigcup_{r'\in R_o}L_{t+1}^{o,r'}.
\]
Perception is therefore return-constrained discrimination of the live field.
\end{definition}

\begin{law}[Perception is not warrant]
A raw perceptual return is authoritative as the return that occurred through its actualization channel.  It does not thereby warrant any semantic interpretation of that return.
\end{law}

\section{Semantic question, executable probe, and occurrence}

A semantic question, its compiled operator, and an actual occurrence must remain distinct:
\[
\boxed{
q
\neq
\Compile(q)
\neq
\operatorname{Occur}(\Compile(q)).
}
\]

They are nevertheless stages of one inquiry chain:
\[
\boxed{
\text{open relation}
\longrightarrow
\text{probe/operator code}
\longrightarrow
\text{realized probe occurrence}.
}
\]

A question is therefore a \emph{probe schema} semantically; a prompt or other executable artifact is its operational realization; the event is its historical occurrence.

\section{Raw return, decoding, partial answers, and completion}

\subsection{Raw return}

For a compiled probe \(o\), the runtime receives:
\[
r:\Raw(R_o).
\]
The raw return is preserved before any decoder, interpreter, or semantic patch is allowed to operate on it.

\subsection{Decoder}

For semantic question \(q\), let \(Y_q\) be a decoder input carrier and define:
\[
\boxed{
\delta_q:
Y_q
\rightharpoonup
\mathcal P_+(A(q)).
}
\]

For a returned \(y\):
\begin{itemize}
\item \(|\delta_q(y)|=1\): a complete semantic answer is supported;
\item \(1<|\delta_q(y)|<|A(q)|\): the response partially resolves the question;
\item \(\delta_q(y)=A(q)\): the decoded return is nondiscriminating for \(q\);
\item \(\delta_q(y)\) undefined: the response is malformed or outside this decoder's declared
      decoding domain; support failure and incomplete coverage are classified only by the
      subsequent resolution judgment.
\end{itemize}

For infinite answer carriers, ``partial'' means a proper nonempty supported subset rather than cardinality comparison.

\subsection{Remaining completion field}

For supported answer set \(S\subseteq A(q)\),
\[
\boxed{
F_q(S)
=
\{x\in X_q:\sigma_q(x)\cap S\neq\varnothing\}.
}
\]

The excluded answer set is:
\[
A(q)\setminus S.
\]

Thus a response need not terminate inquiry.  It may merely contract the live boundary field and thereby generate the next question.

\section{Explicit resolution paths}

\begin{definition}[Resolution path]
Resolution paths are typed:
\[
\boxed{
\ResPath:A\leadsto B.
}
\]
The reference grammar is:
\[
\ResPath::=
\mathsf{Id}_A
\mid
\mathsf{Edge}(e)
\mid
\ResPath_2\circ\ResPath_1,
\]
where each edge is an admitted decoder, parser, checker-facing conversion, or binding-native interpretation step.
Its denotation is a typed relation
\[
\boxed{
\operatorname{Run}(p)\subseteq A\times B
\qquad(p:A\leadsto B),
}
\]
defined by
\[
\begin{aligned}
\operatorname{Run}(\mathsf{Id}_A,a,b)
&\iff a=b,\\
\operatorname{Run}(\mathsf{Edge}(e),a,b)
&\iff \den{e}(a,b),\\
\operatorname{Run}(p_2\circ p_1,a,c)
&\iff \exists b.\,
\operatorname{Run}(p_1,a,b)\land\operatorname{Run}(p_2,b,c).
\end{aligned}
\]
Equivalently,
\(\operatorname{Run}(p,a)=\{b:\operatorname{Run}(p,a,b)\}\); this set may be empty,
singleton, or non-singleton.  Undefined computation is represented by absence of a related output
together with the path's typed resolution residual, never by an untyped exception masquerading as
an empty semantic answer.
\end{definition}

For a raw return \(r:A\), a semantic completion \(a\in A(q)\) is supported through
\(\ResPath\) when:
\[
\boxed{
\operatorname{Run}(\ResPath,r,a)
\land
\Comp_q(a).
}
\]

The candidate completion set supplied to the five-way resolution classifier is therefore
\[
\boxed{
C_{q,p,r}:=
\{a\in A(q):\operatorname{Run}(p,r,a)\land\Comp_q(a)\}.
}
\]
Support, exact emptiness, undefinedness, support failure, and incomplete coverage are determined
from this set together with the path's declared domain, evidence, and coverage; cardinality alone
does not select the constructor.

\begin{law}[Resolution provenance]
The runtime must retain the path by which a raw return was decoded when protected future use may distinguish different resolution routes.
\end{law}


\section{Source inquiry programs}

The runtime effect language below is intentionally lower-level than the source language used to express portable question programs.  The source language makes answer-dependent questioning explicit and is executable in an ordinary chat when its discharge modes require no unavailable external effect.

\begin{definition}[Declared-route support judgment]\label{def:declared-route-support}
For a well-typed question \(q=?_I R[\beta]\), a nonempty candidate set
\(S\subseteq A(q)\), an exact route description \(\vartheta\), and an evidence environment
\(E\), the checked judgment
\[
\SupportedRoute(q,S;\vartheta,E)
\]
holds exactly when every member of \(S\) is a well-typed completion admitted through the declared
mode of each open port, the referenced evidence and resolution paths exist and apply at the declared
scope/grain/horizon, and no member is promoted beyond that route's authority.  The judgment retains
per-component provenance and coverage.  It may support a proper subset of the completion fiber;
unexamined or unresolved members remain \(\mathsf{Unknown}\).  A generated proposal may be admitted
as a generated branch, but it is not thereby actual, checked, warranted, or standing.
\end{definition}

\begin{definition}[First-order support witness]\label{def:support-witness}
For the same \((q,S,\vartheta,E)\), let
\[
\mathsf{SupportWitness}(q,S;\vartheta,E)
\]
be the first-order carrier of checked derivations of
\(\SupportedRoute(q,S;\vartheta,E)\).  An inhabitant records the declared coverage, a
component-indexed provenance map for every represented component of every member of \(S\), and the
checked support derivation.  Its validation judgment is
\[
\B;\Gamma\vdash
w:\mathsf{SupportWitness}(q,S;\vartheta,E).
\]
The judgment \(\SupportedRoute(q,S;\vartheta,E)\) holds exactly when this carrier is inhabited.
Thus support evidence is inspectable proof-relevant program data; it is not an erased truth value or
a self-warranting assertion.
\end{definition}

\begin{definition}[Supported answer type]\label{def:supported-answer-type}
For semantic question \(q\), let \(\mathsf{Route}(q)\) be its first-order route-description type and
let \(\mathsf{EvidenceEnv}(q,\vartheta)\) be the corresponding typed evidence-environment type.
Define
\[
\boxed{
\mathsf{SuppAns}(q)
=
\sum_{\vartheta:\mathsf{Route}(q)}
\sum_{E:\mathsf{EvidenceEnv}(q,\vartheta)}
\sum_{S:\mathcal P_+(A(q))}
\mathsf{SupportWitness}(q,S;\vartheta,E).
}
\]
Write
\[
\widehat S=\langle\vartheta,E,S,w\rangle,
\qquad
|\widehat S|=S.
\]
The projections \(\mathsf{route}(\widehat S)\), \(\mathsf{evidence}(\widehat S)\),
\(\mathsf{coverage}(\widehat S)\), \(\mathsf{componentProv}(\widehat S)\), and
\(\mathsf{support}(\widehat S)\) retain the exact checked route, evidence, declared coverage,
per-component provenance, and support witness carried by \(w\).  Equality of supported answers is
equality of this proof-carrying first-order record, not merely equality of its member projection.
A complete answer has singleton \(|\widehat S|\).  A partial answer has more than one live
completion or leaves declared components unresolved.  The record is supplied whole to its
continuation; no rule implicitly selects one member or erases its ancestry.
\end{definition}

\begin{definition}[Resolution certificate and residual carriers]\label{def:resolution-carriers}
The resolution payloads are first-order, question-indexed carriers:
\begin{itemize}
\item \(\mathsf{EmptyCert}(q)\) contains an admitted exhaustive-coverage certificate and a checked
      derivation that no completion of \(q\) is supported under that coverage;
\item \(\mathsf{ResolutionResidual}(q)\) identifies a malformed, inapplicable, or undefined
      decoding/resolution path together with the event/raw-return and version provenance checked so
      far;
\item \(\mathsf{SupportResidual}(q)\) contains decoded candidate members and the exact declared
      support obligations they fail, without treating failure as semantic negation; and
\item \(\mathsf{CoverageResidual}(q)\) identifies the uncovered completion region or unresolved
      components that prevent either support or exact emptiness.
\end{itemize}
These carriers are disjoint by constructor and retain the exact route, evidence, versions, scope,
grain, horizon, and provenance needed to check or reopen the result.
\end{definition}

\begin{definition}[Resolution outcome]\label{def:resolution-outcome}
Resolution does not fabricate a nonempty supported answer when no such record is established:
\[
\boxed{
\begin{aligned}
\mathsf{ResolutionOutcome}(q)
={}&
\mathsf{Supported}(\mathsf{SuppAns}(q))\\
&+\mathsf{ExactEmpty}(\mathsf{EmptyCert}(q))\\
&+\mathsf{Undefined}(\mathsf{ResolutionResidual}(q))\\
&+\mathsf{Unsupported}(\mathsf{SupportResidual}(q))\\
&+\mathsf{Unknown}(\mathsf{CoverageResidual}(q)).
\end{aligned}
}
\]
\(\mathsf{ExactEmpty}\) requires exhaustive admitted coverage; \(\mathsf{Undefined}\) records a
malformed or inapplicable resolution path; \(\mathsf{Unsupported}\) records a decoded candidate
that fails its declared support route; and \(\mathsf{Unknown}\) records incomplete coverage.
Provider unavailability and resource exhaustion remain the distinct operational stop states
\(\mathsf{Blocked}\) and \(\mathsf{ResourceBounded}\).  Only the
\(\mathsf{Supported}(\widehat S)\) summand may fill a source answer slot.
\end{definition}

\begin{definition}[Inquiry source grammar]
\[
\boxed{
K::=
\mathsf{Return}_I(e)
\mid
\mathsf{Ask}\{q,x,e_K\}.
}
\]
Here \(e\) and \(e_K\) are first-order expressions, and an \(\mathsf{Ask}\) stores an answer slot
\(x\) rather than a host-language callback.  Its checked continuation package is
\[
\boxed{
\kappa=\langle\Gamma,x:\mathsf{SuppAns}(q),e_K\rangle,
\qquad
\kappa(\widehat S)=e_K[\widehat S/x]\ \text{under explicit }\Gamma.
}
\]
The environment, expression, answer slot, and any named subprogram references are inspectable,
serializable program data; substitution is capture-safe.  The notation \(\mathsf{Ask}(q,\kappa)\)
used elsewhere abbreviates this record.  It never denotes an opaque host-language closure or hidden
model policy.  The continuation may construct a new relation/question from the returned distinction
rather than selecting only among a fixed textual checklist.
\end{definition}

\subsection{Source typing}

\[
\frac{\Gamma\vdash a:A}
{\Gamma\vdash\mathsf{Return}_I(a):\mathsf{IProg}(A)}
\quad(\textsc{IReturn})
\]

\[
\frac{
\Gamma\vdash q\;\mathsf{question}
\qquad
\Gamma,\widehat S:\mathsf{SuppAns}(q)\vdash\kappa(\widehat S):\mathsf{IProg}(A)
}{
\Gamma\vdash\mathsf{Ask}(q,\kappa):\mathsf{IProg}(A)
}
\quad(\textsc{Ask})
\]

\begin{definition}[Residual question generator]
A higher-order continuation may be represented by
\[
\boxed{
\mathsf{ResidualQuestion}:
(\Sigma,\AskRef,\widehat S,\mathsf{Residual})
\rightsquigarrow
q'.
}
\]
Here \(\widehat S:\mathsf{SuppAns}(\mathsf{AskQuestion}(\AskRef))\) is supplied whole and
\(\AskRef\) is the checked source \(\mathsf{Ask}\) occurrence, so the generator retains the
continuation and environment that make the route meaningful.  The residual includes surviving
filler classes, open dependencies, broken exclusions, reopened folds, representation defects, or
other consequential structure left unresolved by the answer.  The canonical next question is one
whose possible supported returns can discriminate that residual or whose discharge is explicitly
required by the standing program.
\end{definition}

Hence a rendered inquiry trace may abbreviate the semantic questions as
\[
\boxed{
q_0\xrightarrow{\AskRef_0,E_0,\ResPath_0,\widehat S_0}q_1
\xrightarrow{\AskRef_1,E_1,\ResPath_1,\widehat S_1}q_2\xrightarrow{}\cdots
}
\]
rather than a static list of instructions.  This display is only a projection of the canonical
occurrence chain in \cref{def:resolved-inquiry-occurrence}; it does not identify an event, raw
return, resolution path, or supported set with a question.

\subsection{Derived occurrence-indexed question succession}

\textbf{Status: DERIVED.}  Dynamic succession is a projection of the source program; it is not a
new transition primitive.

\begin{definition}[Checked Ask occurrence]\label{def:ask-occurrence}
For result type \(A\), let \(\mathsf{SourceConfig}_{\B}(A)\) be the first-order dependent carrier
\[
\boxed{
\mathsf{SourceConfig}_{\B}(A)
:=
\left\{
\langle K,\Gamma,b,v,p\rangle
\ \middle|\
\begin{aligned}
&K\in\mathsf{IProg}(A),\quad
 \Gamma\in\mathsf{Env}_{\B},\quad
 b\in\mathsf{BindingVersion}_{\B},\\
&v\in\mathsf{CompilerVersion}_{\B},\quad
 p\in\mathsf{Provenance}_{\B},\\
&\B;\Gamma\vdash K:\mathsf{IProg}(A)
 \land\mathsf{CheckedSourceVersions}(\B,\Gamma,b,v,p)
\end{aligned}
\right\}.
}
\]
Here \(\mathsf{CheckedSourceVersions}\) verifies that \(\Gamma\), binding version \(b\), compiler
version \(v\), and accepted provenance \(p\) are mutually applicable to the stored source program;
it is not inferred from co-location.
For \(\Sigma=\langle K,\Gamma,b,v,p\rangle:\mathsf{SourceConfig}_{\B}(A)\), let
\(\mathsf{Pos}(K)\) be the first-order structural positions of \(K\), let
\(K|_\pi\) be the subterm at \(\pi\), and let
\(\mathsf{EnvAt}(K,\Gamma,\pi)\) and \(\mathsf{ProvAt}(p,\pi)\) be the capture-safe environment and
structural provenance obtained by replaying the source syntax to that position.  These are checked
structural operations, not user-supplied labels.  Define
\[
\boxed{
\mathsf{AskOcc}_A(\Sigma)
:=
\left\{
\begin{gathered}
\langle\pi,q,x,e_K,\Gamma_\pi,b_\pi,v_\pi,p_\pi\rangle
\\[-1mm]
\mathrel{\Big|}\quad
\pi\in\mathsf{Pos}(K),\quad
 K|_\pi=\mathsf{Ask}\{q,x,e_K\},\\
\Gamma_\pi=\mathsf{EnvAt}(K,\Gamma,\pi),\quad
 b_\pi=b,\quad v_\pi=v,\quad p_\pi=\mathsf{ProvAt}(p,\pi),\\
\B;\Gamma_\pi\vdash q\ \mathsf{question},\quad
\B;\Gamma_\pi,x:\mathsf{SuppAns}(q)\vdash e_K:\mathsf{IProg}(A),\\
\mathsf{CaptureSafe}(K,\pi,x,e_K,\Gamma_\pi)
\end{gathered}
\right\}.
}
\]
Thus the formation judgment
\(\Sigma\vdash\AskRef:\mathsf{AskOcc}_A\) is derivable exactly by re-walking the source program
and checking every displayed field.  A copied position, question, slot, environment, continuation,
binding/compiler version, or provenance field does not form an occurrence.
Let \(\mathsf{Question}_{\B}(\Sigma)\) be the dependent carrier of questions appearing with such a
typing derivation in \(\Sigma\), including capture-safe answer extensions reached by pure
normalization of its subprograms.  For
\(\AskRef=\langle\pi,q,x,e_K,\Gamma_\pi,b,v,p_\pi\rangle\), write
\[
\mathsf{AskQuestion}(\AskRef)=q,
\qquad
\mathsf{AskContinuation}(\AskRef)=\kappa_{\AskRef},
\qquad
\kappa_{\AskRef}(\widehat S)=e_K[\widehat S/x]\ \text{under }\Gamma_\pi.
\]
The projections have dependent types
\[
\mathsf{AskQuestion}:\mathsf{AskOcc}_A(\Sigma)\to\mathsf{Question}_{\B}(\Sigma),
\]
\[
\mathsf{AskContinuation}(\AskRef):
\mathsf{SuppAns}(\mathsf{AskQuestion}(\AskRef))\to\mathsf{IProg}(A).
\]
Two occurrences may have the same normalized question while retaining different continuations.
\end{definition}

\begin{definition}[Head question]\label{def:head-question}
\(\HeadQ(K,q')\) holds when pure normalization of source program \(K\) first reaches the checked
semantic question \(q'\).  It is empty when \(K\) returns before another question.
Thus
\(\HeadQ_{\Sigma}\subseteq\mathsf{IProg}(A)\times\mathsf{Question}_{\B}(\Sigma)\) is a partial
functional relation at a fixed deterministic normalization version, restricted to programs and
capture-safe answer extensions belonging to \(\Sigma\).  Write \(\HeadQ\) when \(\Sigma\) is
inferable.
\end{definition}

\begin{definition}[Occurrence-indexed question succession]\label{def:qsucc-occurrence}
For \(\widehat S\in\mathsf{SuppAns}(\mathsf{AskQuestion}(\AskRef))\),
\[
\boxed{
\QSucc(\AskRef,\widehat S,q')
\iff
\HeadQ(\kappa_{\AskRef}(\widehat S),q').
}
\]
The protected successor presentation
\(\QStep(\Sigma,\AskRef,\widehat S)\) is obtained by the same answer substitution and any pure canonical
reconstruction preceding the next open effect.
Equivalently, at fixed versions,
\[
\QSucc\subseteq
\sum_{\AskRef:\mathsf{AskOcc}_A(\Sigma)}
\mathsf{SuppAns}(\mathsf{AskQuestion}(\AskRef))
\times\mathsf{Question}_{\B}(\Sigma).
\]
\end{definition}

\begin{theorem}[Question-succession typing]\label{thm:qsucc-occurrence-typing}
If \(\AskRef\) checks as an \(\mathsf{Ask}\) occurrence, \(\widehat S\) is a supported answer to its
exact question, and \(\QSucc(\AskRef,\widehat S,q')\), then \(q'\) is well typed in the lawful
first-order environment obtained by capture-safe substitution of \(\widehat S\) into
\(\kappa_{\AskRef}\).
\end{theorem}

\begin{proof}
The \textsc{Ask} rule types \(\kappa_{\AskRef}(\widehat S)\) in the explicit continuation
environment.
Pure normalization preserves source-program typing, so its first open question is typed in that
environment.  No additional transition or typing rule is used.
\end{proof}

\begin{law}[Occurrence identity is consequential]\label{law:qsucc-occurrence-identity}
In general,
\[
\mathsf{AskQuestion}(\AskRef_1)=\mathsf{AskQuestion}(\AskRef_2)
\land \widehat S_1=\widehat S_2
\not\Rightarrow
\QStep(\Sigma,\AskRef_1,\widehat S_1)
=
\QStep(\Sigma,\AskRef_2,\widehat S_2).
\]
Therefore no dynamic route, learned question method, trace fingerprint, or replay check may identify
successions only by \((q,\widehat S)\).  It retains the occurrence or an equivalent checked continuation
identity.  A partial supported set is passed whole; succession never selects an arbitrary member.
\end{law}

\begin{law}[Static and dynamic question structure differ]\label{law:q-static-dynamic}
A static relation between reified questions---refinement, shared anchor, converse orientation,
breaker duality, binding-supplied adjunction, or reciprocal pairing---does not create a realized
answer-conditioned route.  Conversely, an occurrence may unlock a successor that is not its static
complement, converse, or refinement.
\end{law}

\subsection{Portable chat execution}

If every open port of a source program is \(\mathsf{Pure}\) or \(\mathsf{Generate}\), and the required context is present in the chat, a backend may execute the source program directly by rendering each current question and using its supported answer to construct the continuation.  If any port is \(\mathsf{Probe}\), \(\mathsf{Check}\), or \(\mathsf{Warrant}\), the source program remains well formed but cannot lawfully discharge that port by generative text alone.

\begin{law}[Question programs are not instruction schemas]
A source inquiry program specifies open typed relations and answer-dependent continuations.  Natural-language instructions may be used by a compiler to establish runtime bindings or formatting, but they do not replace the semantic question relation.  In particular, merely suppressing an alternative in prompt text does not establish the relational exclusion required by a question.
\end{law}

\subsection{Canonical derived combinators}

\begin{definition}[Supported-family lifting]\label{def:supported-family-lifting}
For \(\widehat S:\mathsf{SuppAns}(q)\) and a typed question family
\(a\mapsto q_a=?_{I_a}R_a[\beta_a]\), choose a finite nonempty materialized family
\(F\subseteq|\widehat S|\) with declared coverage \(c_F\).  Put
\[
\boxed{
J_F:=\sum_{a\in F}I_a,
\qquad
X^F_{(a,i)}:=X^a_i,
\qquad
A_F:=\prod_{(a,i)\in J_F}X^a_i
\cong\prod_{a\in F}A(q_a).
}
\]
Let \(R^\Pi_{F,q_\bullet}\) be the ordinary product/conjunction of the child relation instances,
with completion predicate
\[
\boxed{
\Comp_{R^\Pi_{F,q_\bullet}}\!\bigl((b_a)_{a\in F}\bigr)
\iff
\forall a\in F.\,\Comp_{q_a}(b_a).
}
\]
The finite ordinary dependent question is
\[
\boxed{
\mathsf{LiftQ}_F(\widehat S,q_\bullet;c_F)
:=
?_{J_F}R^\Pi_{F,q_\bullet}.
}
\]
Its answer carrier is \(A_F\), and its discharge authority is inherited pointwise:
\[
\boxed{
\operatorname{mode}(\mathsf{LiftQ}_F(\widehat S,q_\bullet;c_F),(a,i))
:=
\operatorname{mode}(q_a,i).
}
\]
The checked family route records \((a,i)\) before discharging that position; it may not replace
heterogeneous child modes by one common mode or use a generated tag to discharge a
\(\mathsf{Probe}\), \(\mathsf{Check}\), or \(\mathsf{Warrant}\) position.
Every completion remains tagged by its predecessor \(a\), the proof-carrying parent
\(\widehat S\), and its own declared-route evidence.  Iterated lifting therefore produces a
dependent family of complete tagged paths.  It neither selects one member nor introduces a source
constructor or runtime opcode.  The lift has whole-parent coverage exactly when \(c_F\) checks
\(F=|\widehat S|\).  Otherwise every member in \(|\widehat S|\setminus F\) remains explicitly
\(\mathsf{Unknown}\), and no consumer may treat \(\mathsf{LiftQ}_F\) as the complete family.
\end{definition}

\begin{definition}[Dependent reciprocal program]\label{def:recip-program}
For \(\alpha\in\{X,Y\}\), source presentation \(W=W_D^\alpha(s)\), protected horizon
\(\mathcal H\), and declared semantic/execution coverage \(c_{sem},c_{exec}\), use the following
ordinary typed question families:
\[
\begin{aligned}
q^{use}_\alpha(W,s)
&:=?u[\mathsf{ApplicableNegationUse}(D,\alpha,W,s,u)],\\
q^{ext}_\alpha(W,s,u)
&:=?(e,w)[\PosNeg(u;s,e,w)],\\
q^{fib}_\alpha(u,e)
&:=?F[\mathsf{SameUseFiber}(u,e,F)],\\
q^{ret}_\alpha(u,e,F)
&:=?(R,m_R)[R\in F\land m_R:\mathsf{Member}(R,F)],\\
q^{rec}_\alpha(W,\mathcal H,u,e,F,c_{sem},c_{exec})
&:=?\rho[\mathsf{RecoveryProfile}(W,\mathcal H,u,e,F,c_{sem},c_{exec},\rho)].
\end{aligned}
\]
Here \(\mathsf{SameUseFiber}(u,e,F)\) admits exactly a first-order fiber record whose membership
relation is \(\RetField_u(e)\).  The relation \(\mathsf{RecoveryProfile}\) returns the complete
three-valued map of \cref{def:recovery-loss-profile} for every declared \(r\in W\), retaining
\(s,e,u,F,\mathcal H,c_{sem},c_{exec}\), occurrence evidence, and local/web/family provenance.
Thus recovery is never a function of \(F\) alone.

The reorientation questions are
\[
\begin{aligned}
q^{seed}_{X\to Y}(u_X,O_X,w_X,F_X,\rho_X)
&:=?(S_Y,\sigma,w_\sigma)
  [\mathsf{AdmittedSeed}(u_X,O_X,w_X,F_X,\rho_X,S_Y,\sigma,w_\sigma)],\\
q^{pres}_Y(S_Y)
&:=?(W_Y,w_W)[\mathsf{AdmittedPresentation}(D,Y,S_Y,W_Y,w_W)],
\end{aligned}
\]
where \(w_\sigma\) and \(w_W\) are exact support/warrant records.  No undeclared
\(\mathsf{Present}\) predicate is used.  Once the complete dependent occurrence \(\omega\) exists,
the downstream question is
\[
\boxed{
q^\Gamma_D(\omega)
:=?\gamma[\mathsf{GammaOutcome}_D(\omega,\gamma)].
}
\]
Its answer carrier is the tagged sum of a compatibility certificate, an incompatibility witness,
and an \(\mathsf{Unknown}\) coverage residual.  The open port of \(q^\Gamma_D\) has mode
\(\mathsf{Check}\); \(\mathsf{Check}\) is not a result constructor.  Likewise \(q^{fib}_\alpha\)
uses \(\mathsf{Pure}\), \(q^{rec}_\alpha\) uses \(\mathsf{Check}\), and
\(q^{pres}_Y\) uses \(\mathsf{Warrant}\).  Every other port keeps the exact mode declared by its
relation use and route.

The derived first-order program \(\mathsf{Recip}(D,W_D^X(S_X))\) is the nested
\(\mathsf{Ask}\)/named-continuation expansion of the dependency graph
\[
\boxed{
\begin{gathered}
q^{use}_X
\to q^{ext}_X
\to q^{fib}_X
\to (q^{ret}_X,q^{rec}_X)
\to q^{seed}_{X\to Y}
\to q^{pres}_Y\\
\to q^{use}_Y
\to q^{ext}_Y
\to q^{fib}_Y
\to (q^{ret}_Y,q^{rec}_Y)
\to q^\Gamma_D.
\end{gathered}
}
\]
Every arrow is implemented by the checked finite
\(\mathsf{LiftQ}_F(\widehat S,q_\bullet;c_F)\), so all proof-carrying predecessor answers and tags
are preserved.  If \(F\neq|\widehat S|\), the uncovered family remains \(\mathsf{Unknown}\) and no
complete downstream role may be claimed.  The second use is independently admitted from the
supported seed presentation; one-way success cannot fill it.  A boundary point contributes only
provenance, and \(q^\Gamma_D\) is reachable only after every dependent role, use, fiber, recovery,
seed, and residual field exists.  This is a transparent macro over ordinary first-order source
syntax and checks, never a runtime opcode, hidden closure, or six-slot oracle.
\end{definition}

\begin{definition}[Hole solving]
Given residual web \(W\) with open carrier \(X\), \(\mathsf{HoleSolve}(W)\) computes
\[
S=\operatorname{Sol}^{X}_{W}.
\]
If \(S/{\equiv_{\mathcal H}}\) is a singleton, it returns the protected filler class.  Otherwise it asks for a relation that separates two surviving protected classes, extends \(W\) by the returned relation, recomputes the meet, and recurs subject to the effectivity horizon.
\end{definition}

\begin{definition}[Representation search]
Given a witnessed pair \(x,y\) with \(\eta(x)=\eta(y)\) but \(x\not\equiv_{\mathcal H}y\), \(\mathsf{RepresentDiff}(x,y)\) is the derived inquiry program that searches, in binding-appropriate order, for a separating context, representation, grain, probe, language extension, or rebind.  A successful candidate is tested for protected consequence relevance and then for unnecessary excess detail.
\end{definition}

\section{Program core}

\subsection{Core grammar}

The executable inquiry core is:
\[
\boxed{
P::=
\mathsf{Return}(a)
\mid
\mathsf{Branch}\{P_i\}_{i\in I}
\mid
\mathsf{Probe}(o,\kappa^{raw}).
}
\]

Interpretation:
\begin{itemize}
\item \(\mathsf{Return}(a)\): pure return with no external actuality event;
\item \(\mathsf{Branch}\{P_i\}\): internal/generative alternatives, none of which is thereby actual;
\item \(\mathsf{Probe}(o,\kappa^{raw})\): actualizable interaction whose lower-level operational
      handler receives the immutable raw return only after its ordinary event has been appended.
      The handler is not a source answer continuation.
\end{itemize}

\subsection{Typing rules}

\[
\frac{\Gamma\vdash a:A}
{\Gamma\vdash\mathsf{Return}(a):\Prog(A)}
\quad(\textsc{Return})
\]

\[
\frac{
I\neq\varnothing
\qquad
\forall i\in I.\ \Gamma\vdash P_i:\Prog(A)
}{
\Gamma\vdash\mathsf{Branch}\{P_i\}_{i\in I}:\Prog(A)
}
\quad(\textsc{Branch})
\]

If
\[
\Gamma\vdash o:\ProbeIR[R]
\]
and
\[
\Gamma,r:\Raw(R)\vdash\kappa^{raw}(r):\mathsf{Fin}(\Prog(A)),
\]
then:
\[
\frac{
\Gamma\vdash o:\ProbeIR[R]
\qquad
\Gamma,r:\Raw(R)\vdash\kappa^{raw}(r):\mathsf{Fin}(\Prog(A))
}{
\Gamma\vdash\mathsf{Probe}(o,\kappa^{raw}):\Prog(A)
}
\quad(\textsc{Probe})
\]

A deterministic continuation is the singleton specialization.

\begin{definition}[Source-safe Ask lowering]\label{def:source-safe-ask-lowering}
For \(\AskRef:\mathsf{AskOcc}_A(\Sigma)\) and
\(q=\mathsf{AskQuestion}(\AskRef)=?_I R[\beta]\), define the evidence carrier for each open port
by its declared mode:
\[
\boxed{
\mathsf{PortEvidence}(\AskRef,i)
:=
\begin{cases}
\mathsf{PureEvidence}(\AskRef,i),
  &\operatorname{mode}(q,i)=\mathsf{Pure},\\
\mathsf{GenerationEvidence}(\AskRef,i),
  &\operatorname{mode}(q,i)=\mathsf{Generate},\\
\mathsf{ProbeEvidence}(\AskRef,i),
  &\operatorname{mode}(q,i)=\mathsf{Probe},\\
\mathsf{CheckEvidence}(\AskRef,i),
  &\operatorname{mode}(q,i)=\mathsf{Check},\\
\mathsf{WarrantEvidence}(\AskRef,i),
  &\operatorname{mode}(q,i)=\mathsf{Warrant}.
\end{cases}
}
\]
Pure, generation, check, and warrant evidence retain their exact typed result, route, resolution
path, versions, and provenance; generation evidence carries no actuality authority.  Probe evidence is the dependent
record
\[
\boxed{
\begin{aligned}
\mathsf{ProbeEvidence}(\AskRef,i)
&:=\sum_{R_i}\sum_{o_i:\ProbeIR[R_i]}\\[-1mm]
&\quad\sum_{E_i:\mathsf{ActualEvent}}\sum_{r_i:\Raw(R_i)}
\sum_{p_i:\Raw(R_i)\leadsto X_i}
\bigl(\mathsf{EventFor}(E_i,\AskRef)
\times\mathsf{PortFor}(E_i,\AskRef,i)\\
&\times\mathsf{OperatorFor}(E_i,o_i)
\times\mathsf{RawFor}(E_i,r_i)
\times\mathsf{ResolutionRouteFor}(i,E_i,p_i)\bigr).
\end{aligned}
}
\]
The \(\mathsf{PortFor}\) checker follows the accepted source lowering from \(i\) to the exact
operator/request; it is not inferred from event order.  \(\mathsf{OperatorFor}\) checks that exact
compiled operator/request, \(\mathsf{RawFor}\) checks that the event's immutable
\(\ReturnRef\) names exactly \(r_i\), and \(\mathsf{ResolutionRouteFor}\) checks the exact stored
path \(p_i\) from that raw-return type to the port carrier.  Byte equality without event identity
is insufficient.

The complete port-indexed bundle and program-wide resolver are
\[
\boxed{
\begin{aligned}
\mathsf{DischargeBundle}(\AskRef)
&:=\prod_{i\in I}\mathsf{PortEvidence}(\AskRef,i),\\
\mathsf{Resolver}(\AskRef)
&:=\mathsf{DischargeBundle}(\AskRef)
  \to\mathsf{ResolutionOutcome}(q).
\end{aligned}
}
\]
The bundle may contain zero, one, or several ordinary actual events.  Shared events require an
explicit checked multi-port lowering; separate events retain separate identities and route
provenance.  The resolver checks cross-port relation constraints and support after all components
are available; it does not manufacture a Cartesian answer when \(R\) forbids one.
Write
\[
\B;\Gamma\vdash
P\;\mathsf{lowers}_{\Sigma}\;(\AskRef,d):A
\]
only under the premises
\[
\Sigma:\mathsf{SourceConfig}_{\B}(A),
\qquad
\AskRef:\mathsf{AskOcc}_A(\Sigma),
\qquad
\B;\Gamma\vdash d:\mathsf{Resolver}(\AskRef),
\]
and when \(P:\Prog(A)\) expands only to the existing
\(\mathsf{Return}/\mathsf{Branch}/\mathsf{Probe}\) core, covers every \(i\in I\) exactly once under
its declared mode, and every probe component first enters the ordinary event spine.  Every complete
execution follows the typed factorization
\[
\mathcal D:\mathsf{DischargeBundle}(\AskRef)
\longrightarrow
d(\mathcal D)=\zeta:\mathsf{ResolutionOutcome}(q)
\]
before any source continuation is reached, with
\[
\begin{array}{rcl}
\zeta=\mathsf{Supported}(\widehat S)
  &\Rightarrow& \Compile_{\mathsf{compiler}(\Sigma)}(\kappa_{\AskRef}(\widehat S)),\\
\zeta=\mathsf{ExactEmpty}(\nu)
  &\Rightarrow& \text{typed residual or justified stop carrying }\nu,\\
\zeta=\mathsf{Undefined}(\nu)
  &\Rightarrow& \text{typed residual or justified stop carrying }\nu,\\
\zeta=\mathsf{Unsupported}(\nu)
  &\Rightarrow& \text{typed residual or justified stop carrying }\nu,\\
\zeta=\mathsf{Unknown}(\nu)
  &\Rightarrow& \text{typed residual or justified stop carrying }\nu.
\end{array}
\]
The lower-level handler may preserve and route raw actuality, but it is ill typed for it to fill
the source answer slot directly.  No non-\(\mathsf{Supported}\) outcome may invoke
\(\kappa_{\AskRef}\).  A single-Probe-port lowering is only the one-probe specialization of this
bundle judgment.  This is a compiler/lowering judgment, not a new runtime constructor.
\end{definition}

\begin{law}[Source Ask compilation is source-safe]\label{law:source-ask-safe}
Every compilation of a source \(\mathsf{Ask}\) establishes
\cref{def:source-safe-ask-lowering}.  A direct predecessor runtime
\(\mathsf{Probe}(o,\kappa^{raw})\) remains a runtime program but is not thereby a lawful lowering of
an occurrence-specific source \(\mathsf{Ask}\).
\end{law}

\section{Sequencing and composition of programs}

Define:
\[
P\bind k
\]
to run \(P\) and then continue every returned value \(a\) with \(k(a)\).

\subsection{Equations}

\[
\boxed{
\mathsf{Return}(a)\bind k
=
k(a).
}
\]

\[
\boxed{
\mathsf{Branch}(S)\bind k
=
\mathsf{Branch}\{P\bind k:P\in S\}.
}
\]

\[
\boxed{
\mathsf{Probe}(o,\kappa^{raw})\bind k
=
\mathsf{Probe}
\left(
o,
r\mapsto
\{P\bind k:P\in\kappa^{raw}(r)\}
\right).
}
\]

On the fragment where the operations are defined, require monad-style identities modulo protected program equivalence:
\[
\mathsf{Return}(a)\bind k=k(a),
\]
\[
P\bind\mathsf{Return}\equiv P,
\]
\[
(P\bind k)\bind h
\equiv
P\bind(\lambda x.\,k(x)\bind h).
\]

\section{Operational semantics}

A runtime configuration is:
\[
\boxed{
\langle P,X_t\rangle
}
\]
where \(X_t\) is the authoritative runtime state defined later.

\subsection{Return}

\[
\langle\mathsf{Return}(a),X_t\rangle
\Downarrow a.
\]

\subsection{Internal branch}

\[
\langle\mathsf{Branch}\{P_i\},X_t\rangle
\longrightarrow
\{\langle P_i,X_t\rangle:i\in I\}.
\]
This transition changes represented possibility; it creates no actuality event.

\subsection{Probe request}

For
\[
P=\mathsf{Probe}(o,\kappa^{raw}),
\]
the runtime:
\begin{enumerate}
\item validates applicability and discharge authority;
\item actualizes \(\den{o}_{\B}\);
\item receives raw return \(r\) and successor coupled configuration \(z'\);
\item appends the raw occurrence to authoritative event history before decoding;
\item supplies \(r\) to the lower-level \(\kappa^{raw}\); for a compiled source Ask, the
      source-safe lowering judgment requires event linkage and a declared resolution outcome
      before the source continuation can be reached.
\end{enumerate}

Schematically:
\[
\frac{
\den{o}_{\B}(z,r,z')
}{
\langle\mathsf{Probe}(o,\kappa^{raw}),X_t[z]\rangle
\longrightarrow
\{
\langle P,X_{t+1}[z']\rangle:
P\in\kappa^{raw}(r)
\}
}.
\]

\begin{law}[Raw-return immutability]
A decoder, semantic interpreter, or later patch may not mutate the authoritative record of the raw return that occurred.
\end{law}

\section{Guarded recurrence}

Potentially unbounded inquiry is permitted only through guarded recurrence.  Every productive corecursive cycle must cross:
\begin{itemize}
\item an actual probe boundary;
\item a finite internal branch whose state changes; or
\item another binding-authorized progress boundary with explicit semantics.
\end{itemize}

Pure invisible recurrence such as
\[
\mathsf{rec}\ f=f
\]
is rejected as nonproductive in the inquiry runtime.

\part{Compiler and LLM Binding}

\section{Compiler architecture}

Compilation is staged:
\[
\boxed{
\begin{gathered}
\text{typed forms/relation schema}\to\text{open relation/question}\to\mathsf{IProg}\\
\to\OpenIR\to\text{positive-negation/dependent reciprocal chart}\to\text{active view}\\
\to\ProbeIR\to\text{backend code}\to\Prog.
\end{gathered}
}
\]

No compiler stage may silently strengthen semantic authority.

\subsection{Stage 1: normalize the open relation}

\[
\Normalize:
\OpenIR\to\OpenIR.
\]

Normalization may alpha-rename ports, normalize joins/hiding, expand derived logical syntax, and canonicalize type references.  It must preserve the completion relation:
\[
\Comp_{\Normalize(q)}\cong\Comp_q.
\]

\subsection{Stage 2: construct the local reciprocal chart}

The compiler constructs or selects:
\[
\boxed{
\mathcal B_q=
(
D_q,
W_q,
\partial_q,
\mathcal U_q^X,
\mathcal U_q^Y,
\mathcal N_q^X,
\mathcal N_q^Y,
\mathsf{coverage}_{sem},
\mathsf{coverage}_{exec},
\sigma_q,
\mathcal F_q,
\mathcal R_q,
\Gamma_q,
T_q,
g_q,
\mathcal H_q
).
}
\]

Here \(W_q\) is the versioned source determination presentation; \(\partial_q\) contains candidate boundary incidence and status; \(\mathcal U_q^\alpha\) and \(\mathcal N_q^\alpha\) retain typed use identities and tagged frontiers; \(\mathcal F_q\) and \(\mathcal R_q\) retain use-specific return fibers and recovery results.  Seed and reciprocal fields may remain open.  For a multi-answer question, this may be a family of overlapping charts.  The compiler must not invent a global binary partition, a reverse negation, an exterior from projection, or a completed sixfold when the semantic answer structure does not provide one.

\subsection{Stage 3: compute the active view}

\[
v_q=\eta_{q,g}(s_t).
\]

The view builder:
\begin{enumerate}
\item includes structure known to affect the live question or protected future;
\item preserves relevance-undetermined structure as reachable reserve;
\item occludes only structure with an applicable exact or approximate licence;
\item records all reopening predicates used by the occlusion.
\end{enumerate}

\subsection{Stage 4: select discharge route}

For every open port \(i\), the compiler respects:
\[
\operatorname{mode}(q,i).
\]

Pure ports lower to deterministic computation.  Generate ports may lower to an LLM or another semantic generator.  Probe ports lower to actualizable interactions.  Check ports lower to admitted checkers.  Warrant ports lower to the standing engine.

A mixed question may compile to a program that discharges different ports through different routes.

\section{LLM possibility relation}

For model parameters \(\Theta\), define the binding-level conditional possibility relation:
\[
\boxed{
G_\Theta:
\mathsf{ContextIR}
\times
\PromptIR
\times
\mathsf{Budget}
\rightsquigarrow
\mathsf{RawCompletion}.
}
\]

The weights \(\Theta\) parameterize the latent possibility field.  The calculus neither duplicates nor claims direct symbolic access to that field.

\begin{law}[Weights/trace separation]
\[
\boxed{
\text{weights parameterize possibility;}
\quad
\text{prompt/response occurrences constitute realized history.}
}
\]
\end{law}

\section{Prompt as executable operator}

For semantic question \(q\), boundary chart \(\mathcal B_q\), and active view \(v_q\), the LLM compiler emits:
\[
\boxed{
u=
\Render_{\Theta,\B}(q,\mathcal B_q,v_q,\mathsf{policy})
:
\PromptIR.
}
\]

The prompt is the executable representation of the probe operator for the LLM backend.

A runtime LLM occurrence is:
\[
\boxed{
r
\in
G_\Theta(v_q,u,b).
}
\]

The semantic question is not identified with the prompt text:
\[
q\neq u.
\]
Prompt wording may change while preserving the same typed operator contract.

\section{Canonical rendering and elaboration}

Define:
\[
\Render_{\B}:
\OpenIR\to\PromptIR
\]
and partial elaboration:
\[
\Elab_{\B}:
\PromptIR\rightharpoonup\mathcal P_f(\OpenIR).
\]

For controlled canonical renderings, target:
\[
\boxed{
\Normalize(
\Elab_{\B}(\Render_{\B}(q))
)
=
\{\Normalize(q)\}.
}
\]
Free prose receives no such guarantee.

\subsection{Rendering obligations}

A canonical renderer must preserve:
\begin{itemize}
\item relation schema identity;
\item bound referents;
\item open port names and types;
\item scope and applicability;
\item boundary orientation when protected;
\item source determination-presentation identity and version;
\item every negation-use identity and its departure-witness provenance;
\item semantic coverage independently from execution/materialization coverage;
\item candidate, open, generated, actual, checked, standing, and unknown status without strengthening;
\item use-specific return-fiber, selected-return, seed, and recovery provenance when present;
\item discharge route;
\item resolution criterion;
\item any comparator required for recurrent probing.
\end{itemize}

\subsection{Derived interrogative rendering and erasure}

\textbf{Status: DERIVED.}  Surface forms such as \emph{what}, \emph{why}, \emph{how},
\emph{why not}, \emph{under what conditions}, and \emph{what if} are controlled renderings of
ordinary typed open relations.  They are not semantic question species.  The same word may render
support, mechanism, contrast, necessity, provenance, applicability, or another binding-native
relation, and different words may normalize to the same question.

An implementation may attach an implementation-only root or route annotation to guide rendering
or bounded search, provided
\[
\boxed{
\Normalize(\mathsf{eraseAnn}(q,\mathsf{ann}))=\Normalize(q).
}
\]
The annotation cannot alter the answer carrier, occurrence-specific continuation, discharge mode,
standing, use-specific reciprocal provenance, or coverage.

Controlled rendering preserves whether a backward relation is an existential converse/preimage, a
binding-supplied universal guarantee or adjoint, a causal/prerequisite relation, or the reverse
section of an exact admitted reciprocal use.  In particular, rendering ``under what conditions''
may not strengthen an existential preimage into a weakest sufficient condition, and surface
opposition may not synthesize a contextual \(\mathsf{NegationUse}\).  Classical formula negation
inside a breaker remains the \(\mathsf{FormulaIR}\) constructor, distinct from contextual typed
negation.

\begin{theorem}[Conservative interrogative lowering]\label{thm:interrogative-lowering}
Let \(K_Q\) be a well-typed source presentation using the derived root and route annotations of
\cref{sec:derived-interrogative-algebra}.  Expanding every transparent macro, compiling each
dependent route to occurrence-indexed \(\mathsf{Ask}(q,\kappa)\), and erasing explanatory
annotations yields an ordinary v1.1 \(\mathsf{IProg}\) with the same semantic questions, whole
supported-answer branches, discharge modes, actuality obligations, use-specific provenance, and
protected continuations.
\end{theorem}

\begin{proof}
Expose is ordinary partial binding; orientation uses a relation or its admitted converse; factor
uses join, exposure, and hiding; polarization opens an ordinary logical breaker, separator, or
already admitted positive-negation relation; variation opens a binding-supplied transform; and
grounding opens ordinary support, check, or warrant relations.  Dependent composition is the
canonical \textsc{Ask} rule.  Route labels add no transition and erase.  Reciprocal cases remain the
derived program of \cref{def:recip-program}.  Thus no semantic constructor or runtime effect is
added.
\end{proof}

\begin{corollary}[No interrogative runtime opcode]\label{cor:no-interrogative-opcode}
No runtime opcode for \textsf{Why}, \textsf{How}, \textsf{WhatIf}, \textsf{Unlock},
\textsf{Root}, \textsf{Orbit}, or \textsf{Saturate} is required.  Any implementation convenience
with such a name must erase to the ordinary source/runtime core while preserving authority and
actuality.
\end{corollary}

\section{Behavioral compiler correctness}

Literal text equality is not the relevant criterion.  For protected continuation family \(\mathcal H\), define:
\[
u_1\equiv_{\mathcal H}^{LLM}u_2
\]
when no protected LLM-facing continuation distinguishes the operational inquiry behavior induced by the two prompts.

For operator \(o\), compiler correctness is:
\[
\boxed{
\operatorname{Behavior}(\Compile_\Theta(o))
\equiv_{\mathcal H}
\den{o}
}
\]
or, where exact equality is impossible, an explicitly licensed approximation:
\[
\operatorname{Behavior}(\Compile_\Theta(o))
\approx_{\mathcal H,\epsilon}
\den{o}.
\]

The approximation relation must name its discriminator set, error relation, and reopening condition.

\section{Question-conditioned LLM field}

A practical LLM context may be constructed from:
\[
\boxed{
\mathcal F_t(q)
=
 \mathsf{Active}_t(q)
+
 \mathsf{Prior}_t(q)
+
Cross_t(q)
+
Retrieve_t(q)
+
Methods_t(q)
+
Reopen_t(q).
}
\]

These symbols denote derived projections:
\begin{itemize}
\item \(\mathsf{Active}_t(q)\): currently active relation-bearing structure;
\item \(\mathsf{Prior}_t(q)\): prior comparable occurrences of the same probe contract;
\item \(Cross_t(q)\): other probe traces relevant to covariation or dissociation;
\item \(Retrieve_t(q)\): distant retained/documentary structure adjacent under the open relation;
\item \(Methods_t(q)\): native methods whose typed input/output relation may discharge a port;
\item \(Reopen_t(q)\): unlock conditions that may reactivate currently occluded structure.
\end{itemize}

This is not a separate memory ontology.  It is a question-conditioned projection of retained structure.

\part{History, State, and Recurrent Probing}

\section{Three orders that must not collapse}

The runtime distinguishes:
\[
\boxed{
<_{\mathrm{ledger}},
\qquad
\lambda_{\mathrm{real}},
\qquad
T_D.
}
\]

\begin{itemize}
\item \(<_{\mathrm{ledger}}\): append/storage order of authoritative event records;
\item \(\lambda_{\mathrm{real}}\): binding-defined realized/domain succession;
\item \(T_D\): inquiry traversal along a distinction boundary.
\end{itemize}

No equation among these is constitutional.

\section{Single authoritative ancestry and typed projections}

\begin{definition}[Authoritative ancestry]\label{def:authoritative-ancestry}
\[
\boxed{
H^{auth}
=
\mathsf{AcceptedRoots}
\cup H^{evt}
\cup H^{acc}
}
\]
is the one authoritative provenance structure: accepted content-addressed source/program/artifact
roots, ordinary actual events, and independently accepted revisions linked by immutable typed
references.  The event and revision sequences below are non-collapsing typed projections of
\(H^{auth}\), not competing histories.  Derived traces, routes, presents, memories, methods, and
reciprocal views can point into this ancestry but cannot append semantic actuality or acceptance on
their own authority.
\end{definition}

\begin{definition}[Event history]
\[
\boxed{
H_t^{evt}=E_0\cdot E_1\cdots E_{t-1}
}
\]
is the append-only authoritative history of realized probe occurrences and other explicitly admitted external returns.
\end{definition}

\begin{definition}[Accepted revision history]
\[
\boxed{
H_t^{acc}=P_0\cdot P_1\cdots P_{m-1}
}
\]
is the authoritative history of accepted semantic/traversal/compiler/binding/protection patches.
\end{definition}

The event and accepted-revision projections have independent edge semantics inside one ancestry.
A semantic interpretation may be revised without rewriting the external event from which it was
derived.

\section{Event record}

A canonical actual event record is:
\[
\boxed{
E_t=
\left\langle
\StateRef_t,
\mathsf{sourceAsk}_t,
q_t,
\BoundaryRef_t,
\OpRef_t,
\ReturnRef_t,
\StateRef_{t+1},
g_t,
\mathsf{route}_t,
\mathsf{prov}_t
\right\rangle.
}
\]
The source-occurrence field has the dependent option type
\[
\boxed{
\mathsf{sourceAsk}_t:
\mathbf 1+
\sum_{A:\mathsf{Ty}_{\B}}\sum_{\Sigma:\mathsf{SourceConfig}_{\B}(A)}
\mathsf{AskOcc}_A(\Sigma).
}
\]
The \(\mathbf 1\) summand is \(\mathsf{none}\) for a direct or legacy runtime probe; every event
compiled from a source \(\mathsf{Ask}\) uses \(\mathsf{some}(\AskRef)\).

The record preserves:
\begin{itemize}
\item the state/configuration reference before the occurrence;
\item the exact source \(\mathsf{Ask}\) occurrence when the event was compiled from one;
\item the semantic question/probe-schema reference;
\item the active boundary chart reference;
\item the exact compiled operator/prompt reference;
\item the raw return reference;
\item the successor state/configuration reference;
\item the representation grain;
\item the discharge/actualization route;
\item provenance, backend version, binding version, and required replay metadata.
\end{itemize}

\begin{definition}[Source event linkage]\label{def:event-for-ask}
For \(A:\mathsf{Ty}_{\B}\), \(\Sigma:\mathsf{SourceConfig}_{\B}(A)\), and
\(\AskRef:\mathsf{AskOcc}_A(\Sigma)\), source event linkage is the checked proposition
\[
\boxed{
\mathsf{EventFor}(E,\AskRef)
}
\]
holds exactly when
\[
\mathsf{sourceAsk}(E)=\mathsf{some}(\AskRef),
\qquad
\mathsf{question}(E)=\mathsf{AskQuestion}(\AskRef),
\]
and the event's operator/request, binding and compiler versions, ledger membership, and
provenance are the checked lowering and dispatch context of that exact occurrence.  Equality of
question, operator, raw return, or endpoint without this link is insufficient.
\end{definition}

\begin{law}[Reciprocal actuality is ordinary actuality]\label{law:reciprocal-actuality}
A generated negation use, exterior candidate, return candidate, seed, recovery claim, or sixfold view is not thereby actual.  Every external/model/tool result enters history through the ordinary event spine before decoding or interpretation.  Partial, failed, or one-sided reciprocal attempts remain reconstructible as ordinary events and residuals.  A completed \(\Xi_D(\omega)\) is a derived view over those events, relation uses, fibers, and support; it is never a substitute authoritative history.
\end{law}

\section{Derived paired actuality and question-route projection}

\textbf{Status: DERIVED.}  The paired view below consolidates the question-side and return-side
projections of an inquiry occurrence without adding another event species or history.

\begin{definition}[Resolved inquiry occurrence]\label{def:resolved-inquiry-occurrence}
For result type \(A\), define the normalized source-control result
\[
\boxed{
\mathsf{Next}_A
:=
\left(\sum_{\Sigma':\mathsf{SourceConfig}_{\B}(A)}\mathsf{AskOcc}_A(\Sigma')\right)
+\den{A}_{\B}.
}
\]
The checked relation \(\mathsf{HeadResult}_A(\Sigma,K,n)\) uses the environment, versions, and
provenance inherited from \(\Sigma\) to normalize \(K:\mathsf{IProg}(A)\) capture-safely.  Its
result \(n:\mathsf{Next}_A\) is either a newly checked \(\mathsf{AskOcc}_A(\Sigma')\) or the value
of \(\mathsf{Return}_I(a)\), in the corresponding summand.

A resolved occurrence is the dependent first-order record
\[
\boxed{
\begin{aligned}
\mathsf{ResolvedOccurrence}_A
:=\{\langle\Sigma,\AskRef,\mathcal D,d,\widehat S,K,n\rangle\mid{}&
 \Sigma:\mathsf{SourceConfig}_{\B}(A),\quad
 \AskRef:\mathsf{AskOcc}_A(\Sigma),\\
&\mathcal D:\mathsf{DischargeBundle}(\AskRef),\quad
 d:\mathsf{Resolver}(\AskRef),\\
&d(\mathcal D)=\mathsf{Supported}(\widehat S),\\
&\widehat S:\mathsf{SuppAns}(\mathsf{AskQuestion}(\AskRef)),\\
&K=\kappa_{\AskRef}(\widehat S),\quad
 \mathsf{HeadResult}_A(\Sigma,K,n)\}.
\end{aligned}
}
\]
Write the field projections as \(\mathsf{ask}(\omega)\), \(\mathsf{bundle}(\omega)\),
\(\mathsf{answer}(\omega)\), \(\mathsf{continuation}(\omega)\), and
\(\mathsf{next}(\omega)\).  The probe components of \(\mathcal D\) retain every exact
\(E/\ReturnRef/\ResPath\) chain; non-probe components retain their own declared routes.  No
\(\mathsf{ExactEmpty}\), \(\mathsf{Undefined}\), \(\mathsf{Unsupported}\), or
\(\mathsf{Unknown}\) outcome inhabits this carrier.

Thus its provenance-retaining projection is
\[
\AskRef_t
\longrightarrow\mathcal D_t
\longrightarrow\mathsf{Supported}(\widehat S_t)
\longrightarrow\kappa_{\AskRef_t}(\widehat S_t)
\longrightarrow n_t:\mathsf{Next}_A.
\]
Every arrow is checked against the same occurrence and binding/version/provenance context.  The
abbreviated tuple \((z,q,r,z')\), when useful for a one-probe specialization, is only a projection
of this record and cannot replace its complete identity.
\end{definition}

\begin{definition}[Paired actuality view]\label{def:paired-actuality-view}
For a resolved occurrence \(\omega_t\), let \(\mathsf{ProbeComponents}(\omega_t)\) be the
port-and-event-indexed family of Probe components in its discharge bundle and define
\[
\boxed{
\mathsf{PairedActuality}(\omega_t)
=
\bigl(
\mathsf{QuestionTrace}(\omega_t,i,E),
\mathsf{ReturnTrace}(\omega_t,i,E)
\bigr)_{(i,E)\in\mathsf{ProbeComponents}(\omega_t)}.
}
\]
For each \((i,E)\), the question projection retains the source occurrence, port, question, pre-state,
boundary, operator, route, binding, and event.  The return projection retains the same event, raw
return, port resolution path, full program-supported answer, continuation, successor/return,
post-state, route, and binding.  They are coupled by exact event identity; neither is independently
authoritative.  Multiple ports may share one event only through an explicit checked multi-port
lowering; their port-indexed evidence and paired traces remain distinct.  A no-Probe occurrence has
an empty paired-actuality family rather than a fabricated event.
\end{definition}

\begin{law}[One authoritative occurrence spine]\label{law:one-occurrence-spine}
\[
\boxed{
\text{question trace}\neq\text{return trace},
\qquad
\text{paired view}\neq\text{second history}.
}
\]
Endpoint equality does not imply equal event, raw-return, resolution-path, supported-answer, or
continuation provenance.  Conversely, replaying one historical event reconstructs a view of that
same occurrence and does not append a new event.
\end{law}

\begin{law}[Resume, replay, and order separation]\label{law:resume-replay-order}
Resume supplies a supported answer to an already checked continuation.  Replay reconstructs the
historical chain in \cref{def:resolved-inquiry-occurrence} from accepted roots and ordinary event
history, without redispatch.  Therefore
\[
\boxed{
\mathsf{Resume}\neq\mathsf{Replay},
\qquad
<_{\mathrm{ledger}}\neq\lambda_{\mathrm{real}}\neq\QSucc.
}
\]
Ledger membership or order never manufactures a causal edge or a question-successor occurrence.
\end{law}

\begin{definition}[Distinct regenerative cue fibers]\label{def:paired-cue-fibers}
Because raw returns, supported answers, and question occurrences occupy different carriers, their
reconstruction fibers remain distinct.  Write \(\mathsf{NextAsk}(\omega,\AskRef')\) when
\(\mathsf{next}(\omega)=\mathsf{inl}(\Sigma',\AskRef')\) for some checked successor configuration
\(\Sigma'\).  For checked neighboring occurrence information, define:
\[
\CueRaw(\AskRef,\AskRef')
=
\{\mathsf{bundle}(\omega):\exists A.\,
\omega\in\mathsf{ResolvedOccurrence}_A
\land\mathsf{ask}(\omega)=\AskRef
\land\mathsf{NextAsk}(\omega,\AskRef')\},
\]
\[
\CueSupp(\AskRef,\AskRef')
=
\{\mathsf{answer}(\omega):\exists A.\,
\omega\in\mathsf{ResolvedOccurrence}_A
\land\mathsf{ask}(\omega)=\AskRef
\land\mathsf{NextAsk}(\omega,\AskRef')\},
\]
and
\[
\mathsf{CueAsk}(\widehat S_{t-1},\widehat S_t)
=
\{\AskRef:\exists\omega.\,
\mathsf{ContinuationHead}(\widehat S_{t-1},\AskRef)
\land
\mathsf{OccurrenceAnswer}(\AskRef,\widehat S_t,\omega)\}.
\]
A missing position is regenerated exactly only when the appropriate fiber collapses to one
protected class under declared coverage.  Collapse of \(\CueSupp\) does not reconstruct a unique
raw event or resolution path; collapse of a raw-return fiber does not identify a unique
continuation occurrence.
\end{definition}

\begin{law}[Paired regeneration and reopening]\label{law:paired-regeneration}
Question occurrences constrain which supported returns and actual provenance could lawfully connect
them; supported returns constrain which occurrence-specific continuations could connect neighboring
positions.  A recurrent segment may be folded only when the relevant fibers, authority,
provenance, and protected consequences regenerate.  A new separator, probe, binding, or protected
continuation that splits a folded fiber reopens the fold through ordinary ancestry.
\end{law}

\section{State and working presentation}

A state snapshot may be physically large.  The calculus imposes no requirement that it be globally compressed.

Define the authoritative accepted presentation:
\[
\boxed{
A_t=\operatorname{VersionedFold}(H_t^{acc}).
}
\]

Define the current working presentation as:
\[
\boxed{
\Pi_t=\operatorname{Presentation}(A_t,H_t^{evt},\epsilon_t).
}
\]

Any question-conditioned active view is derived from retained state/history:
\[
v_t(q)=\eta_{q,g}(\Pi_t,H_t^{evt}).
\]

No hidden mutable ``memory state'' may become authoritative when the same information is supposed to be reconstructible from event ancestry and accepted patches.

\section{Operator occurrence index}

The compact operator layer is:
\[
\boxed{
\mathsf{OpOcc}_t
=
(
\StateRef_t,
\OpRef_t,
\ReturnRef_t,
\StateRef_{t+1},
\BoundaryRef_t
).
}
\]

For operator \(o\),
\[
\boxed{
Occurrences(o)
=
\{
\EventRef_t:
\OpRef(E_t)=o
\}.
}
\]

The operator index is a lookup over history; it does not duplicate the state snapshots to which it points.

\section{Recurrent probes}

\begin{definition}[Probe contract]
A recurrent probe contract is:
\[
\boxed{
\mathcal C_o=
(
\mathsf{RelRole},
\B,
g,
\mathsf{Applicability},
\mathsf{Comparator},
\mathcal H,
\mathsf{DecoderVersion}
).
}
\]
\end{definition}

Surface wording is neither necessary nor sufficient for probe identity.

\begin{definition}[Occurrence comparability]
Two occurrences \(E_i,E_j\) of an operator are comparable when:
\[
\boxed{
Comparable(E_i,E_j)
\iff
SameContract(E_i,E_j)
\lor
Bridge(\mathcal C_i,\mathcal C_j).
}
\]
A bridge is itself a standing typed relation establishing the protected correspondence between changed contracts.
\end{definition}

\begin{law}[Probe identity]
\[
\boxed{
\text{same wording}
\neq
\text{same probe}
}
\]
and:
\[
\boxed{
\text{same operator identifier}
\not\Rightarrow
\text{comparable occurrence}
}
\]
when binding, grain, comparator, applicability, or protected meaning changed without a bridge.
\end{law}

\section{Probe traces}

For recurrent probe contract \(o\), define the ordered comparable trace:
\[
\boxed{
H_o^t=
\langle E_{t_1},\ldots,E_{t_n}\rangle
}
\]
such that each pair admitted for comparison satisfies the probe contract or an explicit bridge.

Probe traces are ordinary relational operands.  Questions over them can ask for:
\begin{itemize}
\item persistence;
\item change;
\item change point;
\item recurrence;
\item covariation with another trace;
\item dissociation;
\item path dependence;
\item boundary movement;
\item comparator failure.
\end{itemize}

No perceived pattern becomes warranted solely because it is salient in the trace.

\section{Fresh probe before historical comparison}

When the environment/model interface permits a current return to be obtained without injecting prior semantic interpretations, prefer:
\[
\boxed{
\text{fresh probe}
\longrightarrow
\text{trace comparison}.
}
\]

Formally, obtain:
\[
r_t^0=\operatorname{ProbeFresh}(o,z_t,v_t^0),
\]
then:
\[
p_t=
\operatorname{Compare}(H_o^{t-1},E_t^0).
\]

This is a method-level control, not a constitutional primitive.  Its purpose is to prevent continuity in model wording from being mistaken for persistence in the independently probed environment.

\section{Question-order diagnostics}

Let \(U_q\) denote the state/update relation induced by executing and resolving \(q\).  For questions believed to be independent, test:
\[
U_r\circ U_q
\quad\text{against}\quad
U_q\circ U_r.
\]

If:
\[
\boxed{
U_r\circ U_q
\not\equiv_{\mathcal H}
U_q\circ U_r,
}
\]
inquire whether the difference reflects:
\begin{enumerate}
\item a genuine relational interaction/path dependence; or
\item an irrelevant order effect of the LLM/compiler/runtime.
\end{enumerate}

Question-pattern learning must not silently quotient away protected order dependence.


\part{Protected Behavior, Folding, Learning, and Recovery}

\section{Protected behavioral equivalence}

Let \(\mathcal H\) be the protected family of future contexts/continuations applicable to parallel terms.

\begin{definition}[Protected equivalence]
For \(f,g:A\to B\) or parallel relational terms,
\[
\boxed{
f\equiv_{\mathcal H}g
\iff
\forall K\in\mathcal H,\,
\Con_K(f)=\Con_K(g).
}
\]
\end{definition}

\begin{definition}[Separator family]
\[
\boxed{
\Sep_{\mathcal H}(f,g)
=
\{
K\in\mathcal H:
\Con_K(f)\neq\Con_K(g)
\}.
}
\]
Hence:
\[
f\not\equiv_{\mathcal H}g
\iff
\Sep_{\mathcal H}(f,g)\neq\varnothing.
\]
\end{definition}

\begin{law}[Horizon monotonicity]
If
\[
\mathcal H\subseteq\mathcal H',
\]
then:
\[
\boxed{
\equiv_{\mathcal H'}
\subseteq
\equiv_{\mathcal H}.
}
\]
Adding protected future use may split old classes but cannot merge classes already distinguished by the retained horizon.
\end{law}

\begin{definition}[Working nondistinction]
For finite tested discriminator set \(\mathcal D\subseteq\mathcal H\),
\[
x\approx_{\mathcal H,\mathcal D}y
\]
means no discriminator in \(\mathcal D\) has separated \(x,y\).  It is weaker than exact protected equivalence unless the binding licenses completeness of \(\mathcal D\).
\end{definition}

\section{Hom-wise operator quotient}

Because direct equivalence compares parallel terms, quotient typed operator spaces hom-wise:
\[
\boxed{
\mathsf Q_{\mathcal H}(A,B)
=
\mathsf{Tm}_{\B}(A,B)/
{\equiv_{\mathcal H}}.
}
\]

Let:
\[
\eta_{\mathcal H}^{A,B}:
\mathsf{Tm}_{\B}(A,B)
\to
\mathsf Q_{\mathcal H}(A,B)
\]
be the quotient map.

\begin{definition}[Congruence]
Protected equivalence is compositional over an admitted regime when:
\[
f\equiv_{\mathcal H}f',
\qquad
g\equiv_{\mathcal H}g'
\]
imply:
\[
\boxed{
g\circ f
\equiv_{\mathcal H}
g'\circ f'.
}
\]
\end{definition}

When congruence holds, composition descends:
\[
\boxed{
[g]_{\mathcal H}\circ[f]_{\mathcal H}
=
[g\circ f]_{\mathcal H}.
}
\]

\section{Operator descent}

A present-output quotient is not automatically executable recurrent state.

Let:
\[
q:X\to S
\]
be a representation quotient and let:
\[
a:X\rightsquigarrow X'
\]
be a protected continuation.

\begin{definition}[Continuation descent]
Continuation \(a\) descends through \(q\) when there exist:
\[
q':X'\to S'
\]
and:
\[
\bar a:S\rightsquigarrow S'
\]
such that the square commutes relationally:
\[
\boxed{
q'\circ a
=
\bar a\circ q.
}
\]
\end{definition}

For every protected continuation in scope, executable retained state requires such a descent.

\begin{law}[Current-output sufficiency is not recursive-state sufficiency]
\[
\boxed{
\text{present consequence sufficiency}
\not\Rightarrow
\text{continuation sufficiency}.
}
\]
\end{law}


\section{Regenerative sufficiency, understanding, and learning}

\begin{definition}[Regenerative sufficiency]
Let \(m:M\) be a retained representation of a protected form or distinction \(z\).  Write
\[
\boxed{
Regen_{\mathcal H}(m,z)
}
\]
when there exist typed reconstruction relations that recover every protected component required by the current horizon from \(m\), and the reconstructions are protected-equivalent to the corresponding components of \(z\).  For a reciprocal occurrence, this includes its candidate side classes and, when protected, its determination presentations, use-tagged dependent sixfold view, return fibers, recovery judgments, coverage, seed, residuals, and provenance:
\[
\rho_X(m)\equiv_{\mathcal H}x_z,
\qquad
\rho_Y(m)\equiv_{\mathcal H}y_z,
\qquad
\rho_{\Xi}(m)\equiv_{\mathcal H}\Xi_D(\omega),
\qquad
\rho_{use/fib}(m)\equiv_{\mathcal H}(u_X,u_Y,F_X,F_Y).
\]
\end{definition}

\begin{definition}[Inquiry-regenerative sufficiency]
\[
\boxed{
Regen^{inq}_{\mathcal H}(m,z)
}
\]
holds when \(Regen_{\mathcal H}(m,z)\) holds and \(m\) also retains or can regenerate the protected discriminators, residuals, support dependencies, or reopening routes needed to test and revise \(z\).
\end{definition}

\begin{definition}[Regenerative economy frontier]\label{def:regenerative-economy}
Let \(\preceq\) be a binding-supplied resource preorder and let
\(\mathsf{Licensed}_{\mathcal H}(m,z)\) include the required scope, applicability,
authority, provenance, continuation, recovery, and unlock contracts.  Define
\[
\boxed{
\mathsf{Economy}_{\mathcal H,\preceq}(z)
=
\operatorname{Min}_{\preceq}
\{
m:
Regen^{inq}_{\mathcal H}(m,z)
\land
\mathsf{Licensed}_{\mathcal H}(m,z)
\}.
}
\]
This is a minimal or nondominated frontier, not a promised unique global
\(\arg\min\).  It may be empty when no admitted representation satisfies the protected
contract.  Exact consequence factorization such as
\(\ker m\subseteq\ker\chi\) discharges only the named consequence; it is not equivalent
to inquiry-regenerative sufficiency unless the discriminators, continuations,
provenance, residuals, and reopening obligations also factor through \(m\).
\end{definition}

\begin{law}[Differentiate only enough to regenerate]\label{law:regenerative-economy}
Inquiry may add distinctions until protectedly different live completions are
separable, then subtract distinctions whose removal preserves the declared
inquiry-regenerative contract.  A further subtraction is blocked by a witnessed loss
of protected regeneration, discrimination, continuation behavior, warrant provenance,
or reopening.  This law governs licensed active representations and folds; it does not
authorize deletion or rewriting of authoritative history.
\end{law}

\begin{definition}[Understanding]
Understanding is a derived role, not a new semantic species:
\[
\boxed{
Understand_{\mathcal H}(m,z)
\iff
Regen^{inq}_{\mathcal H}(m,z).
}
\]
This definition does not claim that biological or computational memory is unnecessary.  It distinguishes stored historical content from the regenerative competence by which a protected relational structure can be reconstructed and revised.
\end{definition}

\begin{definition}[Ablative regeneration test]
For retained relational presentation \(M\) and form \(x\), define \(\operatorname{Ablate}_x(M)\) by replacing protected occurrences of \(x\) with open ports while retaining their governing relations.  Define
\[
\boxed{
\operatorname{Regen}_x(M)
=
\operatorname{Sol}_{\operatorname{Ablate}_x(M)}.
}
\]
A direct protected regeneration witness is
\[
\operatorname{Regen}_x(M)/{\equiv_{\mathcal H}}
=
\{[x]_{\mathcal H}\}.
\]
This is an executable test of whether the retained web determines \(x\) at the protected grain.
\end{definition}

\begin{law}[Learning as regenerative gain]
\[
\boxed{\textbf{LEARNING IS A STANDING CHANGE IN REGENERATIVE CAPACITY.}}
\]
A learned change may add or refine a distinction, improve protected reconstruction, invent a nonredundant representation/probe, compile a recurrent question/operator path, or reduce the active cost of regenerating protected structure.  Historical events provide evidence for the change but are not identical to the learned competence.
\end{law}

\begin{definition}[Learning gain]
A candidate learned form \(m\) has a regenerative/traversal gain when, under a declared applicability region and resource preorder, it strictly improves at least one of:
\[
\mathsf{Discriminate},\quad
\mathsf{Regenerate},\quad
\mathsf{Traverse},\quad
\mathsf{Reacquire},\quad
\mathsf{Cost},\quad
\mathsf{Robustness}
\]
while preserving protected behavior and warrant boundaries.  Promotion still requires independent standing.
\end{definition}

\section{Method and operator-path folding}

Suppose a recurrent typed path:
\[
P=o_n\circ\cdots\circ o_1
\]
is repeatedly instantiated.

\begin{definition}[Method promotion]
A method \(m:A\to B\) may be promoted with:
\[
m\doteq_{\mathcal H}P
\]
only when:
\begin{enumerate}
\item \(m\) and \(P\) are parallel or explicitly bridged;
\item \(m\equiv_{\mathcal H}P\) over the declared applicability region;
\item required continuations descend;
\item the promotion has a typed operational gain or necessary implementation purpose;
\item defining path/provenance remains recoverable as required;
\item an unlock/reopening contract is stored;
\item method utility does not confer semantic warrant on future outputs.
\end{enumerate}
\end{definition}

The promotion record contains:
\[
\boxed{
\mathsf{MethodFold}
=
(
m,P,\mathcal H,\mathsf{Applicability},
\mathsf{EvidenceRefs},\mathsf{Gain},
\mathsf{Recovery},\mathsf{Unlock}
).
}
\]

\section{Two forms of traversal learning}

\subsection{Path folding}

A recurring question/probe sequence:
\[
o_1;o_2;\cdots;o_n
\]
is compressed to a reusable operator/method under the fold law above.

\subsection{Probe-basis extension}

A new probe \(p'\) may be admitted when the old probe repertoire cannot expose a protected distinction that \(p'\) can expose.

A canonical nonredundancy witness is:
\[
\boxed{
\exists x,y:
\left[
\forall p\in\mathcal P_{\mathrm{old}},\,
p(x)=p(y)
\right]
\land
p'(x)\neq p'(y)
\land
\Con_{\mathcal H}(x)\neq\Con_{\mathcal H}(y).
}
\]

A generated probe remains inert until subsequent independent use demonstrates the claimed nonredundant protected discrimination, retrieval, reasoning, or control benefit.

Thus:
\[
\boxed{
\text{learning what relation stands}
\neq
\text{learning how to expose that relation}.
}
\]

\section{Question-pattern learning}

A realized question-pattern occurrence is represented first by the complete occurrence-indexed
chain:
\[
\boxed{
\mathcal T_t^{occ}
:=
\left\{
\omega_t:\mathsf{ResolvedOccurrence}_A
\right\}.
}
\]
Let
\[
\mathsf{ProbeView}:=
\mathsf{Fin}(\mathsf{PortName}\times\OpRef\times\ReturnRef).
\]
Its compact operator and state/boundary views are the lossy projections
\[
\pi_{op}(\mathcal T_t^{occ})
\subseteq
\OpenIR
\times\mathsf{ProbeView}
\times\OpenIR
\]
and
\[
\pi_{state}(\mathcal T_t^{occ})
\subseteq
\StateRef
\times
\OpenIR
\times
\BoundaryRef
\times
\mathsf{ProbeView}
\times
\StateRef
\times
\OpenIR.
\]
Neither projection identifies the occurrence.  In particular, equal open queries, endpoints,
operators, boundaries, or content-equal raw returns do not erase the checked \(\AskRef_t\), any
event/path member of its discharge bundle, \(\widehat S_t\), the checked continuation, or its
normalized next-control result.  Operator-pattern aggregation may
index either projection only after a protected-equivalence and applicability check licenses that
generalization.

A learned policy/method describes:
\[
(\mathcal T_t^{occ},s_t)\rightsquigarrow q_{t+1}
\]
only under an explicit applicability relation.  It may expose a compact projection for lookup, but
must recover the occurrence chain before selecting an occurrence-specific continuation.  It is not
a universal scheduler or a semantic truth relation.

\section{Memory and recovery}

Memory is not introduced as an independent semantic state kind.  Retained realized history and accepted operator/semantic folds are authoritative.  Future use may nevertheless require several distinct recovery relations.

For retained material \(m\), context \(c\), and competence \(k\), distinguish:
\[
Use_{\mathcal H}(m,c,k),
\]
\[
Reconstruct_{\mathcal H}(m,c,k),
\]
\[
Evaluate_{\mathcal H}(m,c,k),
\]
\[
Reacquire_{\mathcal H}(m,c,k).
\]

No universal implication among these relations is assumed.

\subsection{Use}

The competence can presently be exercised using retained material.

\subsection{Reconstruction}

The protected competence/representation can be regenerated from retained material without substantial new external evidence.

\subsection{Evaluation without reconstruction}

A protected consequence can be evaluated even when the full historical representation is not reconstructed.

\subsection{Reacquisition}

Retained structure lowers the future cost/frontier for regaining competence, even if direct use and reconstruction fail.

Let:
\[
\mathfrak C_{\mathcal H}(m,c,k)
\]
be a binding-relative nondominated reacquisition frontier and:
\[
\mathfrak C^0_{\mathcal H}(c,k)
\]
a pinned baseline without \(m\).  A reacquisition advantage exists when:
\[
\boxed{
\mathfrak C_{\mathcal H}(m,c,k)
\prec
\mathfrak C^0_{\mathcal H}(c,k).
}
\]

\begin{law}[Keep experiences; compile routes]
The default architecture preserves realized history as fully as the storage contract permits and learns smaller routing/operator structure over it.  State compression is an independent licence, not a precondition of learning.
\end{law}

\section{Historical fact versus generated reconstruction}

A reconstruction candidate must distinguish:
\[
\boxed{
\mathsf{HistoricalFactRefs}
\quad\text{from}\quad
\mathsf{GeneratedReconstruction}.
}
\]

Generated continuity may fill a working context, but it does not rewrite historical event ancestry.

When reconstruction leaves several protected-distinct candidates:
\[
\widehat x_1\not\equiv_{\mathcal H}\widehat x_2,
\]
the ambiguity becomes an inquiry obligation rather than being silently collapsed.

\part{Claims, Support, Warrant, and Standing}

\section{Claim lifecycle}

A semantic generator may produce:
\[
\boxed{
\mathsf{Claim}[\phi]
}
\]
without establishing \(\phi\).

A reifier may partially map:
\[
\mathsf{Reify}:
\mathsf{Claim}[\tau]
\rightharpoonup
\mathsf{Candidate}[\tau].
\]

Failure to reify is lawful and may open decomposition, clarification, re-expression, binding, or unknown-status inquiry.

A candidate relation becomes standing only through the warrant system.

\begin{law}[Claim/fact separation]
\[
\boxed{
\mathsf{Claim}[\phi]
\not\Rightarrow
\Stand(\phi).
}
\]
Question syntax, answer type, or model confidence cannot promote a claim.
\end{law}

\section{Support environments}

\begin{definition}[Support environment]
For candidate or retained relation \(\lambda\), a support environment \(E\) is a finite or finitely represented set of typed premises, actual returns, checker results, assumptions, and applicable standing relations claimed to support \(\lambda\).
\end{definition}

\begin{definition}[Minimal support family]
\[
\boxed{
\Supp(\lambda)
=
\operatorname{Min}_{\subseteq}
\{
E:
E\text{ is a candidate support environment for }\lambda
\}.
}
\]
Multiple incomparable support environments may exist.
\end{definition}

\section{Open dependency boundary}

\begin{definition}[Open dependency]
For support environment \(E\),
\[
\boxed{
\partial^-_E\lambda
=
\{
d:
d\text{ is required for }\lambda
\text{ but is neither supplied by }E
\text{ nor independently discharged}
\}.
}
\]
\end{definition}

Every open dependency is itself a lawful question target:
\[
\boxed{
?d[
d\in\partial^-_E\lambda
].
}
\]

\section{Closed support}

Relative to standing set \(X\), define:
\[
\Closed_X(E,\lambda)
\]
when:
\begin{enumerate}
\item all premises of \(E\) requiring standing belong to \(X\);
\item applicability/scope conditions hold;
\item \(\partial^-_E\lambda=\varnothing\);
\item admitted independent checks required by the route succeed;
\item no explicit inconsistency policy invalidates \(E\).
\end{enumerate}

\section{Least-fixed-point standing}

Let \(\mathsf{Ingress}_t\) contain grounded standing facts available independently of inference cycles: preserved actual-return facts, trusted configuration facts, accepted predecessor relations, checker axioms, or other binding-authorized roots.

Define the monotone operator:
\[
\boxed{
T_t(X)
=
\mathsf{Ingress}_t
\cup
\{
\lambda:
\exists E\in\Supp(\lambda),\
\Closed_X(E,\lambda)
\}.
}
\]

Then:
\[
\boxed{
\Stand_t
=
\mu T_t.
}
\]

\begin{theorem}[No rootless positive support cycle]
A finite set of relations whose only support routes depend positively on one another and which contains no grounded ingress element does not enter \(\mu T_t\).
\end{theorem}

\begin{proof}
Starting from the least element \(\varnothing\), no member of the cycle can be added until its dependencies are already present.  By induction on fixed-point iteration, none appears without an external ingress/discharge route.
\end{proof}

\section{Applicability is not evidential support}

An applicability condition \(g\) says when a relation may be used.  It is not a substitute for support.

A retained relation therefore keeps at least:
\[
\boxed{
(
\lambda,
\mathsf{scope},
\mathsf{applicability},
\Supp(\lambda),
\partial^-,
\mathsf{warrantClass},
\mathsf{certificateRefs}
).
}
\]

A change in applicability may deactivate a relation without erasing its historical warrant record.

\part{Compression, Residue, and Reopening}

\section{Exact representation quotient}

Let:
\[
c:X\to Y
\]
be a proposed exact representation quotient.

\subsection{Consequence sufficiency}

Require:
\[
\boxed{
c(x)=c(y)
\Longrightarrow
x\equiv_{\mathcal H}y.
}
\]

If an exact coarsest quotient exists:
\[
c(x)=c(y)
\iff
x\equiv_{\mathcal H}y.
\]

\subsection{Continuation sufficiency}

For every protected continuation \(a\), require a descended continuation \(\bar a\) as specified earlier.


\subsection{Regenerative preservation}

When a quotient or fold replaces an actively represented form \(z\) by \(m\), and protected future use may require reopening \(z\), require
\[
\boxed{
Regen_{\mathcal H}(m,z)
}
\]
or an explicit reacquisition contract whose loss is already admitted by an approximate licence.  Present-output sufficiency alone is not enough to call the retained form regeneratively sufficient.

\subsection{Recovery/reopening contract}

The quotient must retain enough:
\begin{itemize}
\item provenance;
\item residual distinctions;
\item factorization route;
\item recovery or reacquisition route;
\item unlock triggers
\end{itemize}
for every protected future requirement of the licence.

\section{Compression license}

A canonical exact compression licence is:
\[
\boxed{
\Lambda_c
=
(
c,\mathcal H,\mathcal A,
\mathsf{Scope},
\mathsf{Recovery},
\mathsf{Residual},
\Unlock(c),
\mathsf{Evidence}
).
}
\]

Here \(\mathcal A\) is the protected continuation/operator family.

No quotient is called exact merely because current tests failed to distinguish its classes.

\section{Approximate compression}

When exact quotienting does not satisfy the resource contract, a binding may define:
\[
\widetilde c:X\to\widehat X
\]
and consequence distortion:
\[
\boxed{
\Delta_{\mathcal H}(x,\widetilde c(x)).
}
\]

An approximate licence must state:
\begin{enumerate}
\item the protected consequence family;
\item the distortion relation or preorder;
\item admissible error;
\item resource gain;
\item residual/error structure retained;
\item reopening triggers;
\item whether continuation behavior is also approximated.
\end{enumerate}

The distortion relation or preorder must preserve the direction of the licensed
inference.  An over-approximation, an under-approximation, and a binding-specific
approximation with the same scalar error are not interchangeable merely because that
number agrees.  Any implementation encoding approximation by a scalar alone must be
rejected when the directional contracts discharge different protected claims.

Approximate compression may never be reported as exact protected equivalence.

\section{Unlock field}

\begin{definition}[Unlock]
\[
\boxed{
\Unlock(c)
=
\Unlock_{\mathrm{obs}}(c)
\cup
\Unlock_{\mathrm{dyn}}(c)
\cup
\Unlock_{\mathrm{ctx}}(c).
}
\]

\[
\Unlock_{\mathrm{obs}}(c)
=
\{
K:
\Con_K
\text{ no longer factors through }c
\}.
\]

\[
\Unlock_{\mathrm{dyn}}(c)
=
\{
a:
a
\text{ no longer descends through }c
\}.
\]

\(\Unlock_{\mathrm{ctx}}\) contains binding, scope, applicability, assumptions, authority, protection, or recovery-contract changes that invalidate the original licence.
\end{definition}

Thus:
\[
\boxed{
\text{separator-triggered reopening}
\subseteq
\text{unlock-triggered reopening}.
}
\]

\section{Reopening}

If a trigger \(u\in\Unlock(c)\) becomes active:
\[
\boxed{
\Reopen(c,u).
}
\]

Reopening:
\begin{enumerate}
\item preserves the old quotient as historical ancestry;
\item locates the smallest implicated identification;
\item restores residue/provenance/recovery information;
\item refines the active representation;
\item records the new protection or continuation that broke the old licence.
\end{enumerate}

Reopening is not epistemic amnesia and does not require deleting a once-lawful compression from history.

\section{Pattern separation and completion}

If retained events are similar under a working representation but have protected-different consequences:
\[
Similar(e_i,e_j)
\land
\Con_{\mathcal H}(e_i)\neq\Con_{\mathcal H}(e_j),
\]
then the representation must separate them.

Conversely, pattern completion across compressed histories is licensed only when all completions still live under the reconstruction are protected-equivalent or when an explicit approximate licence permits the ambiguity.

\part{Question Frontiers, Productivity, and Search}

\section{Semantic question universe}

Fix \(A:\mathsf{Ty}_{\B}\) and a source-bearing presentation
\(\Sigma:\mathsf{SourceConfig}_{\B}(A)\).  Let:
\[
Q^\infty(\Sigma)
\]
denote the mathematical set of all well-typed questions constructible from the current relation/operator presentation, open-port grammar, admitted bindings, and unbounded semantic composition.

It need not be computably enumerable.

\begin{theorem}[Conservative question-language growth]\label{thm:conservative-question-growth}
Suppose a presentation change \(j:\Sigma_t\hookrightarrow\Sigma_{t+1}\) is a
conservative extension: every old type, relation, question constructor, typing rule,
and protected interpretation remains available through \(j\).  Then structural
transport induces an embedding
\[
\boxed{
\operatorname{EmbedQ}:Q^\infty(\Sigma_t)\hookrightarrow Q^\infty(\Sigma_{t+1}).
}
\]
The image inclusion is strict only when the successor admits at least one well-typed
question not in the image of \(\operatorname{EmbedQ}\).  Rebinding, replacement, or semantic revision
is not necessarily conservative, so no unqualified inclusion of question universes is
licensed for those changes.
\end{theorem}

\begin{proof}
Transport each old question syntax tree and its typing derivation through \(j\).
Preservation of constructors, typing, and interpretation makes the transport
well-defined and injective.  Strictness requires a successor question without an old
preimage.  If old constructors or meanings are removed, the transport need not exist.
\end{proof}

\section{Executable question universe}

An effectivity profile is:
\[
\boxed{
\epsilon=
(
b_G,b_D,b_T,
\mathsf{Deciders},
\mathsf{SemiDeciders},
\mathsf{Policy}
).
}
\]

A finite/resource-bounded generator produces:
\[
Q^\epsilon(\Sigma)\subseteq Q^\infty(\Sigma).
\]

Failure to generate a question under \(\epsilon\) is not proof that no such semantic question exists.

The corresponding checked occurrence universes are the dependent subtypes
\[
\boxed{
\mathsf{AskOcc}^{\infty}(\Sigma)
=
\left\{
\AskRef:\mathsf{AskOcc}_A(\Sigma)
\mid
\mathsf{AskQuestion}(\AskRef)\in Q^\infty(\Sigma)
\right\},
}
\]
\[
\boxed{
\mathsf{AskOcc}^{\epsilon}(\Sigma)
=
\left\{
\AskRef\in\mathsf{AskOcc}^{\infty}(\Sigma)
\mid
\mathsf{AskQuestion}(\AskRef)\in Q^\epsilon(\Sigma)
\right\}.
}
\]
Write \(q_{\AskRef}:=\mathsf{AskQuestion}(\AskRef)\).  The superscripts classify existing checked
source occurrences; they do not introduce new occurrence species.

\section{Derived question availability and required discharge}

\textbf{Status: DERIVED.}  For a checked occurrence \(\AskRef\) with question \(q\), distinguish:
\begin{itemize}
\item \(\Formable_{\B,\Sigma}(q)\): the relation, bindings, guards, and open carriers are
      representable and well typed;
\item \(\Applicable_{\B,\Sigma}(q)\): the question applies at the declared scope, grain, binding,
      version, and horizon;
\item \(\Executable_{\B,\Sigma,\epsilon}(q)\): every open port has a lawful route within the
      current effectivity horizon;
\item \(\Answerable_{\B,\Sigma}(q)\): the declared semantic/domain state admits at least one
      supported completion;
\item \(\Productive_{\mathcal H}(\Sigma,\AskRef)\): supported answers can induce
      protected-different occurrence-specific successors;
\item \(\ResolvedQ_{\mathcal H}(\Sigma,\AskRef)\): all live supported answers induce
      protected-equivalent successors under sufficient declared coverage; and
\item \(\Ready^{\chi}_{\Sigma}(q)\): formability, applicability, and the declared use contract
      \(\chi\) hold.
\end{itemize}
These predicates do not collapse.  A formable check can be nonexecutable; an executable exhaustive
probe can establish an empty completion field; incomplete coverage cannot establish resolvedness.

\begin{definition}[Required discharge]\label{def:required-discharge}
\(\RequiredDischarge(\Sigma,\AskRef,d)\) holds when an explicit standing or source-program
dependency \(d\) cannot lawfully continue until the open relation of \(\AskRef\) is discharged through
its declared \(\mathsf{Probe}\), \(\mathsf{Check}\), \(\mathsf{Warrant}\), exact support, or
reconstruction route.  It is not a new question species and does not assert that the obligation is
executable or successfully discharged.
\end{definition}

\begin{law}[Productivity and required discharge differ]\label{law:productive-required-discharge}
A singleton-result check or warrant may be nonproductive as a discretionary discriminator while
remaining obligatory for lawful continuation.  Conversely, a productive generated question is
not thereby required or authorized.  Therefore a controller may omit a continuation-equivalent
question only when no explicit dependency requires its discharge.
\end{law}

\section{Question productivity}

Let \(\QStep^\ast(\Sigma,\AskRef,\widehat S)\) be the finite or binding-native successor behavior
after supplying proof-carrying supported answer \(\widehat S\) to the exact occurrence
\(\AskRef\).

An occurrence is productive relative to \(\mathcal H\) when there exist lawful supported answer
records \(\widehat S_1,\widehat S_2\in
\mathsf{SuppAns}(\mathsf{AskQuestion}(\AskRef))\) such that:
\[
\boxed{
\QStep^\ast(\Sigma,\AskRef,\widehat S_1)
\not\equiv_{\mathcal H}^{\mathcal P}
\QStep^\ast(\Sigma,\AskRef,\widehat S_2).
}
\]

The occurrence index is required: the same normalized \(q\) and return class under another checked
continuation may have a different successor.

The lifting \(\equiv_{\mathcal H}^{\mathcal P}\) for plural/nondeterministic successors is binding-supplied.  The reference finite symmetric lifting is:
\[
\widehat{\Con}(K,U)
=
\{\Con_K(u):u\in U\},
\]
and:
\[
U\equiv_{\mathcal H}^{\mathcal P}V
\iff
\forall K\in\mathcal H,\,
\widehat{\Con}(K,U)=\widehat{\Con}(K,V).
\]

Bindings may instead supply Hoare, Smyth, Plotkin, probabilistic, trace, game, or concurrent liftings.

\section{Semantic and executable frontiers}

For a binding-supplied resource/risk preorder \(\preceq\), let \(\ND_{\preceq}(C)\) be the
members of \(C\) not strictly dominated by another member.  Define
\[
\mathsf{Req}_{\Sigma}(C)
=
\{\AskRef\in C:\exists d.\RequiredDischarge(\Sigma,\AskRef,d)\},
\]
\[
\boxed{
\ND_{\preceq,\Sigma}^{req}(C)
=
\mathsf{Req}_{\Sigma}(C)
\cup
\ND_{\preceq}(C).
}
\]
When no preorder is supplied, \(\ND_{\Sigma}^{req}(C)=C\).  A stronger rule may remove a required
occurrence only with a typed proof that another retained occurrence discharges the same dependency;
ordinary cost dominance never deletes it.  Since nondominance is still computed over all of \(C\),
a required occurrence may dominate and remove an optional occurrence.  Requiredness protects its
bearer from removal; it does not protect unrelated optional work from comparison against it.

Define:
\[
\boxed{
\begin{aligned}
\mathcal C^\infty_{\mathcal H}(\Sigma)
=\bigl\{\AskRef\in\mathsf{AskOcc}^\infty(\Sigma):{}&\\[-1mm]
&\Formable_{\B,\Sigma}(q_{\AskRef})
 \land\Applicable_{\B,\Sigma}(q_{\AskRef})\\
&\land\bigl(
  (\Productive_{\mathcal H}(\Sigma,\AskRef)
   \land\neg\ResolvedQ_{\mathcal H}(\Sigma,\AskRef))\\
&\hspace{21mm}\lor
  \exists d.\,\RequiredDischarge(\Sigma,\AskRef,d)
\bigr)\bigr\}.
\end{aligned}
}
\]
and
\[
\boxed{
Frontier^\infty_{\mathcal H}(\Sigma)
=\ND_{\preceq,\Sigma}^{req}\bigl(\mathcal C^\infty_{\mathcal H}(\Sigma)\bigr).
}
\]

Define the executable classification:
\[
Class^\epsilon(\Sigma,\AskRef)
\in
\{
\mathsf{productive},
\mathsf{idle},
\mathsf{unknown}
\}.
\]

Then:
\[
\boxed{
\begin{aligned}
\mathcal C^\epsilon_{\mathcal H}(\Sigma)
=\bigl\{\AskRef\in\mathsf{AskOcc}^\epsilon(\Sigma):{}&\\[-1mm]
&\Formable_{\B,\Sigma}(q_{\AskRef})
 \land\Applicable_{\B,\Sigma}(q_{\AskRef})\\
&\land\Executable_{\B,\Sigma,\epsilon}(q_{\AskRef})\\
&\land\bigl(Class^\epsilon(\Sigma,\AskRef)=\mathsf{productive}
  \lor\exists d.\,\RequiredDischarge(\Sigma,\AskRef,d)\bigr)
\bigr\}.
\end{aligned}
}
\]
and
\[
\boxed{
Frontier^\epsilon_{\mathcal H}(\Sigma)
=\ND_{\preceq,\Sigma}^{req}\bigl(\mathcal C^\epsilon_{\mathcal H}(\Sigma)\bigr).
}
\]

A required but nonexecutable obligation remains an explicit \(\mathsf{Blocked}\),
\(\mathsf{Unknown}\), or resource residual outside the executable frontier; it is not erased as
idle.  A resolvedness claim cannot suppress an undischarged obligation.

\begin{law}[Unknown legality]
Failure to find a separator, path, counterexample, proof, or useful question under current resources does not establish equivalence, impossibility, necessity, or irrelevance.
\end{law}

\section{Derived interrogative algebra}\label{sec:derived-interrogative-algebra}

\textbf{Status: DERIVED.}  Natural-language question roots and route classifications are
transparent normal forms over the canonical relation/question/source-program language.  They add
no interrogative ontology, actuality species, authority route, or scheduler primitive.

\subsection{Root normal forms}

The retained derived root family is
\[
\boxed{
\Omega_Q=
\{\mathsf{Expose},\mathsf{Orient},\mathsf{Factor},
\mathsf{Polarize},\mathsf{Vary},\mathsf{Ground}\}.
}
\]
Its meanings are:
\begin{itemize}
\item \(\mathsf{Expose}_I(R,\beta)=?_I R[\beta]\), the canonical question constructor;
\item \(\mathsf{Orient}^{+}\) and \(\mathsf{Orient}^{-}\), exposure of an admitted relation in
      forward or converse orientation;
\item \(\mathsf{Factor}\), exposure of factors or mediators through join, composition, hiding, and
      open ports;
\item \(\mathsf{Polarize}\), construction of a binding-appropriate positive alternative,
      classical breaker, separator, or already admitted positive-negation field;
\item \(\mathsf{Vary}\), opening a binding-supplied admissible removal, substitution,
      intervention, perturbation, strengthening, weakening, reversal, or reordering relation; and
\item \(\mathsf{Ground}\), reifying a claim/relation/question obligation and opening its ordinary
      support, check, or warrant relation.
\end{itemize}
Reification is ambient recursive form closure rather than a seventh root.  Questions about
questions reify the checked question occurrence and apply the same roots.

\begin{law}[No universal polarization]\label{law:no-universal-polarization}
Logical complement, counterexample, alternative design, separator, and contextual
positive-negation incidence remain different relations unless a binding establishes their
correspondence.  \(\mathsf{Polarize}\) cannot manufacture a \(\mathsf{NegationUse}\), exterior,
departure witness, or standing breaker.
\end{law}

Operational names such as \textsf{Contrast}, \textsf{Backchain}, \textsf{Preimage},
\textsf{Project}, \textsf{Ablate}, \textsf{Substitute}, \textsf{Localize}, \textsf{Construct},
\textsf{Scrutinize}, \textsf{WhyNot}, and \textsf{HowCan} are transparent compositions of these
roots.  Deleting a name and expanding its definition must preserve typing, applicability, whole
supported-answer behavior, authority, failure exits, provenance, and reopening.

\subsection{Static pair discipline}

Keep distinct:
\[
\boxed{
\operatorname{PairwiseDistinct}
\left(
\begin{gathered}
\text{inverse},\ \text{converse},\ \text{logical breaker dual},\\
\text{binding-supplied adjoint},\ \text{same-use reciprocal return}
\end{gathered}
\right).
}
\]
A displayed \(\operatorname{PairwiseDistinct}\) obligation means that every two applicable members
must retain different relation identity and cannot be substituted without a binding-supplied
equivalence theorem preserving the protected continuation.
A many-to-one relation has a lawful converse question with a non-singleton reverse fiber but no
strict inverse.  An adjoint or weakest-condition question requires a checked binding-supplied law,
not merely converse incidence.  Reciprocal return requires the exact admitted oriented use and its
departure, coverage, and fiber provenance.

\subsection{Unlock and route explanations}

\begin{definition}[Answer-conditioned unlock]\label{def:q-unlock}
For target use contract \(\chi\),
\[
\boxed{
\Unlock^{\chi}(\Sigma;\AskRef,\widehat S\Rightarrow r)
\iff
\neg\Ready^{\chi}_{\Sigma}(r)
\land
\Ready^{\chi}_{\QStep(\Sigma,\AskRef,\widehat S)}(r).
}
\]
The supported answer may supply a dependent index, guard, foil, mediator, probe, support/warrant
route, representation, or reciprocal seed.  It does not answer \(r\) merely by making it ready.
\end{definition}

Route annotations such as narrowing, discharge, defeat, reframe, reorientation, or extension state
why an ordinary continuation is useful.  They may overlap and are not a runtime enum.  A defeat of
standing requires sufficient authority; a reframe uses ordinary representation/binding revision;
reorientation retains the canonical seed relation; extension requires its own accepted successor.
Mere adjacency does not establish \(\Unlock\), and deleting a route label cannot change execution.

\subsection{Root frontier and local closure}

For reified held structure \(W\), root family \(\Omega\), and effectivity horizon \(\epsilon\),
define
\[
\boxed{
\begin{aligned}
\RootFrontier^{\epsilon}_{\mathcal H}(W,\Sigma)
=\ND^{req}_{\preceq,\Sigma}\!\bigl(\{\AskRef:{}&
              \AskRef\text{ is produced by an admitted root over }W,\\
              &\Formable_{\B,\Sigma}(q_{\AskRef})
               \land\Applicable_{\B,\Sigma}(q_{\AskRef})
               \land\Executable_{\B,\Sigma,\epsilon}(q_{\AskRef})\\
              &\land(\Productive_{\mathcal H}(\Sigma,\AskRef)
               \lor\exists d.\,\RequiredDischarge(\Sigma,\AskRef,d))
               \}\bigr).
\end{aligned}
}
\]
The required-safe nondominated operator is relative to a binding-supplied resource/risk preorder and
becomes the identity when none is supplied.  Cost dominance never removes an occurrence carrying an
undischarged dependency.  The roots generate candidates; they are not a wheel to execute
ritually.  After a return, reconstruct the residual and frontier.

\begin{definition}[Local interrogative fixed point]\label{def:local-ifp}
\(\IFP_{\Omega,\mathcal H,\epsilon}(W,\Sigma)\) holds under explicit coverage when every relevant
root-generated occurrence is determined by the retained profile, factorably redundant,
inapplicable, protected-continuation-equivalent with no required discharge, explicitly blocked,
resource-bounded, or dependent on a represented language/probe/binding extension.  No executable,
applicable, nonredundant productive occurrence or executable required discharge remains.
\end{definition}

\begin{law}[Local closure reopens]\label{law:local-ifp-reopens}
Local interrogative closure is not global completion.  A new separator, relation, probe, foil,
binding, representation, protected horizon, or resource regime may make a new occurrence formable,
executable, productive, or required and thereby reopen the fixed point.
\end{law}

\subsection{Question-route occurrence and regeneration}

A compact derived route record is
\[
\boxed{
\mathsf{QRouteOcc}
=
\langle
\AskRef,E,r,\ResPath,\widehat S,\AskRef',\chi,
\Delta_{in},\Delta_{out},\mathcal H,
\mathsf{modes},\mathsf{coverage},\mathsf{prov}
\rangle,
}
\]
with absent actual fields left explicitly nonactual for pure/generated routes.  It points to
canonical program and event ancestry and is not a second history.  The occurrence identity prevents
two equal semantic \((q,|\widehat S|)\) projections with different occurrences, continuations, or
support ancestry from collapsing.

Removing a supported-answer position yields the occurrence-specific cue fiber
\[
\{\widehat S:\QSucc(\AskRef,\widehat S,q')\},
\]
while removing a question position yields a fiber of checked occurrences, not merely surface
question forms.  Regeneration is exact only when the appropriate fiber is one protected class and
its typing, authority, event/path provenance, and reopening obligations survive.  Ambiguity remains
a separator residual.

A recurrent route may be proposed as a transparent method containing its trigger residual, source
\(\mathsf{IProg}\), applicability, protected gain, failure residuals, coverage, provenance, and
reopen condition.  Before folding it, ablate or lawfully bypass the question occurrence and seek a
protected continuation that distinguishes the original and reduced programs.  Frequency alone
cannot warrant the method, and an answer-independent host policy cannot replace the stored
continuation.

\begin{law}[No second question language or history]\label{law:no-second-question-language}
Root syntax, route labels, local-frontier views, and \(\mathsf{QRouteOcc}\) are authoring,
analysis, or controller projections.  Their conservative expansion contains only ordinary typed
relations, questions, \(\mathsf{Return}_I\), and occurrence-specific \(\mathsf{Ask}\), and their
realized evidence remains ordinary event history.
\end{law}

\section{Relative completeness of the question grammar}

\begin{theorem}[Bounded separator completeness]
Assume:
\begin{enumerate}
\item the current bounded representational repertoire contains every available protected distinguishing context;
\item context application is expressible;
\item every nonempty typed open port set may be exposed as a question;
\item well-typed question enumeration is fair.
\end{enumerate}
If two live completions \(x,y\) are distinguishable by the bounded repertoire, then the generated question language contains a question that separates them.
\end{theorem}

\begin{proof}
By assumption, there exists an expressible protected context \(K\) with distinct consequences on \(x,y\).  Reify the relation schema for \(K\)'s relevant output and expose the corresponding output/discriminator port.  Fair enumeration eventually generates that well-typed open relation, whose support profile differs on \(x,y\).
\end{proof}

This theorem is relative to the binding, represented repertoire, protection horizon, and resource assumptions.  Acquisition of a new relation/context can lawfully reopen an old quotient.

\section{Question selection}

\begin{definition}[Sufficient discriminator basis]\label{def:separator-basis}
Let \(Z\) be a live residual field and let \(C\) be a family of supported, applicable,
well-typed discriminators with declared answer semantics and coverage.  The family is
sufficient for \(Z\) at horizon \(\mathcal H\) when
\[
\boxed{
z_1\not\equiv_{\mathcal H}z_2
\Longrightarrow
\exists c\in C.\,
\mathsf{Ans}_c(z_1)\neq\mathsf{Ans}_c(z_2).
}
\]
For exact total deterministic answer signatures
\(A_C:Z\to\prod_{c\in C}\mathsf{Ans}(c)\), this is equivalent to
\[
\boxed{\ker A_C\subseteq{\equiv_{\mathcal H}}.}
\]
Given a binding-supplied resource preorder, the smallest sufficient bases form
\[
\operatorname{Min}_{\preceq}
\{C:\mathsf{SepBasis}_{\mathcal H}(C,Z)\}.
\]
Several incomparable minimal bases may induce the same protected quotient.  Under
partial generation, observation, or coverage, failure to find such a basis remains
\(\mathsf{Unknown}\) or resource-bounded rather than proving that none exists.
\end{definition}

The calculus does not impose a universal scalar objective.  When several executable questions are incomparable, retain a nondominated frontier over binding-supplied dimensions such as:
\begin{itemize}
\item expected discrimination;
\item probe cost;
\item latency;
\item risk;
\item reversibility;
\item authority strength;
\item information gain;
\item reuse value;
\item reopening value.
\end{itemize}

A scheduling policy is a binding/runtime policy, not part of the semantic kernel.


\part{Prediction, Native Methods, and Control}

\section{Prediction}

Prediction is a represented forward consequence, not actuality.

For:
\[
T:X\rightsquigarrow Y
\]
and target predicate \(Q\subseteq Y\):
\[
\boxed{
Possible_T(Q,x)
\iff
\exists y.\,T(x,y)\land Q(y),
}
\]
\[
\boxed{
Necessary_T(Q,x)
\iff
\forall y.\,T(x,y)\Rightarrow Q(y).
}
\]

The universal precondition transformer is:
\[
\boxed{
Pre_T(Q)
=
\{x:
\forall y.\,T(x,y)\Rightarrow Q(y)
\}.
}
\]

\section{Prediction seal}

When a consequential probe is being used to discriminate a prediction, seal the prediction before the return:
\[
\boxed{
\mathsf{PredictionSeal}
=
(
\mathsf{id},
\mathsf{attemptRef},
\widehat R,
\widehat C,
\epsilon,
\mathsf{scope},
\mathsf{basisRefs}
).
}
\]

The seal satisfies:
\[
\boxed{
\mathsf{PredictionSeal}
\prec_{\mathrm{ledger}}
\mathsf{ExternalReturn}.
}
\]

Once actualization begins producing a return, the prediction, its assumptions, and its acceptance relation may not be retrospectively rewritten for that test.

\section{Mismatch}

A mismatch is a typed relation between a sealed prospective prediction and an actual preserved return:
\[
\boxed{
Mismatch(\pi,r).
}
\]

Mismatch does not by itself identify which premise failed.  It generates localization questions over:
\begin{itemize}
\item scope;
\item applicability;
\item relation/factorization;
\item transition;
\item decoder;
\item representation;
\item hidden dependency;
\item binding.
\end{itemize}

\section{Success attacks necessity}

A confirming return does not end scrutiny.  Failure attacks sufficiency:
\[
?x[P(x)\land\neg Q(x)].
\]
Success still attacks necessity:
\[
?x[\neg P(x)\land Q(x)].
\]

Thus:
\[
\boxed{
\textbf{FAILURE ATTACKS SUFFICIENCY;
SUCCESS STILL ATTACKS NECESSITY.}
}
\]

\section{Description and control}

A relation that predicts or describes a consequence is not automatically a control relation.

If:
\[
D(x,y)
\]
describes a relation and:
\[
C(u,x,y)
\]
claims an intervention/control guarantee, then:
\[
\boxed{
D\not\Rightarrow C
}
\]
without an independently warranted transition/control bridge.

When a controller is synthesized over an abstraction, the abstraction-to-concrete refinement relation must preserve the protected guarantee.  This is binding-specific and may use native feedback-refinement, simulation, refinement, or other established machinery.

\section{Native method contract}

A native method is registered by a compact typed contract:
\[
\boxed{
\begin{aligned}
\mathcal M=(
&\mathsf{id},R_m:X\rightsquigarrow Y,
\mathsf{Applicability},\mathsf{Law},\mathsf{Coverage},\mathsf{Authority},\\
&\mathsf{ExtensionDomain},\mathsf{Backend},\mathsf{Checker},\mathsf{Cost},
\mathsf{FailureSchemas},\mathsf{Provenance}).
\end{aligned}
}
\]

The generic inquiry compiler should use a mature native method when its input/output relation lawfully discharges the open relation.  The calculus does not reproduce domain algorithms merely because inquiry can describe them.

A method is \emph{admitted} when this contract has standing acceptance, \emph{runnable}
when its declared backend is available under the current resource and binding state,
and \emph{usable} for a particular open port only when it is admitted, runnable,
applicable, type-compatible, sufficiently covered, and authorized for that port.
None of these predicates implies either of the others without the stated evidence.

\begin{law}[Failure-closed method boundary]\label{law:method-residual}
A method result re-enters the ordinary return/decode/check/warrant lifecycle.  When
actual execution occurs, its raw return is preserved before classification.  A checked
semantic non-discharge---for example, a certified empty solution field or a retained
counterexample---may produce a typed residual or an exact terminal status.  Backend
absence, crash, timeout, malformed return, or undecided checking does not establish
semantic emptiness, impossibility, or contradiction; it remains an operational
\(\mathsf{Blocked}\), \(\mathsf{ResourceBounded}\), or \(\mathsf{Unknown}\) result as
applicable.  Residual handling and resumption compile through ordinary first-order
answer-dependent programs and introduce no new runtime effect.
\end{law}

\part{Learning, Patches, and Self-Revision}

\section{Reasoning presentation}

At version \(t\), define the reifiable presentation:
\[
\boxed{
\begin{aligned}
\mathsf{Pres}_t
=(&\mathsf G_t,\mathsf E_t,\mathsf{RelSchemas}_t,\mathsf{TypeRules}_t,\\
 &\mathsf{QuestionRules}_t,\mathsf{ProbeContracts}_t,\mathsf{Compiler}_t,\\
 &\mathsf{Bindings}_t,\mathsf{Protection}_t).
\end{aligned}
}
\]

Here:
\begin{itemize}
\item \(\mathsf G_t\) is the admitted typed generator/operator repertoire;
\item \(\mathsf E_t\) is the admitted family of equations, inequations, applicability constraints, and composition laws.
\end{itemize}

\section{Patch roles}

A candidate change is:
\[
\boxed{
\pi:\PatchIR[k].
}
\]

Roles include:
\[
k\in
\{
\mathsf{Semantic},
\mathsf{Traversal},
\mathsf{Compiler},
\mathsf{Binding},
\mathsf{Protection},
\mathsf{Implementation}
\}.
\]

\subsection{Semantic patch}

Changes standing semantic relations or their support/warrant status.

\subsection{Traversal patch}

Changes question/probe repertoire, operator paths, routing policy, learned methods, or active-view construction without asserting a new fact about the external domain.

\subsection{Compiler patch}

Changes lowering/rendering/decoding implementation while preserving semantic operator contracts or explicitly changing them through a governed successor.

\subsection{Binding patch}

Adds or changes native types, relations, actualization channels, comparison bridges, or composition laws.

\subsection{Protection patch}

Changes the protected horizon.  It is particularly sensitive because it may invalidate earlier folds.  It may not be silently changed by an optimization.

\section{Primitive presentation changes}

A warranted successor may:
\begin{itemize}
\item add a generator/operator;
\item add a standing equation/merge;
\item split an old equation after a new separator;
\item add a question/open-port operation;
\item refine an answer frame;
\item add a probe contract;
\item promote a composite method;
\item import a native method;
\item rebind;
\item change a compiler;
\item change the protected horizon under independent governance.
\end{itemize}

All are represented as ordinary typed patches.

\section{Predecessor-judged successor}

For presentations \(R,R'\), define:
\[
\boxed{
R'\succ_{\mathcal H}R
}
\]
when:
\begin{enumerate}
\item \(\Preserve_{\mathcal H}(R',R)\);
\item every removed predecessor relation/obligation has an explicit warranted disposition;
\item \(\Gain(R',R)\) holds under a declared preorder or \(R'\) repairs a demonstrated defect;
\item regression tests for protected capabilities pass;
\item the support for accepting \(R'\) does not depend solely on claims first introduced by \(R'\).
\end{enumerate}

Improvement is generally a partial order/frontier.  Incomparable successors remain incomparable.

\section{Self-inquiry}

Reify:
\[
\ulcorner\mathsf{Pres}_t\urcorner.
\]

Then ordinary questions may target its relations:
\[
?g[
Extension(\mathsf{Pres}_t,g)
\land
Expresses(\mathsf{Pres}_t+g,\rho)
].
\]

Or:
\[
?K[
Protected(K)
\land
\Con_K(\mathsf{Pres}_t+g)
\neq
Required(K)
].
\]

The same question, probe, return, support, warrant, and successor rules apply.  No privileged meta-language is introduced.

\begin{law}[Anti-oracle]
A candidate revision may propose its own discriminator, proof strategy, probe, or successor.  It may not count its own generated proposal as the independent return required to warrant its own promotion.
\end{law}

\part{Runtime Infrastructure}

\section{Canonical runtime state}

A reference runtime may expose:
\[
\boxed{
X_t=
(
H_t^{evt},
H_t^{acc},
A_t,
\Pi_t,
\Stand_t,
Prot_t,
\Gamma_t,
\Omega_t,
\epsilon_t,
sched_t
).
}
\]

Where:
\begin{itemize}
\item \(H_t^{evt}\): authoritative actual event history;
\item \(H_t^{acc}\): accepted patch history;
\item \(A_t\): versioned accepted presentation;
\item \(\Pi_t\): derived working presentation/context;
\item \(\Stand_t\): least-fixed-point standing relations;
\item \(Prot_t\): protected future declaration;
\item \(\Gamma_t\): live typed environment/referents;
\item \(\Omega_t\): open inquiry obligations;
\item \(\epsilon_t\): effectivity/resource profile;
\item \(sched_t\): non-authoritative scheduling policy/state.
\end{itemize}

The tuple is an implementation seam, not an ontology.

\section{Canonical records}

\subsection{Relation schema record}

\begin{lstlisting}[basicstyle=\ttfamily\small]
RelSchema {
    id
    binding_version
    named_ports[]       // (name, TypeId)
    body_ir
    laws[]
    provenance[]
}
\end{lstlisting}

\subsection{Relation-use and reciprocal records}

These are conceptual typed records.  They do not require new runtime effects or authoritative database tables; content-addressed implementations may represent them as ordinary artifacts and references.

\begin{lstlisting}[basicstyle=\ttfamily\small]
RelationUse {
    id
    relation_schema_ref
    typed_bindings{}
    orientation
    scope
    applicability
    grain
    support_ref
    warrant_ref
}

DeterminationPresentation {
    id
    distinction_ref
    orientation
    source_ref
    relational_web_ref
    scope
    applicability
    grain
    protected_horizon_ref
    support_ref
    status
    predecessor_ref
}

DepartureWitness {
    id
    distinction_ref
    source_ref
    candidate_ref
    source_presentation_ref
    source_observation_use_ref
    candidate_observation_use_ref
    source_answer_ref
    candidate_answer_ref
    incompatibility_use_ref
    support_ref
    scope
    applicability
    grain
}

NegationUse {
    id
    relation_use_ref
    distinction_ref
    orientation
    source_presentation_ref
    candidate_field_ref
    soundness_derivation_ref
    semantic_coverage
    scope
    applicability
    grain
    warrant_ref
}

SixfoldOccurrenceView {
    distinction_ref
    boundary_ref_optional
    source_x_ref
    negation_x_use_ref
    exterior_x_ref
    return_x_fiber_ref
    selected_return_x_ref_optional
    recovery_x_ref
    seed_use_ref
    source_y_ref
    negation_y_use_ref
    exterior_y_ref
    return_y_fiber_ref
    selected_return_y_ref_optional
    recovery_y_ref
    semantic_and_execution_coverage_refs[]
    event_refs[]
    residual_refs[]
    gamma_result_ref_optional
}
\end{lstlisting}

\subsection{Open query record}

\begin{lstlisting}[basicstyle=\ttfamily\small]
OpenQuery {
    id
    relation_schema_id
    bound_ports{}       // port -> typed reference/value
    open_ports[]        // nonempty
    discharge_modes{}   // open port -> mode
    scope
    applicability
    grain
    protected_horizon_ref
}
\end{lstlisting}

\subsection{Boundary chart}

\begin{lstlisting}[basicstyle=\ttfamily\small]
BoundaryChart {
    id
    query_id
    distinction_ref
    source_presentation_ref
    reciprocal_presentation_ref_optional
    plus_carrier
    boundary_carrier
    minus_carrier
    plus_projection
    minus_projection
    orientation
    candidate_status
    negation_x_use_refs[]
    negation_y_use_refs[]
    semantic_coverage_refs[]
    execution_coverage_refs[]
    seed_use_ref_optional
    return_fiber_refs[]
    recovery_refs[]
    reciprocal_status
    downstream_gamma_result_ref_optional
    traversal_relation
    grain
    protected_horizon_ref
}
\end{lstlisting}

\subsection{Probe operator}

\begin{lstlisting}[basicstyle=\ttfamily\small]
ProbeOperator {
    id
    query_id
    boundary_chart_id
    active_view_ref
    backend
    executable_code_ref   // prompt, tool call, method invocation, etc.
    return_type
    decoder_contract_ref
    probe_contract_ref
    compiler_version
}
\end{lstlisting}

\subsection{Probe contract}

\begin{lstlisting}[basicstyle=\ttfamily\small]
ProbeContract {
    id
    relational_role
    binding_version
    grain
    applicability
    comparator
    protected_horizon_ref
    decoder_version
    bridge_policy
}
\end{lstlisting}

\subsection{Actual event}

\begin{lstlisting}[basicstyle=\ttfamily\small]
ActualEvent {
    id
    ledger_parent
    state_before_ref
    source_ask_occurrence_ref_optional
    question_ref
    boundary_ref
    operator_ref
    raw_return_ref
    state_after_ref
    grain
    actualization_route
    backend_version
    provenance
}
\end{lstlisting}

\subsection{Operator occurrence}

\begin{lstlisting}[basicstyle=\ttfamily\small]
OperatorOccurrence {
    event_ref
    operator_ref
    state_before_ref
    raw_return_ref
    state_after_ref
    boundary_ref
}
\end{lstlisting}

\subsection{Claim}

\begin{lstlisting}[basicstyle=\ttfamily\small]
Claim {
    id
    proposition_or_payload
    source_question_ref
    supporting_return_refs[]
    resolution_path_refs[]
    scope
    applicability
    status            // candidate, checked, standing, rejected, unknown
}
\end{lstlisting}

\subsection{Support environment}

\begin{lstlisting}[basicstyle=\ttfamily\small]
SupportEnvironment {
    id
    claim_or_relation_ref
    premise_refs[]
    actual_return_refs[]
    checker_refs[]
    assumption_refs[]
    open_dependency_refs[]
    applicability
    scope
}
\end{lstlisting}

\subsection{Compression license}

\begin{lstlisting}[basicstyle=\ttfamily\small]
CompressionLicense {
    id
    quotient_or_fold_ref
    exact_or_approximate
    protected_horizon_ref
    protected_continuations[]
    scope
    evidence_refs[]
    residual_ref
    recovery_contract_ref
    unlock_conditions[]
    distortion_contract_ref
}
\end{lstlisting}

\subsection{Patch}

\begin{lstlisting}[basicstyle=\ttfamily\small]
Patch {
    id
    role
    predecessor_version
    changes[]
    preserved_obligations[]
    dispositions[]
    strict_gain_or_defect_ref
    evidence_refs[]
    regression_refs[]
    reopening_effects[]
}
\end{lstlisting}

\section{State storage and references}

The recommended architecture separates:
\[
\boxed{
\text{large retained state/history}
\quad+\quad
\text{small typed operator graph/index}.
}
\]

State snapshots and raw returns may be content-addressed or otherwise immutable.  Operators, prompts, relation schemas, and methods are stored once and occurrences point to them by reference.

This implements:
\[
\boxed{
\begin{gathered}
\textbf{THE STATE HISTORY IS THE MASS;}\\
\textbf{THE OPERATOR HISTORY IS THE INDEX.}
\end{gathered}
}
\]

The phrase is architectural, not semantic: the important law is that operator occurrences retain exact links to the states and returns in which they occurred.

\section{Checkpoint and replay}

A checkpoint is a performance optimization.  It is authoritative only relative to a replay/version contract.

A replayable checkpoint records:
\begin{itemize}
\item event-history prefix reference;
\item accepted-patch prefix reference;
\item compiler/binding versions;
\item deterministic fold version;
\item external returns by immutable reference;
\item any nondeterministic choices that must be replayed.
\end{itemize}

If replay semantics changes, the checkpoint requires an explicit migration/bridge.  Old raw events remain untouched.

\section{Concurrency}

The kernel does not impose total domain order.  Concurrent bindings may use partial orders, event structures, vector clocks, traces, or another native model.

Ledger append serialization, if used for persistence, must not be interpreted as causal succession unless the binding establishes that relation.

\part{Canonical Inquiry Process}

\section{One master regenerative recurrence}\label{sec:master-recurrence}

\textbf{Status: DERIVED CONTROL NORMAL FORM.}  The one overarching recurrence of v2.0 is
\[
\boxed{
\mathsf{BIND}
\to
\mathsf{OPEN}
\to
\mathsf{VARY}
\to
\mathsf{RETURN}
\to
\mathsf{DETERMINE}
\to
\mathsf{REFACTOR}
\circlearrowleft.
}
\]
These names summarize compositions already available in the canonical language.  They are not
runtime opcodes, new discharge modes, or a second program grammar.

Let a bound inquiry configuration \(\Delta_t\) retain the current binding and versions, accepted
presentation, ordinary event/revision ancestry, standing support, explicit source program and
environment, protected scope/grain/horizon, effectivity profile, live residual, and declared
coverage.  One recurrence has the following typed responsibilities.

\begin{description}
\item[\(\mathsf{BIND}\)] Reconstruct \(\Delta_t\) from accepted roots and actuality; establish the
      carriers, relations, arrangement, scope, applicability, authority, effectivity, and protected
      horizon.  Binding does not infer standing or causality from availability or order.
\item[\(\mathsf{OPEN}\)] Construct the smallest consequential checked \(\mathsf{Ask}\) occurrence
      \(\AskRef_t\) from the occurrence-aware live frontier.  It is retained because it is productive
      or because an explicit dependency requires discharge.
\item[\(\mathsf{VARY}\)] Construct or admit the lawful candidate field needed by that open relation:
      a contrast, separator, positive-negation exterior candidate, transformation, factor,
      representation, probe, or other binding-native variation.  Generation proposes occupants and
      never creates actuality, support, or warrant.
\item[\(\mathsf{RETURN}\)] Discharge each open component through its declared mode:
      \(\mathsf{Pure}\), \(\mathsf{Generate}\), \(\mathsf{Probe}\), \(\mathsf{Check}\), or
      \(\mathsf{Warrant}\).
      Preserve an actual raw return and ordinary event before resolution.  Preserve partial answer
      sets and per-component provenance.
\item[\(\mathsf{DETERMINE}\)] Resolve exactly what the return supports under the declared route and
      coverage as \(\zeta_t:\mathsf{ResolutionOutcome}(q_{\AskRef_t})\); check protected
      consequence, departure, recovery, support, or acceptance only where the relevant relation
      applies; and construct the explicit residual.  Exact emptiness, undefined resolution,
      unsupported completion, and unknown coverage remain distinct.
\item[\(\mathsf{REFACTOR}\)] If
      \(\zeta_t=\mathsf{Supported}(\widehat S_t)\), supply the whole proof-carrying answer to
      \(\kappa_{\AskRef_t}\).  For \(\mathsf{ExactEmpty}\), \(\mathsf{Undefined}\),
      \(\mathsf{Unsupported}\), or \(\mathsf{Unknown}\), construct the corresponding typed residual
      or justified stop without invoking that source continuation.  Then refine, fold, reopen,
      reconstruct, reorient, reconcile, extend, or rebind only as warranted; retain ancestry and
      unlocks; and produce the successor bound configuration \(\Delta_{t+1}\) when a lawful
      continuation exists.
\end{description}

Thus a realized step factors as
\[
\boxed{
\begin{aligned}
\Delta_t
\xrightarrow{\mathsf{BIND}}
\widehat\Delta_t
\xrightarrow{\mathsf{OPEN}}
\AskRef_t
\xrightarrow{\mathsf{VARY}}
\mathsf{route}_t\\
{}
\xrightarrow{\mathsf{RETURN}}
Y_t
\xrightarrow{\mathsf{DETERMINE}}
(\zeta_t,\mathcal R_t)
\xrightarrow{\mathsf{REFACTOR}}
\Delta_{t+1}\ \text{or a typed stop}.
\end{aligned}
}
\]
Here the per-port return carrier is the tagged sum
\[
\boxed{
Y_i=
\mathsf{PureResult}_i
+\mathsf{GeneratedProposal}_i
+\mathsf{ActualReturn}_i
+\mathsf{CheckedResult}_i
+\mathsf{WarrantResult}_i,
}
\]
selected only by \(\operatorname{mode}(q,i)\).  An actual-return summand contains an immutable
ordinary event and raw-return reference; resolution then contributes its path and supported
completion set.  Other summands do not fabricate an event.  A mixed question returns the dependent
product of its per-port summands and retains each route and provenance.  \(\mathcal R_t\) is the
explicit protected residual, not an alternate history.

Only
\[
\zeta_t=\mathsf{Supported}(\widehat S_t)
\quad\Longrightarrow\quad
\kappa_{\AskRef_t}(\widehat S_t)
\]
enters the occurrence-specific source continuation.  Every other resolution summand preserves its
own evidence and exits through a typed residual or stop; recurrence may later open a new question
from that residual, but it never fabricates \(\widehat S_t\).

\begin{law}[One recurrence, many specializations]\label{law:one-master-recurrence}
Every other loop, rhythm, clock, traversal, or repeated program in this specification is a
specialization, projection, binding, or rendering of
\cref{sec:master-recurrence}.  None constitutes a second controller or history.
\end{law}

\begin{law}[Guarded regenerative progress]\label{law:master-progress}
Recurrence crosses a new actual return, an admitted distinction, a representation/binding or
authority change, an observed state change, or a strict reduction of a declared finite frontier.
The fingerprint retains occurrence/continuation identity, bindings, evidence, repository/domain
actuality, horizon, coverage, and frontier.  The same supported answer in the same state may not
recur indefinitely; finite fuel yields \(\mathsf{ResourceBounded}\).
\end{law}

\section{Specializations of the master recurrence}

The positive-negation reciprocal program specializes the six positions as follows:
\[
\begin{array}{c|l}
\mathsf{BIND} & \text{present }W_D^X(S_X)\text{ and admitted oriented use }u_X\\
\mathsf{OPEN} & \text{open the use-tagged positive-negation relation}\\
\mathsf{VARY} & \text{construct a positively supported exterior candidate }O_X\\
\mathsf{RETURN} & \text{obtain the exterior and reverse section of the same use}\\
\mathsf{DETERMINE} & \text{check departure, the whole fiber, recovery, residuals, and coverage}\\
\mathsf{REFACTOR} & \text{establish }S_Y\text{ by seed/reorientation and rebuild the frontier.}
\end{array}
\]
The reciprocal orientation repeats through the same master recurrence.  The resulting sixfold and
\(\Gamma_D\) result are derived projections after both orientations; they are not another loop.

Hole solving, separator inquiry, representation/probe invention, standing closure, paired trace
reconstruction, sufficient-present updating, method learning, folding/reopening, and
predecessor-judged revision likewise occupy the six positions with their own typed relations and
authority.  The interrogative roots are available derived operators used within positions, not a
fixed root wheel.  The software-engineering clock is a binding-specific rendering of these
positions and roots, not semantic ontology.

\section{Stopping states}

Inquiry may relatively stop as:
\begin{itemize}
\item \(\mathsf{Satisfied}\): the current obligation is standing-satisfied under its declared scope;
\item \(\mathsf{Impossible}\): impossibility is independently warranted under the declared binding/scope;
\item \(\mathsf{Equivalent}\): remaining alternatives are exactly protected-equivalent under the admitted horizon;
\item \(\mathsf{Blocked}\): a required external return/capability is unavailable;
\item \(\mathsf{Unknown}\): executable inquiry cannot currently discriminate the live alternatives.
\item \(\mathsf{ResourceBounded}\): declared fuel or resource limits were reached with a live finite or open frontier remaining.
\end{itemize}

\(\mathsf{Unknown}\) is not failure of the calculus.  It is the lawful result when the semantic frontier is not executable or no warranted separator is currently available.

\part{Derived Structures}

\section{Problems}

A problem is not primitive.  Given desired consequence region \(G\subseteq X\), define problem region:
\[
P_G=X\setminus G
\]
when complement is meaningful in the chosen field, or more generally:
\[
Problem(x;G)\iff \neg G(x)
\]
inside a declared applicability field.

A solution is not ``the negative'' as an ontology.  It is any admissible succession/operator path that reaches a target satisfying the required consequence.

\section{Prerequisites}

For target relation \(T\), missing predecessor/dependency \(d\) is simply an open port:
\[
?d[
Required(d,T)
\land
\neg Discharged(d)
].
\]

Recursive prerequisite descent is therefore ordinary inquiry over dependency relations.

\section{Factorization}

For composition:
\[
h=g\circ f,
\]
factor discovery opens \(f\), \(g\), or a mediator.  Factorization is not a primitive question species.

\section{Necessity and sufficiency}

For candidate implication \(P\Rightarrow Q\):
\[
\boxed{
?x[P(x)\land\neg Q(x)]
}
\]
is the sufficiency breaker, and:
\[
\boxed{
?x[\neg P(x)\land Q(x)]
}
\]
is the necessity/converse breaker.

These are ordinary open relations generated by logical composition.

\section{Variation}

Let:
\[
V_\nu:X\rightsquigarrow X
\]
be a binding-supplied admissible variation relation and \(C\) a protected consequence.

Stretch:
\[
?y[
V_\nu(x,y)
\land
C(y)=C(x)
]
\]
seeks variation preserving consequence.

Squeeze/crossing:
\[
?y[
V_\nu(x,y)
\land
C(y)\neq C(x)
]
\]
seeks variation changing consequence.

Neither requires a universal metric or minimum.  If a cost/order is supplied, a nondominated or minimal frontier may be selected.


\section{Abstraction and hole reconstruction}

Abstraction is derived from ordinary partial binding.  Given a relational web \(W\), replacing a protected filling \(x\) by an open port yields \(\operatorname{Hole}_x(W)\).  Its lawful filler field is the indexed meet of the residual solution fibers.  If that field is one protected equivalence class, the web determines \(x\) at the protected grain; otherwise the surviving ambiguity generates a separator question.

\section{Representation and explanation}

Let \(\kappa:X\to K\) be a protected consequence and \(\eta:X\to S\) a candidate representation.  An exact explanation-by-representation over scope \(U\subseteq X\) is a factorization
\[
\boxed{
\kappa|_U=h\circ\eta|_U
}
\]
for some \(h:S\to K\).  The representation is insufficient when a pair in one \(\eta\)-cell has protected-different consequences.  It contains potentially unnecessary detail when a coarser representation still supports the same protected factorization.  Thus explanatory inquiry alternates refinement until consequence separates with coarsening until nonconsequential detail disappears.

\section{Linear dot-product consequence binding}

The following is a binding-specific exact specialization, not constitutional linear or
probabilistic structure.

\begin{theorem}[Consequence subspace]\label{thm:consequence-subspace}
Let \(\Omega=\mathbb R^d\), let \(P_Q\) have finite second moment
\[
M_Q=\mathbb E_{q\sim P_Q}[qq^\top],
\]
and protect the dot-product consequence family
\(\mathsf{Cons}_q(x)=q^\top x\) up to \(P_Q\)-almost-everywhere equality.  Then
\[
x\sim x'
\quad\Longleftrightarrow\quad
x-x'\in\ker M_Q.
\]
Hence the coarsest exact vector-space quotient is
\[
\boxed{\mathbb R^d/\ker M_Q,}
\]
and the induced map
\[
[x]\longmapsto M_Qx
\]
is a canonical linear isomorphism onto \(\operatorname{im}M_Q\).  Its vector-space
dimension is \(\operatorname{rank}M_Q\).
\end{theorem}

\begin{proof}
For \(v=x-x'\),
\[
\mathbb E[(q^\top v)^2]=v^\top M_Qv.
\]
The integrand is nonnegative, so the expectation vanishes exactly when
\(q^\top v=0\) almost everywhere.  Since \(M_Q\) is positive semidefinite,
\(v^\top M_Qv=0\) exactly when \(v\in\ker M_Q\).  The remaining claims follow from
the first isomorphism theorem and rank--nullity.
\end{proof}

\begin{corollary}[Minimal linear consequence dimension]
If \(h:\mathbb R^d\to\mathbb R^k\) is linear and exact for this consequence family,
then
\[
\ker h\subseteq\ker M_Q,
\qquad
k\geq\operatorname{rank}M_Q.
\]
The bound is attained by projection onto \(\operatorname{im}M_Q\), or equivalently by
any linear coordinate representation of the quotient.  No dimension lower bound is
claimed for arbitrary untyped, discontinuous, or otherwise unrestricted encodings.
\end{corollary}

\begin{remark}[Second moment, evidence, and reopening]
Centered covariance cannot replace \(M_Q\) unless the query mean is zero.  For the
deterministic query \(q=\mu\neq0\), covariance is zero while
\(M_Q=\mu\mu^\top\) retains the consequence-relevant mean direction.  The binding must
therefore preserve the query-distribution identity/version and the exactness evidence
for \(M_Q\).  A sample estimate or floating-point rank threshold supplies only a
working/approximate licence unless certified bounds establish the protected claim, and
a distribution change is an unlock condition for the quotient.
\end{remark}

\section{Cause}

Cause is not a primitive of the kernel.  A causal binding must supply stronger relations than correlation or succession: intervention semantics, counterfactual semantics, structural equations, potential outcomes, process mechanisms, or another native theory.  The Inquiry Calculus opens and tests those relations but does not infer causality from co-occurrence alone.

\section{Memory}

``Memory'' is a derived role of retained history, accepted folds, recovery relations, and question-conditioned projections.  It does not require a separate mutable semantic store.

\section{Method}

A method is a promoted/reified operator path with:
\[
(\mathsf{typed\ expansion},\mathsf{applicability},\mathsf{behavioral\ contract},
\mathsf{provenance},\mathsf{reopening}).
\]

\section{Attention}

Attention is not a primitive.  Question-conditioned activation/occlusion is:
\[
\eta_{q,g}:S\rightsquigarrow V_{q,g}.
\]
Calling the resulting active view ``attention'' is a binding/interpretive choice.

\section{Perceptual square and cube}

The square/cube constructions defined earlier may be used as explanatory or visualization devices.  They add no new semantic power beyond reciprocal boundary arrangement plus realized succession.

\part{Metatheory and Reference Mathematical Homes}

\section{Algebraic effects and interaction trees}

The core:
\[
\mathsf{Return}
\mid
\mathsf{Branch}
\mid
\mathsf{Probe}
\]
admits a free-effect interpretation in which probe requests are algebraic operations with return-indexed continuations.  Potentially unbounded inquiry can be modeled by interaction trees or related coinductive structures.

This correspondence is a reference model, not a new primitive requirement.

\section{Containers and dependent question signatures}

A question repertoire with answer family \(A(q)\) has container/polynomial form:
\[
\boxed{
P_{\mathcal Q}(X)
=
\sum_{q\in Q}X^{A(q)}.
}
\]

Question identity is the shape; complete answers are positions; answer-dependent continuations occupy those positions.

The open-relation account is more general because a raw return may support multiple semantic completions and because discharge authority is typed separately.

\section{Coalgebra and behavioral equivalence}

State-based recurring behavior can be represented:
\[
\gamma:X\to FX.
\]

Coalgebraic behavioral equivalence and partition refinement provide mature reference structures for exact retained state under repeated future interaction.

\section{Syntax/behavior compatibility}

Let \(T\) encode generated typed operator syntax and \(F\) encode observable behavior.  A distributive law:
\[
\boxed{
\lambda:TF\Rightarrow FT
}
\]
together with:
\[
\alpha:TX\to X,
\qquad
\gamma:X\to FX
\]
satisfying:
\[
\boxed{
\gamma\circ\alpha
=
F\alpha\circ\lambda_X\circ T\gamma
}
\]
is a strong metatheoretic target for proving that generated syntax composes compatibly with behavior.

It is not constitutionally required of every implementation.

\section{Abstract interpretation and refinement}

Exact/approximate representation, abstraction, refinement, factorization, consequence preservation, and reopening have native analogues in abstract interpretation, CEGAR, partition refinement, behavioral automata learning, simulation/refinement, and related methods.  These are method bindings and proof precedents, not parallel semantics.

\part{Conformance Laws}

\section{Static conformance}

A conforming implementation must:
\begin{enumerate}
\item type every relation schema and named port;
\item reject ill-typed joins/compositions;
\item record nonempty open-port sets for questions;
\item enforce per-port discharge authority;
\item distinguish raw returns from decoded completions;
\item preserve question/relational identity independently of surface wording.
\end{enumerate}

\section{Dynamic conformance}

A conforming runtime must:
\begin{enumerate}
\item preserve raw actual returns before interpretation;
\item record exact state/operator/return linkage at the declared grain;
\item keep ledger order separate from binding-defined domain succession;
\item retain resolution-path provenance when protected;
\item preserve \(\mathsf{Unknown}\) when the executable frontier cannot decide;
\item prevent generative branches from masquerading as actual events.
\end{enumerate}

\section{Question conformance}

Every active question must satisfy:
\begin{enumerate}
\item exact relation-schema and bound/open port identity;
\item scope/applicability;
\item answer carrier;
\item discharge mode;
\item local question equivalence;
\item boundary chart or an explicit statement that no two-sided local chart is presently formed;
\item active-view/occlusion contract;
\item residual/reopening route for any consequential occlusion.
\end{enumerate}


\section{Distinction and regenerative conformance}

A conforming implementation of a protected distinction must:
\begin{enumerate}
\item preserve the typed distinction schema and boundary-point identity;
\item retain both candidate side projections without treating them as exteriority evidence;
\item represent the live determination presentation with its version, support, scope, applicability, grain, and horizon;
\item require a positive, relevant, non-circular departure witness for every exact negation incidence;
\item preserve oriented negation-use identity and semantic coverage independently of execution coverage;
\item distinguish a selected return filling from the same-use reverse return fiber and record protected recovery;
\item generate the six oriented roles dependently with seed and use provenance when protected;
\item represent the four-cell square only as a projection unless a proof shows every omitted use/fiber/recovery distinction redundant;
\item apply \(\Gamma_D\) only as a downstream compatibility check;
\item keep orientation reversal distinct from along-boundary advance and from domain/world crossing;
\item preserve enough residual structure to regenerate protected sides, relation uses, fibers, recovery, and reopening routes at finer grain.
\end{enumerate}

\section{Source question-program conformance}

A conforming source inquiry runtime/compiler must:
\begin{enumerate}
\item represent \(\mathsf{Ask}(q,\kappa)\) as a typed open relation plus answer-dependent continuation rather than as an untyped text checklist;
\item permit partial supported answers to generate continuations without forcing premature singleton completion;
\item preserve per-port discharge authority when lowering to the runtime effect language;
\item distinguish relational exclusion established by question semantics from surface suppression performed by prompt rendering;
\item compile reciprocal completion as the answer-dependent program of \cref{def:recip-program}, not as six independent openings;
\item retain use provenance and partial answer sets through every reciprocal continuation;
\item allow unresolved residuals and representation defects to generate new questions.
\end{enumerate}

\section{v2 interrogative-language fixture obligations}\label{sec:v2-question-fixtures}

The following executable breakers protect the newly explicit v2.0 derived language/control
contracts.  Their presence in the canonical specification creates obligations, not passing evidence:
\begin{enumerate}
\item \texttt{QSUCC-OCC-001}. Equal semantic questions and the same whole proof-carrying
      \(\widehat S\) under two checked continuations may yield protected-different successors; a key
      containing only \((q,\widehat S)\) is rejected.
\item \texttt{QSUCC-PARTIAL-001}. A proof-carrying answer whose member projection is non-singleton is
      supplied whole to the first-order continuation; implicit member selection is rejected.
\item \texttt{QREADY-UNLOCK-001}. A supported answer supplies the exact missing dependency and
      changes a target question from not ready to ready with occurrence provenance.
\item \texttt{QREADY-NONUNLOCK-001}. Temporal adjacency or later execution without a readiness
      change is rejected as an unlock.
\item \texttt{QSTATIC-DYNAMIC-001}. A static refinement, converse, common anchor, or breaker pairing
      does not create a realized \(\QSucc\) edge; a dynamic edge retains its answer and continuation.
\item \texttt{QCONVERSE-NOT-INVERSE-001}. A many-to-one relation admits a converse/preimage with
      multiple answers while a claimed inverse is rejected.
\item \texttt{QADJOINT-001}. A universal guarantee, weakest condition, or adjoint is admitted only
      through its binding-supplied checked law, never inferred from converse incidence.
\item \texttt{QRECIP-PROV-001}. Reciprocal return retains the exact oriented negation-use,
      departure, coverage, and same-use fiber provenance and cannot be substituted by a generic
      converse route.
\item \texttt{QFRONTIER-REQDISCHARGE-001}. A continuation-equivalent check/warrant/support
      occurrence remains live when explicitly required even when a cheaper candidate strictly
      dominates it in the resource preorder, while a productive but unauthorized generated question
      is not promoted to required; conversely, an optional occurrence strictly dominated by a
      required occurrence is removed.
\item \texttt{QIFP-LOCAL-001}. Local interrogative closure is established only for a declared
      finite root/effectivity/coverage frontier and retains blocked, bounded, extension, and unknown
      exits.
\item \texttt{QIFP-REOPEN-001}. A new separator, probe, representation, binding, horizon, or
      resource regime positively reopens a previously established local fixed point.
\item \texttt{QROUTE-REGEN-001}. Removing a question or answer position reconstructs it only when
      the occurrence-specific residual fiber collapses protectedly with typing, authority, and
      provenance intact.
\item \texttt{QROUTE-ABLATE-001}. A recurrent route node is foldable only when ablation preserves
      protected behavior, reopening, first-order control, and actual ancestry.
\item \texttt{QRENDER-001}. Controlled interrogative rendering/elaboration preserves open ports,
      bindings, modes, occurrence/continuation references, converse/adjoint/reciprocal distinctions,
      and route provenance.
\item \texttt{QLOWER-001}. Erasing every used derived root, macro, and route annotation produces an
      ordinary typed relation/question/\(\mathsf{IProg}\) with equal supported branches, authority,
      actuality obligations, provenance, and protected continuations.
\item \texttt{QACTUAL-SEPARATION-001}. Two checked \(\mathsf{Ask}\) occurrences with equal semantic
      question, compiled operator, raw bytes, and semantic member projection but
      protected-different continuations produce events linked to their distinct \(\AskRef\) values.
      Occurrence formation rejects each independently forged structural position, question, answer
      slot, environment, continuation, binding/compiler version, and provenance field.
      Cold replay reconstructs the occurrence-specific \(\QSucc\); a missing or swapped event link
      is rejected, as is a mismatch in request/operator, raw-return reference, binding/compiler
      version, ledger membership, or provenance.  A mixed-port case retains every port-indexed
      route/path/event chain in one discharge bundle, including two ports sharing one explicitly
      checked event without collapsing their evidence; replay performs zero provider redispatch.
\item \texttt{QLIFT-ALLPATHS-001}. A non-singleton proof-carrying parent followed by a second
      non-singleton dependent question preserves every supported tagged path and each child port's
      discharge mode; arbitrary parent/member selection and common-mode collapse are rejected.  A
      proper finite materialization \(F\subsetneq|\widehat S|\) retains uncovered members as
      \(\mathsf{Unknown}\) and cannot claim a complete family.
\item \texttt{QRESOLUTION-GATE-001}. The same declared question exercises
      \(\mathsf{Supported}\), \(\mathsf{ExactEmpty}\), \(\mathsf{Undefined}\),
      \(\mathsf{Unsupported}\), and \(\mathsf{Unknown}\) with their distinct checked payloads; only
      \(\mathsf{Supported}(\widehat S)\) may reach \(\kappa_{\AskRef}\) or inhabit
      \(\mathsf{ResolvedOccurrence}_A\).  A non-singleton
      \(\operatorname{Run}(p,r,-)\) result retains every compatible completion; selecting only its
      first output is rejected.
\item \texttt{QCODE-TYPING-001}. Source and runtime quotation preserve result type; interpretation
      succeeds only under an admitted binding/compiler version and remains undefined for a
      mismatched version without executing or warranting the code.
\end{enumerate}

This named actuality-separation fixture traverses the non-collapse chain
\[
\AskRef\to\mathcal D\to\zeta=\mathsf{Supported}(\widehat S)
\to\kappa_{\AskRef}(\widehat S)
\to(\AskRef'\ \text{or}\ \mathsf{Return}_I),
\]
restarts from accepted references, invokes no provider again, and rejects any implementation in
which raw bytes directly choose the source answer continuation.

\section{Successor reciprocal-boundary fixture obligations}\label{sec:successor-fixtures}

The following seventy fixtures are additional to, and do not replace, the predecessor typing, compiler, actuality, history, standing, fold, binding, and self-revision fixtures.

\subsection{Determination and departure}
\begin{enumerate}
\item A source determination presentation is explicit and versioned.
\item Unrelated retained facts are not automatically constitutive.
\item Same-carrier exact cell exclusion has a constitutive separator witness.
\item Exact finite cell exclusion and constitutive separator existence coincide.
\item The 65,536-case finite feature fixture has zero mismatches.
\item Unknown observation does not establish departure.
\item Raw signature difference caused only by unknown remains unresolved.
\item Protected non-equivalence alone does not establish departure.
\item Departure may hold while source and candidate remain protected-equivalent.
\item Boundary crossing contains departure plus traversal provenance.
\item Departure does not imply observed crossing.
\item Cross-typed departure may use two observations plus an incompatibility relation.
\item Binding-native direct incompatibility may supply the witness without a shared observation codomain.
\end{enumerate}

\subsection{Typed negation}
\begin{enumerate}\setcounter{enumi}{13}
\item Boundary projection does not imply positive negation.
\item A negation use maps every exact admitted edge to a positive departure witness.
\item A proposed negation relation cannot use itself as its sole departure warrant.
\item Exact exhaustive partition complement forms lawful exact negation.
\item Certified partial opposition forms lawful partial negation.
\item Unknown is not converted into negative incidence.
\item Inapplicability is not converted into negative incidence.
\item One oriented negation use does not synthesize its reverse.
\item Exact semantic coverage and execution coverage remain distinct.
\item Empty exact exhaustive field differs from empty unsearched field.
\item No breaker under partial negation coverage does not establish reciprocal closure.
\end{enumerate}

\subsection{Multiple negation uses}
\begin{enumerate}\setcounter{enumi}{24}
\item The same exterior form through different negation uses remains distinct when return fibers differ.
\item Duplicate exterior grouping retains every witness and use.
\item Heterogeneous target carriers combine through a tagged dependent family.
\item \(\mathsf{CertifiedPartial}+\mathsf{CertifiedPartial}\) does not imply collective exhaustive coverage.
\item Collective exactness requires a cover certificate.
\item Untagged union fails the return-provenance breaker.
\item Intersection and union introduce no new exterior candidate absent from their members.
\item Ordinary relation composition may create a candidate relation but not automatic negation authority.
\item Open negation-family traversal uses generic generator fairness.
\end{enumerate}

\subsection{Return and recovery}
\begin{enumerate}\setcounter{enumi}{33}
\item For every admitted incidence, the source belongs to the reverse return section.
\item Source membership does not imply unique return determination.
\item A return fiber containing two protected source classes yields a separator obligation.
\item One selected stable \(R_X\) does not establish exact return closure if another protected class survives.
\item An exact singleton protected return fiber establishes return stability.
\item Raw relation differences ignored by the horizon do not constitute recovery failure.
\item Protected relation difference inside a return fiber does constitute recovery failure.
\item Local recovery equals protected-signature constancy on the return fiber.
\item Raw containment recovery is only a special case.
\item A source web can be partially recovered without unique source regeneration.
\item Family schema recovery may succeed where each member signature fails.
\item The three-state joint-recovery witness passes.
\item Adding a return-signature coordinate refines but does not coarsen exact family observational equivalence.
\item Historical local recovery is not retroactively strengthened by later learned negation uses.
\end{enumerate}

\subsection{Return versus reconciliation}
\begin{enumerate}\setcounter{enumi}{47}
\item Compatible monotone constraint addition cannot produce a protectedly different filling from an already determining web.
\item The general monotonicity theorem is tested against finite fixtures.
\item Generated positive exterior does not mutate standing state.
\item Actual and warranted exterior may open reconciliation.
\item State-changing redetermination exposes predecessor revision/retraction/splitting, changed applicability, grain or binding, or prior underdetermination.
\end{enumerate}

\subsection{Dependent sixfold}
\begin{enumerate}\setcounter{enumi}{52}
\item Sixfold roles are generated dependently, not independently.
\item \(O_X\) is generated from a specific negation use.
\item \(R_X\) arises from that use's reverse section.
\item \(O_X\) and \(S_Y\) retain distinct roles even if seed identity is used.
\item \(O_Y\) is independently generated from the \(Y\)-oriented frontier.
\item \(R_X\neq O_Y\) and \(R_Y\neq O_X\) remain role distinctions.
\item \(\Gamma_D\) cannot supply missing role fillings.
\item Stable \(X\)-return may coexist with unstable \(Y\)-return.
\item One-way negation does not imply reciprocal negation.
\item The sixfold occurrence view reconstructs from ordinary history and fibers.
\item Exact closure is indexed by negation semantic and execution coverage.
\item No breaker under a partial frontier remains \(\mathsf{Unknown}\), not absolute closure.
\end{enumerate}

\subsection{Representation and learning}
\begin{enumerate}\setcounter{enumi}{64}
\item A sixfold/recovery residual can generate \(\mathsf{RepresentationGap}\).
\item A new attribute can separate a previously collapsed protected pair.
\item A conservatively admitted attribute preserves old constructible questions, and
strictly enlarges the question space only when it constructs a well-typed question
outside the old image; rebinding alone establishes no such inclusion.
\item A new probe capability remains unadmitted until binding-extension checks pass.
\item A new negation-use signature may reopen a previous fold.
\item A recurrent separator may become a candidate attribute without automatic standing.
\end{enumerate}

\section{Cross-cutting derived-breaker obligations}\label{sec:derived-breakers}

The following fixtures protect derived contracts used across several later phases.
They do not add semantic primitives and do not count as additional reciprocal roles:
\begin{enumerate}
\item exact deterministic \(\chi\) factors through \(\sigma\) exactly when
\(\ker\sigma\subseteq\ker\chi\), while an incomplete signature remains
\(\mathsf{Unknown}\);
\item two individually supported signatures with mutually exclusive applicability
contexts fail joint actualization without a compatibility/jointness witness;
\item a conservative representation extension embeds every old well-typed question,
while a rebinding fixture that removes an old constructor defeats unqualified
question-universe inclusion;
\item a certified empty method solution produces a typed semantic residual, whereas a
backend crash with the same surface error text produces only an operational stop;
\item equal scalar distortion for an over-approximation and an under-approximation
does not license the same protected inference;
\item current-consequence kernel inclusion holds while a protected reopening
discriminator fails to factor, so consequence sufficiency does not masquerade as
inquiry-regenerative sufficiency;
\item a relation with neither a recovery certificate nor a non-recovery witness remains
\(\mathsf{Unknown}\) instead of entering a complement-defined irrecoverable residue;
\item one four-class protected quotient admits several incomparable minimal two-cue
separator bases;
\item an exact finite rational dot-product binding reconstructs
\(\mathbb R^d/\ker M_Q\cong\operatorname{im}M_Q\) and its rank;
\item a deterministic nonzero-mean query has zero centered covariance but a rank-one
second moment and therefore defeats covariance substitution;
\item a changed query distribution reopens the consequence-subspace quotient, while a
sample-estimated or floating-rank result remains working/approximate without certified
bounds.
\end{enumerate}

\section{Probe conformance}

Every recurrent probe occurrence must record enough to establish or reject comparability across occurrences.  Same wording is insufficient.  Changed binding or grain requires a standing bridge before direct comparison.

\section{Compression conformance}

No exact fold/quotient may be admitted without:
\begin{enumerate}
\item protected consequence sufficiency;
\item protected continuation descent where executable state/reuse is claimed;
\item hom-wise typing for operator folds;
\item provenance/residual/recovery contract;
\item unlock conditions.
\end{enumerate}

Approximate folds must state distortion.
Their conformance fixtures must also distinguish the direction or binding-specific
soundness relation of that distortion whenever direction changes what may be inferred.

\section{Warrant conformance}

Every standing hard relation must expose:
\begin{enumerate}
\item at least one closed support route;
\item scope and applicability;
\item open-dependency status;
\item warrant/checker provenance;
\item independent ingress sufficient to prevent rootless positive self-support.
\end{enumerate}

\section{Learning conformance}

A learned method/probe/operator must not confer warrant on semantic outputs merely because the traversal method is useful.  Semantic and traversal patches remain distinct.  A learned regenerative form must expose the protected candidate sides, determination presentation, negation-use provenance, return/recovery behavior, and reopening routes it claims to regenerate, or explicitly declare an approximate/reacquisition-only contract.  A new probe/tool is not admitted as a semantic discriminator merely because it was generated; the nonredundant distinction must be demonstrated through independently grounded occurrences.

\section{Self-revision conformance}

A successor presentation must be judged under predecessor-established or independently changed acceptance rules.  Candidate self-description is never sufficient evidence of strict gain.

\part{Primitive Elimination and Regenerative Basis}

\section{What is not primitive}

The following are not universal semantic primitives:
\begin{itemize}
\item absolute atomic object independent of grain;
\item Object and operator as disjoint ontological species;
\item Question as a noun/type separate from open relation;
\item Ask as a semantic oracle;
\item observation and action as disjoint interaction kinds;
\item cue;
\item memory state;
\item problem;
\item cause;
\item factor;
\item necessity;
\item sufficiency;
\item four-cell square as the complete definition of a distinction;
\item cube;
\item method;
\item controller;
\item scheduler;
\item attention;
\item learning;
\item update.
\end{itemize}

Implementation IR sorts may use these names.  Their semantics must be regenerated from the relational basis.

\section{Regenerative basis}

The regenerative basis fixed by this v2.0 document is:
\[
\boxed{
\begin{gathered}
\mathsf{Ty}_{\B},
\qquad
\mathsf{Form}_{\B}=\sum_{A:\mathsf{Ty}_{\B}}\den A_{\B},\\
\top_{\B}=\mathsf{Form}_{\B}\times\mathsf{Form}_{\B},\\
R:A\rightsquigarrow B,
\qquad
R\preceq_{\mathrm{rel}}S,\\
\text{named typed relation schemas, partial binding, exposure, hiding, and composition},\\
?_IR[\beta],
\qquad
\operatorname{Sec},\operatorname{Sol},\operatorname{Hole},\\
D=(X,Y,B_D,\pi_X,\pi_Y,\Gamma_D),\\
W_D^\alpha(s),
\qquad
\Depart_{D,W}^{\alpha}(s,e),\\
u=(N_u,D,\alpha,W,\mathsf{coverage}_{sem},\mathsf{support}),\\
\mathcal N_D^\alpha(s)=\sum_u\{u\}\times N_u[s],
\qquad
\RetField_u(e)=N_u^{-1}[e],\\
\Recover_{\mathcal H}(r,F),
\qquad
\Xi_D(\omega)=(S_X,O_X,R_X;S_Y,O_Y,R_Y),\\
\widetilde B_D=B_D\times\{+,-\},
\qquad
T_D:B_D\rightsquigarrow B_D,\\
o:Z\rightsquigarrow R_o\times Z,\\
K::=\mathsf{Return}_I\mid\mathsf{Ask},\\
P::=\mathsf{Return}\mid\mathsf{Branch}\mid\mathsf{Probe},\\
\text{reification and typed operational interpretation},\\
\text{authoritative linked occurrence history},\\
Regen_{\mathcal H},\\
\text{support environments and }\Stand=\mu T,\\
\text{typed folds with residual/recovery/unlock},\\
\text{reification and predecessor-judged revision}.
\end{gathered}
}
\]

Everything else in this document is a law over this basis, a derived construction, a binding-supplied specialization, a metatheoretic correspondence, or implementation infrastructure.

\section{Compact invariant}

\begin{quote}
\small
Every consequentially represented value, relation, distinction, question, operator, or program is a typed form and may enter further lawful relations without losing its type.  Any two represented forms share at least the coarsest relation of joint admission to the current field; inquiry seeks consequentially finer relations.  When a current representation identifies forms that a protected consequence distinguishes, the calculus preserves an obligation to invent or discover a separating context, representation, grain, probe/tool, language extension, or binding.  When an active distinction cannot affect any protected continuation, it is a candidate for licensed folding.

A live determination is represented by an explicit, supported relational presentation.  Its exterior is not presumed from a boundary projection and is not obtained by Boolean complement.  A candidate is positively established as outside the determination when a standing determination-relevant discriminator places source and candidate in incompatible typed cells.  An oriented relation whose admitted incidences have such departure witnesses may serve as a typed negation use, with semantic coverage stated independently of execution coverage.

A determination may admit multiple negation uses.  They remain a tagged family: the relation use that licensed an exterior remains part of the reciprocal occurrence because different uses can return different source fields.  Inquiry partially binds one such relation to the source and positively determines an exterior filling.  Pure return is the reverse section of that same relation.  The return fiber recovers a protected source relation exactly when every source still possible through that fiber agrees on its protected signature.  Multiple negation uses accumulate reusable information through the product of their return signatures, not through an untagged semantic union; that product does not become one actual composite return without supported jointness evidence.

The exterior filling is then taken as a new center of determination through an explicit seed/reorientation relation, and the same operation is repeated in the reciprocal orientation.  The six dependent roles \((S_X,O_X,R_X;S_Y,O_Y,R_Y)\) are therefore a derived reciprocal trace, not six independent openings and not a separate history.  Exact reciprocal closure is fiber- and coverage-relative.  \(\Gamma_D\) checks joint compatibility only after the dependent roles have been generated.

A generated exterior does not revise standing semantics.  If warranted actuality invalidates the standing source web, reconciliation/revision produces a successor web; compatible monotone constraint addition alone cannot transform an already determined source into a protectedly different source.  Every surviving return ambiguity, recovery failure, seed mismatch, reciprocal mismatch, or compatibility failure becomes an ordinary protected residual and feeds the generic separator engine.  When a required separator is absent only from current materialization, generation continues within the admitted language.  When no admitted question, probe, or representation can express a witnessed protected distinction, the result is a representation or binding gap.

A coverage-indexed characterization may be derived from the supported determination,
certified internal variation, tagged exterior frontier, and three-valued recovery/loss
profiles.  It remains relative and versioned; generated constitutive discriminators
require independent admission.  Active representations approach a regenerative
economy frontier by differentiating until protected live classes are separable and
subtracting until any further removal would lose regeneration, continuation,
provenance, warrant, or reopening.

A question is a partially bound typed relation.  Binding what is known while preserving the relation leaves a residual section/fiber: a structured hole whose lawful fillers are the answer field.  A web of such residual relations determines a form to the extent that the intersection of its solution fibers collapses to one protected equivalence class.  Abstraction is therefore regenerative when it removes a filling but preserves enough relational constraint to regenerate the protected form and the discriminators needed to reopen it.  A source inquiry program composes these open relations through answer-dependent continuations and lowers by typed discharge authority to pure computation, generation, actual probing, checking, and warrant.

Dynamic question succession is indexed by a checked \(\mathsf{Ask}\) occurrence \(\AskRef\), not
only by normalized question text or relation identity.  A supported answer is a proof-carrying
record with a nonempty member set of typed completions, retained whole with its declared route,
evidence, coverage, and per-component
provenance.  Capture-safe substitution into the occurrence-specific first-order continuation yields
\(\QSucc(\AskRef,\widehat S,q')\).  Static question relations do not manufacture this dynamic edge.  The
live frontier includes both protectedly productive occurrences and explicit required-discharge
occurrences; required but nonexecutable obligations survive as typed residuals.

The runtime preserves realized state--operator--return history before interpretation.  Raw return, decoded completion, interpretation, checking, warrant, and standing remain distinct.  Learned distinctions, representations, probes, and operator paths are compact regenerative structures over authoritative history; they never replace the evidence that warranted them and never confer warrant on future outputs merely by being useful.  Every quotient or fold is protected-behavior relative, continuation-compatible where reusable state is claimed, regeneratively sufficient or explicitly approximate, and equipped with residual, recovery/reacquisition, and unlock conditions.

There is one authoritative actuality spine: accepted source/program/artifact ancestry plus ordinary
event and accepted-revision histories.  Paired actuality, question traces, return traces, route
records, sufficient presents, cue fibers, learned methods, and reciprocal views are derived checked
projections.  Resume is not replay, append order is not causal order, and endpoint equality is not
event/path identity.  Every recurring inquiry is a specialization of
\(\mathsf{BIND}\to\mathsf{OPEN}\to\mathsf{VARY}\to\mathsf{RETURN}\to
\mathsf{DETERMINE}\to\mathsf{REFACTOR}\); no specialization adds a second controller or history.

The same forms, distinctions, questions, programs, probes, compiler rules, folds, bindings, and protection declarations may be reified and subjected to the same inquiry.  A successor remains subject to independently grounded, predecessor-preserving judgment and cannot license itself.
\end{quote}

\section{Fixed-point target}

The project approaches a regenerative fixed point when:
\begin{center}
\fbox{\parbox{0.90\linewidth}{\centering\bfseries
THE CALCULUS IS CLOSED UNDER REIFICATION, RELATION REFINEMENT, DISTINCTION FORMATION, POSITIVE TYPED NEGATION, SAME-USE RETURN, DEPENDENT RECIPROCAL OPENING, HOLE FORMATION, QUESTION-PROGRAM COMPOSITION, COMPILATION, ACTUALIZATION, REPRESENTATION/PROBE INVENTION, TRAVERSAL, FOLDING, REOPENING, AND RE-ENTRY OF ITS OWN TYPED FORMS WITHOUT SELF-WARRANT.
}}
\end{center}
No claim is made that global convergence, absolute minimality, or a universally coarsest sufficient representation is decidable.

\clearpage
\appendix

\section{Canonical notation table}

\begin{longtable}{>{\raggedright\arraybackslash}p{0.34\textwidth}
                  >{\raggedright\arraybackslash}p{0.59\textwidth}}
\toprule
Symbol & Meaning\\
\midrule
\endhead
\(\B\) & current binding\\
\(\mathsf{Form}_{\B}\) & dependent universe of typed represented forms\\
\(\top_{\B}\) & coarsest relation of joint admission to the represented field\\
\(R\preceq_{\mathrm{rel}}S\) & \(S\) is a refinement/more discriminating relation than \(R\)\\
\(\operatorname{Sec}^{X}_{\rho}(y)\) & residual section after binding the \(Y\)-port\\
\(\operatorname{Sol}^{X}_{\rho}(y)\) & lawful filler fiber for the remaining \(X\)-port\\
\(\operatorname{Hole}_x(W)\) & relational web with protected occurrences of \(x\) exposed as holes\\
\(Determines_{\mathcal H}(W,x)\) & residual web regenerates one protected class \([x]_{\mathcal H}\)\\
\(W_D^\alpha(s)\) & supported, versioned determination presentation for source \(s\) in orientation \(\alpha\)\\
\(\Depart_{D,W}^{\alpha}(s,e)\) & positive determination-relative departure with a standing witness\\
\(N_u\) & ordinary oriented typed relation under immutable negation-use identity \(u\)\\
\(\NegField_u(s)=N_u[s]\) & source-bound positive-negation candidate field for use \(u\)\\
\(\mathcal N_D^\alpha(s)\) & tagged dependent family of admitted negation-use frontiers\\
\(\RetField_u(e)=N_u^{-1}[e]\) & same-use reverse return fiber\\
\(\Recover_{\mathcal H}(r,F)\) & protected signature of relation \(r\) is constant over return fiber \(F\)\\
\(\mathsf{RecoveryStatus}_{\omega}(r,e;u)\) & certified recovery, witnessed non-recovery, or unknown at an occurrence\\
\(\Xi_D(\omega)\) & dependent sixfold role view reconstructed for reciprocal occurrence \(\omega\)\\
\(\widetilde B_D\) & oriented double cover \(B_D\times\{+,-\}\)\\
\(\iota\) & reciprocal orientation involution on the double cover\\
\(Regen_{\mathcal H}(m,z)\) & representation \(m\) regenerates protected structure of \(z\)\\
\(\mathsf{Economy}_{\mathcal H,\preceq}(z)\) & minimal licensed representations that preserve inquiry regeneration\\
\(K::=\mathsf{Return}_I\mid\mathsf{Ask}\) & source inquiry question-program language\\
\(\AskRef\) & checked source \(\mathsf{Ask}\) occurrence with explicit environment, answer slot, continuation, versions, and provenance\\
\(\SupportedRoute\) & \(\SupportedRoute(q,S;\vartheta,E)\) checks admission of one nonempty member set through its declared route and evidence environment\\
\(\widehat S:\mathsf{SuppAns}(q)\) & proof-carrying supported answer record; \(|\widehat S|\) is its nonempty member set\\
\(\QSucc(\AskRef,\widehat S,q')\) & occurrence-indexed derived successor-question relation\\
\(\QStep(\Sigma,\AskRef,\widehat S)\) & protected successor presentation after capture-safe answer substitution and pure reconstruction\\
\(\RequiredDischarge(\Sigma,\AskRef,d)\) & explicit dependency \(d\) requires this occurrence's declared discharge; it does not imply executability or success\\
\(\Omega_Q\) & derived, unordered root-operator family \(\{\mathsf{Expose},\mathsf{Orient},\mathsf{Factor},\mathsf{Polarize},\mathsf{Vary},\mathsf{Ground}\}\)\\
\(\IFP_{\Omega,\mathcal H,\epsilon}\) & coverage-relative local interrogative fixed point, reopenable by new relations, evidence, bindings, horizons, or resources\\
\(\mathsf{MasterRecurrence}\) & \(\mathsf{BIND}\to\mathsf{OPEN}\to\mathsf{VARY}\to\mathsf{RETURN}\to\mathsf{DETERMINE}\to\mathsf{REFACTOR}\), the one derived high-level control recurrence; not runtime opcodes\\

\(A,B,X,Y,Z\) & typed carriers/interfaces\\
\(R:A\rightsquigarrow B\) & typed relation\\
\(R\hookrightarrow\prod_i X_i\) & named-port relation schema\\
\(?_I R[\beta]\) & question obtained by binding closed ports and exposing \(I\)\\
\(A(q)\) & complete semantic answer carrier\\
\(\Fib_I(R\mid\beta)\) & completion fiber\\
\(\operatorname{mode}(q,i)\) & lawful discharge authority for open port \(i\)\\
\(\sigma_q(x)\) & set of complete answers compatible with \(x\)\\
\(x\sim_q y\) & local equivalence induced by question \(q\)\\
\(q\preceq_{\mathrm{prec}}r\) & \(r\) refines/discriminates at least as much as \(q\)\\
\(\eta_{q,g}\) & question-conditioned active view at grain \(g\)\\
\(V_t\) & active representation codomain at step \(t\)\\
\(FD_{\mathcal H}(q)\) & protected distinctions hidden by the local question frame\\
\(D^+\leftarrow B_D\rightarrow D^-\) & candidate boundary factorization; not itself a negation\\
\(T_D\) & along-boundary traversal relation\\
\(\mathsf{coverage}_{sem}/\mathsf{coverage}_{exec}\) & independent semantic-negation and occurrence-execution coverage\\
\(\Prof_D\) & exact boundary-position profile\\
\(o:Z\rightsquigarrow R_o\times Z\) & probe interaction semantics\\
\(z_t\xrightarrow[o]{r}z_{t+1}\) & realized probe occurrence\\
\(\delta_q\) & raw/decoded response to supported answer-set map\\
\(\ResPath:A\leadsto B\) & explicit resolution path; distinct from a generic residual relation \(\rho\)\\
\(\mathsf{ResolvedOccurrence}_A\) & dependent record \(\AskRef\to\mathcal D\to\mathsf{Supported}(\widehat S)\to\kappa_{\AskRef}(\widehat S)\to\mathsf{Next}_A\), retaining every port/event/path and excluding non-Supported outcomes\\
\(\mathsf{Return}\) & pure program return\\
\(\mathsf{Branch}\) & internal/generative alternatives\\
\(\mathsf{Probe}\) & actualizable program interaction\\
\(H^{evt}\) & authoritative actual event history\\
\(H^{acc}\) & accepted patch/revision history\\
\(H^{auth}\) & one authoritative typed ancestry whose event and accepted-revision histories are non-collapsing projections\\
\(<_{\mathrm{ledger}}\) & persistence append order\\
\(\lambda_{\mathrm{real}}\) & binding-defined realized/domain succession\\
\(\equiv_{\mathcal H}\) & exact protected behavioral equivalence\\
\(\approx_{\mathcal H,\mathcal D}\) & finite tested nondistinction\\
\(\mathsf Q_{\mathcal H}(A,B)\) & hom-wise operator quotient\\
\(\Supp(\lambda)\) & minimal support environments for relation \(\lambda\)\\
\(\partial^-_E\lambda\) & open dependency boundary under support environment \(E\)\\
\(\Stand=\mu T\) & least-fixed-point standing\\
\(\Lambda_c\) & compression/fold licence\\
\(\Unlock(c)\) & observation/dynamic/context changes invalidating a fold\\
\(\PatchIR[k]\) & role-indexed candidate revision\\
\(Q^\infty\) & semantic question universe\\
\(Q^\epsilon\) & executable resource-bounded question universe\\
\bottomrule
\end{longtable}

\section{Canonical compiler pipeline}

The single language stack in \cref{sec:language-architecture} has the following typed operational
projection:
\[
\boxed{
\begin{array}{c}
\text{typed forms and named-port relation schema}\\
\downarrow\ \text{partial bind / expose a nonempty port set}\\
q=?_I R[\beta]\quad\text{with }A(q)=\prod_{i\in I}X_i\\
\downarrow\ \text{store first-order answer slot and continuation}\\
\Sigma:\mathsf{SourceConfig}_{\B}(A),\quad
\AskRef:\mathsf{AskOcc}_A(\Sigma)\\
\downarrow\ \text{normalize and compile each port by its declared mode}\\
\OpenIR\to(\text{pure code / generator / probe / checker / warrant route})\\
\downarrow\ \text{execute every port only through its selected lawful route}\\
Y_i:\mathsf{PureResult}_i+\mathsf{GeneratedProposal}_i+\mathsf{ActualReturn}_i
  +\mathsf{CheckedResult}_i+\mathsf{WarrantResult}_i\\
Y:=\prod_{i\in I}Y_i\quad\text{selected pointwise by }\operatorname{mode}(q,i)\\
\downarrow\ \text{for each actual-return component: preserve }E_i/\ReturnRef_i\text{ first}\\
\mathcal D:\mathsf{DischargeBundle}(\AskRef)\\
\downarrow\ \text{resolve and check declared-route support without strengthening}\\
\zeta:\mathsf{ResolutionOutcome}(q)\\
\downarrow\ \text{case analysis retaining the constructor payload}\\
\begin{cases}
\mathsf{Supported}(\widehat S):
  \kappa_{\AskRef}(\widehat S)
  \to \QSucc(\AskRef,\widehat S,q')\ \text{or}\ \mathsf{Return}_I(a),\\
\mathsf{ExactEmpty}/\mathsf{Undefined}/\mathsf{Unsupported}/\mathsf{Unknown}:
  \text{typed residual or justified stop; never }\kappa_{\AskRef}.
\end{cases}\\
\downarrow\ \text{pure normalization / residual construction with the next live frontier.}
\end{array}
}
\]

Positive-negation compilation is one specialization inside this pipeline.  When the open relation
is an admitted oriented negation use, normalization retains the determination presentation,
candidate chart, use identity, departure witness, semantic/execution coverage, same-use return
fiber, recovery signature, seed, reciprocal occurrence, residuals, and downstream-only
\(\Gamma_D\) check.  These fields refine \(\OpenIR\), support, and residual construction; they do
not replace the general pipeline or create another history.

\section{Canonical infrastructure dependency direction}

The dependency direction is:
\[
\boxed{
\text{calculus semantics}
\to
\text{IR contracts}
\to
\text{compiler/runtime}
\to
\text{storage/indexes/backends}.
}
\]

Storage layout, database schema, scheduler, provider SDK, prompt templating library, or agent framework must not redefine semantics in the opposite direction.

\section{Minimum conformance suite}

A future implementation must contain executable witnesses for at least:
\begin{enumerate}
\item typed relation composition and ill-typed rejection;
\item typed reification of an operator as an operand and rejection of an untyped operational interpretation;
\item a pair of represented forms related only by \(\top_{\B}\) and a discovered proper relation refinement;
\item recursive distinction composition in which a distinction point becomes a side of another distinction;
\item positive departure established by a relevant standing discriminator while unknown, failed search, projection, and non-equivalence are rejected as substitutes;
\item same exterior through different use identities yielding different return fibers without provenance collapse;
\item a selected stable return rejected as closure while a second protected class survives in the return fiber;
\item departure coexisting with protected equivalence under the current horizon;
\item family product recovery succeeding where each individual return signature fails;
\item family product information rejected as one actual composite return when the
component applicability contexts are mutually exclusive and no jointness witness exists;
\item exact deterministic determination-through agreeing with kernel inclusion while
an incomplete-coverage case remains \(\mathsf{Unknown}\);
\item a current-consequence-sufficient representation rejected as inquiry-regenerative
when it loses a protected reopening discriminator;
\item a recovery/loss profile retaining \(\mathsf{Unknown}\) separately from witnessed
non-recovery;
\item compatible monotone return preserving an already determined class while warranted reconciliation remains separately capable of revision;
\item a dependent sixfold occurrence retaining \(R_X\neq O_Y\), \(R_Y\neq O_X\), one-sided stability, and a downstream-only \(\Gamma_D\);
\item oriented double-cover traversal alternating \(\iota\) and \(\tau_D\) without synthesizing reverse negation or conflating either with domain crossing;
\item a four-cell square projection demonstrably losing protected use, fiber, recovery, or provenance information from the full dependent occurrence;
\item a residual hole whose web solution set has several fillers but one forced protected property;
\item indexed-meet refinement in which adding one relation strictly narrows the filler field;
\item a protectedly unique filler class reconstructed after ablating the original form;
\item a source \(\mathsf{Ask}\) program whose partial answer constructs a new question rather than selecting a fixed branch;
\item a representation defect \(\eta(x)=\eta(y)\) with protected-different consequence that generates a separating representation/probe inquiry;
\item a new probe/tool basis element exposing a protected distinction unavailable to the old basis;
\item conservative question-language extension transporting every old question and a
rebinding counterexample defeating unqualified inclusion;
\item several incomparable minimal separator bases inducing the same protected
four-class quotient;

\item one relation schema generating several questions by different open ports;
\item discharge-authority rejection of generative filling for an actual/check/warrant port;
\item question refinement with kernel intersection;
\item a nonempty local question-frame defect safely preserved by surrounding history;
\item partial response retaining multiple answer cells;
\item along-boundary traversal with both candidate projections retained;
\item one unified probe whose world projection is observation-like and another action-like;
\item raw-return preservation before decoding;
\item ledger order demonstrably distinct from domain succession;
\item recurrent probe comparability failure after an unbridged binding/grain change;
\item operator path folding with hom-wise typing and provenance;
\item continuation-descent failure preventing an output-sufficient quotient from becoming executable state;
\item multiple independent support environments for one relation;
\item rootless support cycle failing to enter standing;
\item exact compression reopening by a new observation;
\item reopening by a new protected continuation even without a new current observation;
\item approximate compression retaining a distortion/residual contract;
\item equal scalar distortion for over- and under-approximation rejected as authority
to discharge the same protected claim;
\item certified method non-discharge producing a typed residual while an operational
backend failure with the same surface message remains a nonsemantic stop;
\item an exact rational consequence-subspace quotient with rank certificate, a
nonzero-mean covariance breaker, and query-distribution-change reopening;
\item semantic versus traversal patch separation;
\item self-revision candidate failing self-promotion without independent evidence;
\item semantic frontier containing an unavailable question while executable frontier lawfully returns \(\mathsf{Unknown}\);
\item LLM prompt compilation preserving typed open ports across canonical render/elaborate round trip;
\item fresh-probe versus history-conditioned-probe test for anchoring/order sensitivity;
\item learned probe-basis extension demonstrating a protected distinction unavailable to the old basis;
\item state history remaining independently accessible after operator compression.
\end{enumerate}

In addition, every one of the nineteen identified obligations in
\cref{sec:v2-question-fixtures} and every one of the seventy numbered obligations in
\cref{sec:successor-fixtures} requires its declared executable discriminator before the
corresponding conformance may be claimed.

\section{Normative reading rule}

When future work appears to conflict with this document, apply the following order:
\begin{enumerate}
\item identify the exact typed relation and protected consequence in conflict;
\item distinguish semantic law from binding-specific or implementation-only structure;
\item produce a breaker/counterexample or a preservation proof;
\item if the document is wrong, create a typed successor patch rather than silently editing definitions in prose;
\item preserve the old statement as ancestry with an explicit disposition;
\item update this single canonical document only after the successor is independently accepted.
\end{enumerate}

This rule is part of the project's defense against information loss: discoveries do not survive merely as chat prose, and old concepts do not survive merely because a previous file named them.  Consequential structure survives as a typed law, derivation, regression, residual, or reopening path in the current canonical specification.

\end{document}
